%============================================================
%  Bd. 07 - Surjektive Funktionen
%============================================================

\documentclass[main.tex]{subfiles}

\ifSubfilesClassLoaded{
  \BandSetupStandalone{B07}
}{
}

\title{Bd. 07 - Surjektive Funktionen}
\author{Martin Kunze}
\date{}
\setcounter{file}{7}

\hypersetup{
  pdftitle={Bd. 07 - Surjektive Funktionen},
  pdfauthor={Martin Kunze}
}

\begin{document}

\title{Bd. 07 - Surjektive Funktionen}
\author{Martin Kunze}
\date{}
\hypersetup{pageanchor=false}
\maketitle
\hypersetup{pageanchor=true}
\setcounter{tocdepth}{3}
\tableofcontents
\setcounter{file}{7}
\renewcommand{\FormulaBandID}{7}

\chapter{Einleitung}

Dieser Band entwickelt die Surjektivität auf Grundlage der allgemeinen
Funktionstheorie aus Band~05. Einzelne Konstruktionen verwenden außerdem die
in Band~06 eingeführte Injektivität. Begriff und Grundaxiom stehen am Anfang;
anschließend werden die Projektionen erstmals eingeführt und durch Bild-,
Faser-, Einschränkungs- und Kompositionssätze ergänzt.
Das Auswahlaxiom wird hier ausdrücklich nicht vorausgesetzt, sondern erst im
eigenständigen Band~09 untersucht.

\chapter{Axiomatischer Aufbau der surjektiven Funktionen}

\section{Begriff, Grundaxiom und Folgerungen}

\FormulaDefDeltaK[Begriff der surjektiven Funktion]{F\colon A\twoheadrightarrow B}{Surjektive Funktion}{
  \DeltaRow{Mengen}{A\dsep B\dsep x\dsep y}
  \DeltaPrem{Funktionen}{F\colon A\to B}[\FormulaRefAuto{Funktion}]
%
  \DeltaRow{\textbf{Axiome}}{}
  %
  % — Surjekvität —
  \DeltaRow{Surjektivität}
           {y\in B\vdash \exists x\in A\; F(x)=y}
           [\FormulaRefAuto{y\in B\vdash \exists x\in A\; F(x)=y}]
  \DeltaRow{\textbf{Neue Symbole}}{}
  \DeltaPrem{Surjektive Funktionen}{ }
}

\FormulaAxiomDeltaK[Surjektivität]%
{y\in B\vdash \exists x\in A\; F(x)=y}{Surjektivität}
{
\DeltaRow{Mengen}{y\dsep A\dsep B}
\DeltaPrem{Funktionen}{F\colon A\to B}
}

\FormulaThmDelta{F\colon A\sur B\dsep x\in A\vdash F(x)\in B}{
\DeltaRow{Mengen}{A\dsep B}
}
\begin{tabproof}
  \proofstep{1}{F\colon A\sur B}{\rA}
  \proofstep{2}{x\in A}{\rA}
  \proofstep{1}{F\colon A\to B}{\FormulaRefAuto{Surjektive Funktion}{1}}
  \proofstep{1,2}{F(x)\in B}{%
    \FormulaRefAuto{F\colon A\to B\dsep x\in A\vdash F(x)\in B}{3,2}}
\end{tabproof}

\FormulaThmDelta[Elemente einer Teilmenge liegen im Bild der Teilmenge unter surjektiven Funktionen]{%
  C\subseteq A\dsep F\colon A\sur B\dsep x\in C\vdash F(x)\in F[C]
}{
\DeltaRow{Mengen}{A\dsep B\dsep C\dsep x}
\DeltaRow{Funktionen}{F}
}
\begin{tabproof}
  \proofstep{1}{C\subseteq A}{\rA}
  \proofstep{2}{F\colon A\sur B}{\rA}
  \proofstep{3}{x\in C}{\rA}
  \proofstep{2}{F \colon A \to B}{%
      \FormulaRefAuto{Surjektive Funktion}{2}}
  \proofstep{1,2,3}{F(x)\in F[C]}{%
    \FormulaRefAuto{C\subseteq A\dsep F\colon A\to B\dsep x\in C\vdash F(x)\in F[C]}{%
      1,4,3}}
\end{tabproof}

\FormulaThmDeltaR[Nichtsurjektivität liefert ausgelassenes Element]{%
  \begin{aligned}[t]
    &F\colon M\to M\dsep \neg(F\colon M\sur M)\vdash{}\\
    &\exists m_0\in M\,\forall x\in M\,\bigl(F(x)\neq m_0\bigr)
  \end{aligned}
}{%
  F\colon M\to M\dsep \neg(F\colon M\sur M)\vdash
  \exists m_0\in M\,\forall x\in M\,\bigl(F(x)\neq m_0\bigr)
}{%
  \DeltaDecl{Mengen}{M\dsep m_0\dsep x\dsep y}%
  \DeltaDecl{Funktionen}{F}%
}
\begin{tabproof}
  \proofstep{1}{F\colon M\to M}{\rA}
  \proofstep{2}{\neg(F\colon M\sur M)}{\rA}

  \proofstep{3}{\forall y\in M\,\exists x\in M\,\bigl(F(x)=y\bigr)}{\rA}

  \proofstep{1,3}{F\colon M\sur M}{\FormulaRefAuto{Surjektive Funktion}{1,3}}

  \proofstep{1,2,3}{\bot}{\rBI{2,4}}

  \proofstep{1,2}{%
    \neg\forall y\in M\,\exists x\in M\,\bigl(F(x)=y\bigr)
  }{\rCI{3,5}}

  \proofstep{1,2}{%
    \exists m_0\in M\,\forall x\in M\,\bigl(F(x)\neq m_0\bigr)
  }{%
    \FormulaRefAuto{%
      \exists a\in A\,\forall b\in B\,(\neg F(a,b))
      \eqvdash
      \neg\forall a\in A\,\exists b\in B\,(F(a,b))
    }[6]%
  }
\end{tabproof}

\FormulaThmDeltaK[Korestriktion auf das Bild als surjektive Funktion]{%
  F\colon A\to B \vdash F\colon A\sur F[A]%
}{CorestrictionToImageSurjective}{
  \DeltaRow{Mengen}{A\dsep B\dsep x\dsep y}
  \DeltaRow{Funktionen}{F}
}
\begin{tabproofsplit}
\proofpart{Trägerrelation}
\proofstep{1}{F\colon A\to B}{\rA}
  \proofstep{1}{F\subseteq A\times F[A]}{%
    \FormulaRefAuto{F\colon A\to B\vdash F\subseteq A\times F[A]}{1}}

\closeproofpart

\proofpart{Funktionale Eindeutigkeit}
  \setcounter{proofstepnr}{2}% Fortlaufende Schritte dieses Beweises.
\proofstep{1}{\forall x,y\,((x,y)\in F\dsep (x,z)\in F\rightarrow y=z)}{%
    \FormulaRefAuto{(x,y)\in F\dsep (x,z)\in F \vdash y=z}{1}}

\closeproofpart

\proofpart{Totalität}
  \setcounter{proofstepnr}{3}% Fortlaufende Schritte dieses Beweises.
\proofstep{1}{\forall x\,(x\in A\rightarrow \exists y\,(x,y)\in F)}{%
    \FormulaRefAuto{x\in A \vdash \exists y\,(x,y)\in F}{1}}


\closeproofpart

\proofpart{Surjektivität}
  \setcounter{proofstepnr}{4}% Fortlaufende Schritte dieses Beweises.
\proofstep{5}{y\in F[A]}{\rA}
  \proofstep{1,5}{\exists x\in A\,y=F(x)}{%
    \FormulaRefAuto{F\colon A\to B\dsep y\in F[A] \vdash \exists x\in A\,y=F(x)}{1,5}}

  \proofstep{7}{x\in A\land y=F(x)}{\rA}
  \proofstep{7}{x\in A}{\rAEa{7}}
  \proofstep{7}{y=F(x)}{\rAEb{7}}
  \proofstep{7}{F(x)=y}{\FormulaRefAuto{a=b\vdash b=a}{9}}
  \proofstep{7}{x\in A\land F(x)=y}{\rAI{8,10}}
  \proofstep{7}{\exists x\in A\,F(x)=y}{\rEI{11}}

  \proofstep{1,5}{\exists x\in A\,F(x)=y}{\rEE{6,7,12}}
  \proofstep{1}{\forall y\,(y\in F[A]\rightarrow \exists x\in A\,F(x)=y)}{\rUI{\rRI{5,13}}}


\closeproofpart

\proofpart{Zusammenführung}
  \setcounter{proofstepnr}{14}% Fortlaufende Schritte dieses Beweises.
\proofstep{1}{F\colon A\sur F[A]}{\FormulaRefAuto{Surjektive Funktion}{2,3,4,14}}

\closeproofpart
\end{tabproofsplit}

\FormulaThmDeltaK[Surjektivität genau bei vollem Bild]{%
  F\colon A\to B\vdash
  \bigl(F\colon A\sur B\leftrightarrow F[A]=B\bigr)%
}{SurjectiveIffImageEqualsCodomain}{%
  \DeltaRow{Mengen}{A\dsep B\dsep x\dsep y}
  \DeltaRow{Funktionen}{F}
}
\begin{tabproofsplitwide}
  \proofpartwideindR[Surjektivität liefert das volle Bild]{%
    F\colon A\to B\dsep F\colon A\sur B\vdash F[A]=B%
  }{%
    F\colon A\to B\dsep F\colon A\sur B\vdash F[A]=B%
  }[SurjectiveImageEqualsCodomain]
    \proofstepwidestar[1]{F\colon A\to B}{\rA}
    \proofstepwidestar[2]{F\colon A\sur B}{\rA}
    \proofstepwidestar[1]{F[A]\subseteq B}{%
      \FormulaRefAuto{F\colon A\to B \vdash F[A]\subseteq B}[thm]{1}}

    \proofstepwidestar[4]{y\in B}{\rA}
    \proofstepwidestar[2,4]{\exists x\in A\,F(x)=y}{%
      \FormulaRefAuto{Surjektivität}[ax]{2,4}}
    \proofstepwidestar[6]{x\in A}{\rA}
    \proofstepwidestar[7]{F(x)=y}{\rA}
    \proofstepwidestar[7]{y=F(x)}{%
      \FormulaRefAuto{a=b\vdash b=a}[thm]{7}}
    \proofstepwidestar[6,7]{\exists x\in A\,y=F(x)}{%
      \rEI{\rAI{6,8}}}
    \proofstepwidestar[1,6,7]{y\in F[A]}{%
      \FormulaRefAuto{%
        F\colon A\to B\dsep
        \exists x\in A\,y=F(x)\vdash y\in F[A]}[thm]{1,9}}
    \proofstepwidestar[1,2,4]{y\in F[A]}{%
      \rEE{5,6,7,10}}
    \proofstepwidestar[1,2]{B\subseteq F[A]}{%
      \FormulaRefAuto{SubsetIntroductionRule}[thm]{[4]\vdots 11}}
    \proofstepwidestar[1,2]{F[A]=B}{%
      \FormulaRefAuto{%
        A \subseteq B, B \subseteq A \vdash A = B}[thm]{3,12}}
  \closeproofpartwideindR

  \proofpartwideindR[Das volle Bild liefert Surjektivität]{%
    F\colon A\to B\dsep F[A]=B\vdash F\colon A\sur B%
  }{%
    F\colon A\to B\dsep F[A]=B\vdash F\colon A\sur B%
  }[ImageEqualsCodomainSurjective]
    \proofstepwidestar[1]{F\colon A\to B}{\rA}
    \proofstepwidestar[2]{F[A]=B}{\rA}
    \proofstepwidestar[2]{B=F[A]}{%
      \FormulaRefAuto{a=b\vdash b=a}[thm]{2}}

    \proofstepwidestar[4]{y\in B}{\rA}
    \proofstepwidestar[2,4]{y\in F[A]}{\rIE{3,4}}
    \proofstepwidestar[1,2,4]{\exists x\in A\,y=F(x)}{%
      \FormulaRefAuto{%
        F\colon A\to B\dsep y\in F[A]
        \vdash \exists x\in A\,y=F(x)}[thm]{1,5}}
    \proofstepwidestar[7]{x\in A}{\rA}
    \proofstepwidestar[8]{y=F(x)}{\rA}
    \proofstepwidestar[8]{F(x)=y}{%
      \FormulaRefAuto{a=b\vdash b=a}[thm]{8}}
    \proofstepwidestar[7,8]{\exists x\in A\,F(x)=y}{%
      \rEI{\rAI{7,9}}}
    \proofstepwidestar[1,2,4]{\exists x\in A\,F(x)=y}{%
      \rEE{6,7,8,10}}
    \proofstepwidestar[1,2]{%
      \forall y\in B\,\exists x\in A\,F(x)=y}{%
      \rUI{\rRI{4,11}}}
    \proofstepwidestar[1,2]{F\colon A\sur B}{%
      \FormulaRefAuto{Surjektive Funktion}[def]{1,12}}
  \closeproofpartwideindR

  \proofpartwide{Äquivalenzschluss}
    \proofstepwidestar[1]{F\colon A\to B}{\rA}
    \proofstepwidestar[2]{F\colon A\sur B}{\rA}
    \proofstepwidestar[1,2]{F[A]=B}{%
      \FormulaRefAuto{SurjectiveImageEqualsCodomain}[thm]{1,2}}
    \proofstepwidestar[1]{F\colon A\sur B\rightarrow F[A]=B}{%
      \rRI{2,3}}
    \proofstepwidestar[5]{F[A]=B}{\rA}
    \proofstepwidestar[1,5]{F\colon A\sur B}{%
      \FormulaRefAuto{ImageEqualsCodomainSurjective}[thm]{1,5}}
    \proofstepwidestar[1]{F[A]=B\rightarrow F\colon A\sur B}{%
      \rRI{5,6}}
    \proofstepwidestar[1]{%
      F\colon A\sur B\leftrightarrow F[A]=B}{%
      \rLRI{4,7}}
  \closeproofpartwide
\end{tabproofsplitwide}

% ------------------------------------------------------------
% Hilfslemma: der breite Äquivalenz-Block (Diagonal-Schritt)
% ------------------------------------------------------------
\FormulaThmDelta{%
  a\in A \dsep F(a)=\{x\in A\mid x\notin F(x)\}\vdash a\in F(a)\leftrightarrow a\notin F(a)%
}{
  \DeltaRow{Mengen}{A\dsep a\dsep x}
  \DeltaPrem{Funktionen}{F\colon A\to \powerset(A)}
}
\begin{tabproofwide}
  \proofstepwidestar[1]{a\in A}{\rA}
  \proofstepwidestar[2]{F(a)=\{x\in A\mid x\notin F(x)\}}{\rA}

  \proofstepwidestar[2]{\forall x\,(x\in F(a)\leftrightarrow x\in\{x\in A\mid x\notin F(x)\})}{%
      \FormulaRefAuto{\forall x\, (x \in A \leftrightarrow x \in B) \eqvdash A = B}{2}}
  \proofstepwide[1,2]{a\in F(a)}{\leftrightarrow}{%
    \begin{aligned}[t]
      a\in \{x\in A\mid {}\\
      x\notin F(x)\}
    \end{aligned}
  }{\rUE{3}}

  \proofstepwide[1,2]{}{\leftrightarrow}{a\in A \land a\notin F(a)}{%
    \FormulaRefAuto{x \in \{u \in A \mid P(u)\} \eqvdash x \in A \land P(x)}{4}}

  \proofstepwide[1,2]{}{\leftrightarrow}{a\notin F(a)}{%
    \FormulaRefAuto{Q\vdash (P\leftrightarrow Q\land R)\leftrightarrow (P\leftrightarrow R)}{1}}

  \proofstepwide[1,2]{a\in F(a)}{\leftrightarrow}{a\notin F(a)}{\rChain{4,6}}
\end{tabproofwide}


% ------------------------------------------------------------
% Hauptsatz: Cantor (keine Surjektion A -> P(A))
% ------------------------------------------------------------
\FormulaThmDeltaK[Cantors Satz über die Potenzmenge]{%
  F\colon A\to \powerset(A) \vdash \neg(F\colon A\sur \powerset(A))%
}{CantorNoPowerSetSurjection}{
  \DeltaRow{Mengen}{A\dsep F}
}
\begin{tabproofwide}
  \proofstepwidestar[1]{F\colon A\sur \powerset(A)}{\rA}
  \proofstepwidestar[1]{\forall y\in \powerset(A)\,\exists x\in A\,(F(x)=y)}{%
    \FormulaRefAuto{y\in B\vdash \exists x\in A\; F(x)=y}{1}}
  \proofstepwidestar[1]{\exists x\in A\,\bigl(F(x)=\{x\in A\mid x\notin F(x)\}\bigr)}{\rUE{2}}

  \proofstepwidestar[4]{a\in A\land F(a)=\{x\in A\mid x\notin F(x)\}}{\rA}
  \proofstepwidestar[4]{a\in A}{\rAEa{4}}
  \proofstepwidestar[4]{F(a)=\{x\in A\mid x\notin F(x)\}}{\rAEb{4}}

  \proofstepwide[4]{a\in F(a)}{\leftrightarrow}{a\notin F(a)}{%
    \FormulaRefAuto{a\in A \dsep F(a)=\{x\in A\mid x\notin F(x)\}\vdash a\in F(a)\leftrightarrow a\notin F(a)}{5,6}}

  \proofstepwidestar[]{\neg (a\in F(a)\leftrightarrow a\notin F(a))}{%
    \FormulaRefAuto{\neg(P\leftrightarrow\neg P)}}

  \proofstepwidestar[4]{\bot}{\rBI{7,8}}
  \proofstepwidestar[1]{\bot}{\rEE{3,4,9}}
  \proofstepwidestar[]{\neg(F\colon A\sur \powerset(A))}{\rCI{1,10}}
\end{tabproofwide}


% ------------------------------------------------------------
% Surjektion aus Injektion (bei Nichtleere)
% (Einfügen nach \subsection{Surjektivität}, vor \subsection{Bijektivität})
% ------------------------------------------------------------

\subsection{Surjektion aus Injektion bei Nichtleere}


\FormulaDefDeltaR{%
  \begin{gathered}[t]
    F\colon A\inj B \vdash
    \Gsurjfrominj{F}{a_0} \coloneqq\\
    \{\, (b,a)\in B\times A \mid\\
    (b\in F[A]\land F(a)=b)\lor\\
    (b\notin F[A]\land a=a_0)\,\}
  \end{gathered}
}{%
  F\colon A\inj B \vdash
  \Gsurjfrominj{F}{a_0} \coloneqq
  \{\, (b,a)\in B\times A \mid (b\in F[A]\land F(a)=b)\lor (b\notin F[A]\land a=a_0)\,\}
}{%
  \DeltaRow{Mengen}{A\dsep B\dsep a_0}
  \DeltaRow{Funktionen}{F}
}

\FormulaThmDeltaR{%
  F\colon A\inj B \vdash \Gsurjfrominj{F}{a_0}\subseteq B\times A
}{%
  F\colon A\inj B \vdash \Gsurjfrominj{F}{a_0}\subseteq B\times A
}{%
  \DeltaRow{Mengen}{A\dsep B\dsep a_0\dsep a\dsep b}
  \DeltaRow{Funktionen}{F}
}
\begin{tabproof}
  \proofstep{1}{F\colon A\inj B}{\rA}
  \proofstep{}{%
    \begin{gathered}[t]
      \{\, (b,a)\in B\times A \mid\\
      (b\in F[A]\land F(a)=b)\lor\\
      (b\notin F[A]\land a=a_0)\,\}\\
      \subseteq B\times A
    \end{gathered}
  }{\FormulaRefAuto{\{\,x\in A \mid P(x)\,\}\subseteq A}}
  \proofstep{1}{\Gsurjfrominj{F}{a_0}=\{(b,a)\in B\times A\mid(b\in F[A]\land F(a)=b)\lor(b\notin F[A]\land a=a_0)\}}{%
      \FormulaRefAuto{F\colon A\inj B \vdash \Gsurjfrominj{F}{a_0} \coloneqq
    \{\, (b,a)\in B\times A \mid (b\in F[A]\land F(a)=b)\lor (b\notin F[A]\land a=a_0)\,\}}{1}}
  \proofstep{1}{\Gsurjfrominj{F}{a_0}\subseteq B\times A}{%
    \rIE{3,2}}
\end{tabproof}

\FormulaThmDelta{%
  F\colon A\inj B \dsep (b,a)\in \Gsurjfrominj{F}{a_0}\vdash b\in B%
}{
  \DeltaRow{Mengen}{A\dsep B\dsep a_0\dsep a\dsep b}
  \DeltaRow{Funktionen}{F}
}
\begin{tabproof}
  \proofstep{1}{F\colon A\inj B}{\rA}
  \proofstep{2}{(b,a)\in \Gsurjfrominj{F}{a_0}}{\rA}
  \proofstep{1}{\Gsurjfrominj{F}{a_0}=\{(b,a)\in B\times A\mid(b\in F[A]\land F(a)=b)\lor(b\notin F[A]\land a=a_0)\}}{%
      \FormulaRefAuto{F\colon A\inj B \vdash \Gsurjfrominj{F}{a_0} \coloneqq
      \{\, (b,a)\in B\times A \mid (b\in F[A]\land F(a)=b)\lor (b\notin F[A]\land a=a_0)\,\}}{1}}
  \proofstep{1,2}{%
    \begin{aligned}[t]
      (b,a)\in \{\, (b,a)\in B\times A \mid {}&
        (b\in F[A]\land F(a)=b)\\
        &\lor (b\notin F[A]\land a=a_0)\,\}
    \end{aligned}
  }{%
    \rIE{3,2}}
  \proofstep{1,2}{%
    \begin{aligned}[t]
      (b,a)\in B\times A \land \bigl({}&
        (b\in F[A]\land F(a)=b)\\
        &\lor (b\notin F[A]\land a=a_0)\bigr)
    \end{aligned}
  }{%
    \FormulaRefAuto{x \in \{u \in A \mid P(u)\} \eqvdash x \in A \land P(x)}{4}}
  \proofstep{1,2}{(b,a)\in B\times A}{\rAEa{5}}
  \proofstep{1,2}{b\in B\land a\in A}{\FormulaRefAuto{(a,b)\in A\times B\eqvdash a\in A\land b\in B}{6}}
  \proofstep{1,2}{b\in B}{\rAEa{7}}
\end{tabproof}

\FormulaThmDelta{%
  F\colon A\inj B \dsep (b,a)\in \Gsurjfrominj{F}{a_0}\vdash a\in A%
}{
  \DeltaRow{Mengen}{A\dsep B\dsep a_0\dsep a\dsep b}
  \DeltaRow{Funktionen}{F}
}
\begin{tabproof}
  \proofstep{1}{F\colon A\inj B}{\rA}
  \proofstep{2}{(b,a)\in \Gsurjfrominj{F}{a_0}}{\rA}
  \proofstep{1}{\Gsurjfrominj{F}{a_0}=\{(b,a)\in B\times A\mid(b\in F[A]\land F(a)=b)\lor(b\notin F[A]\land a=a_0)\}}{%
      \FormulaRefAuto{F\colon A\inj B \vdash \Gsurjfrominj{F}{a_0} \coloneqq
      \{\, (b,a)\in B\times A \mid (b\in F[A]\land F(a)=b)\lor (b\notin F[A]\land a=a_0)\,\}}{1}}
  \proofstep{1,2}{%
    \begin{aligned}[t]
      (b,a)\in \{\, (b,a)\in B\times A \mid {}&
        (b\in F[A]\land F(a)=b)\\
        &\lor (b\notin F[A]\land a=a_0)\,\}
    \end{aligned}
  }{%
    \rIE{3,2}}
  \proofstep{1,2}{%
    \begin{aligned}[t]
      (b,a)\in B\times A \land \bigl({}&
        (b\in F[A]\land F(a)=b)\\
        &\lor (b\notin F[A]\land a=a_0)\bigr)
    \end{aligned}
  }{%
    \FormulaRefAuto{x \in \{u \in A \mid P(u)\} \eqvdash x \in A \land P(x)}{4}}
  \proofstep{1,2}{(b,a)\in B\times A}{\rAEa{5}}
  \proofstep{1,2}{b\in B\land a\in A}{\FormulaRefAuto{(a,b)\in A\times B\eqvdash a\in A\land b\in B}{6}}
  \proofstep{1,2}{a\in A}{\rAEb{7}}
\end{tabproof}

\FormulaThmDelta{%
  F\colon A\inj B \dsep (b,a)\in \Gsurjfrominj{F}{a_0}\vdash
  (b\in F[A]\land F(a)=b)\lor (b\notin F[A]\land a=a_0)%
}{
  \DeltaRow{Mengen}{A\dsep B\dsep a_0\dsep a\dsep b}
  \DeltaRow{Funktionen}{F}
}
\begin{tabproof}
  \proofstep{1}{F\colon A\inj B}{\rA}
  \proofstep{2}{(b,a)\in \Gsurjfrominj{F}{a_0}}{\rA}
  \proofstep{1}{\Gsurjfrominj{F}{a_0}=\{(b,a)\in B\times A\mid(b\in F[A]\land F(a)=b)\lor(b\notin F[A]\land a=a_0)\}}{%
      \FormulaRefAuto{F\colon A\inj B \vdash \Gsurjfrominj{F}{a_0} \coloneqq
      \{\, (b,a)\in B\times A \mid (b\in F[A]\land F(a)=b)\lor (b\notin F[A]\land a=a_0)\,\}}{1}}
  \proofstep{1,2}{%
    \begin{aligned}[t]
      (b,a)\in \{\, (b,a)\in B\times A \mid {}&
        (b\in F[A]\land F(a)=b)\\
        &\lor (b\notin F[A]\land a=a_0)\,\}
    \end{aligned}
  }{%
    \rIE{3,2}}
  \proofstep{1,2}{%
    \begin{aligned}[t]
      (b,a)\in B\times A \land \bigl({}&
        (b\in F[A]\land F(a)=b)\\
        &\lor (b\notin F[A]\land a=a_0)\bigr)
    \end{aligned}
  }{%
    \FormulaRefAuto{x \in \{u \in A \mid P(u)\} \eqvdash x \in A \land P(x)}{4}}
  \proofstep{1,2}{(b\in F[A]\land F(a)=b)\lor (b\notin F[A]\land a=a_0)}{\rAEb{5}}
\end{tabproof}

\FormulaThmDeltaR{%
  \begin{gathered}[t]
    F\colon A\inj B \dsep b\in B \dsep a\in A \dsep\\
    \bigl((b\in F[A]\land F(a)=b)\lor\\
    (b\notin F[A]\land a=a_0)\bigr)\\
    \vdash (b,a)\in \Gsurjfrominj{F}{a_0}
  \end{gathered}
}{%
  F\colon A\inj B \dsep b\in B \dsep a\in A \dsep
  \bigl((b\in F[A]\land F(a)=b)\lor (b\notin F[A]\land a=a_0)\bigr)
  \vdash (b,a)\in \Gsurjfrominj{F}{a_0}
}{%
  \DeltaRow{Mengen}{A\dsep B\dsep a_0\dsep a\dsep b}
  \DeltaRow{Funktionen}{F}
}
\begin{tabproof}
  \proofstep{1}{F\colon A\inj B}{\rA}
  \proofstep{2}{b\in B}{\rA}
  \proofstep{3}{a\in A}{\rA}
  \proofstep{4}{%
    \begin{gathered}[t]
      (b\in F[A]\land F(a)=b)\lor\\
      (b\notin F[A]\land a=a_0)
    \end{gathered}
  }{\rA}
  \proofstep{2,3}{(b,a)\in B\times A}{\FormulaRefAuto{a\in A\dsep b\in B\vdash (a,b)\in A\times B}{2,3}}
  \proofstep{2,3,4}{%
    \begin{gathered}[t]
      (b,a)\in B\times A \land\\
      \bigl((b\in F[A]\land F(a)=b)\lor\\
      (b\notin F[A]\land a=a_0)\bigr)
    \end{gathered}
  }{\rAI{5,4}}
  \proofstep{2,3,4}{%
    \begin{gathered}[t]
      (b,a)\in \{\, (b,a)\in B\times A \mid\\
      (b\in F[A]\land F(a)=b)\lor\\
      (b\notin F[A]\land a=a_0)\,\}
    \end{gathered}
  }{%
    \FormulaRefAuto{x \in \{u \in A \mid P(u)\} \eqvdash x \in A \land P(x)}{6}}
  \proofstep{1}{\Gsurjfrominj{F}{a_0}=\{(b,a)\in B\times A\mid(b\in F[A]\land F(a)=b)\lor(b\notin F[A]\land a=a_0)\}}{%
      \FormulaRefAuto{F\colon A\inj B \vdash \Gsurjfrominj{F}{a_0} \coloneqq
    \{\, (b,a)\in B\times A \mid (b\in F[A]\land F(a)=b)\lor (b\notin F[A]\land a=a_0)\,\}}{1}}
  \proofstep{1,2,3,4}{(b,a)\in \Gsurjfrominj{F}{a_0}}{%
    \rIE{8,7}}
\end{tabproof}

\FormulaThmDelta{%
  F\colon A\inj B \dsep b\in F[A] \dsep (b,a)\in \Gsurjfrominj{F}{a_0}\vdash F(a)=b%
}{
  \DeltaRow{Mengen}{A\dsep B\dsep a_0\dsep a\dsep b}
  \DeltaRow{Funktionen}{F}
}
\begin{tabproof}
  \proofstep{1}{F\colon A\inj B}{\rA}
  \proofstep{2}{b\in F[A]}{\rA}
  \proofstep{3}{(b,a)\in \Gsurjfrominj{F}{a_0}}{\rA}
  \proofstep{1,3}{(b\in F[A]\land F(a)=b)\lor (b\notin F[A]\land a=a_0)}{%
    \FormulaRefAuto{F\colon A\inj B \dsep (b,a)\in \Gsurjfrominj{F}{a_0}\vdash (b\in F[A]\land F(a)=b)\lor (b\notin F[A]\land a=a_0)}{1,3}}
  \proofstep{5}{b\notin F[A]\land a=a_0}{\rA}
  \proofstep{5}{b\notin F[A]}{\rAEa{5}}
  \proofstep{2,5}{\bot}{\rBI{2,6}}
  \proofstep{2}{\neg(b\notin F[A]\land a=a_0)}{\rCI{5,7}}
  \proofstep{1,2,3}{(b\in F[A]\land F(a)=b)}{%
    \FormulaRefAuto{P \lor Q, \neg Q \vdash P}{4,8}}
  \proofstep{1,2,3}{F(a)=b}{\rAEb{9}}
\end{tabproof}

\FormulaThmDelta{%
  F\colon A\inj B \dsep b\notin F[A] \dsep (b,a)\in \Gsurjfrominj{F}{a_0}\vdash a=a_0%
}{
  \DeltaRow{Mengen}{A\dsep B\dsep a_0\dsep a\dsep b}
  \DeltaRow{Funktionen}{F}
}
\begin{tabproof}
  \proofstep{1}{F\colon A\inj B}{\rA}
  \proofstep{2}{b\notin F[A]}{\rA}
  \proofstep{3}{(b,a)\in \Gsurjfrominj{F}{a_0}}{\rA}
  \proofstep{1,3}{(b\in F[A]\land F(a)=b)\lor (b\notin F[A]\land a=a_0)}{%
    \FormulaRefAuto{F\colon A\inj B \dsep (b,a)\in \Gsurjfrominj{F}{a_0}\vdash (b\in F[A]\land F(a)=b)\lor (b\notin F[A]\land a=a_0)}{1,3}}
  \proofstep{5}{b\in F[A]\land F(a)=b}{\rA}
  \proofstep{5}{b\in F[A]}{\rAEa{5}}
  \proofstep{2,5}{\bot}{\rBI{6,2}}
  \proofstep{2}{\neg(b\in F[A]\land F(a)=b)}{\rCI{5,7}}
  \proofstep{1,2,3}{(b\notin F[A]\land a=a_0)}{%
    \FormulaRefAuto{P \lor Q, \neg P \vdash Q}{4,8}}
  \proofstep{1,2,3}{a=a_0}{\rAEb{9}}
\end{tabproof}

\FormulaThmDelta{%
  F\colon A\inj B \dsep (b,a)\in \Gsurjfrominj{F}{a_0} \dsep (b,c)\in \Gsurjfrominj{F}{a_0}\vdash a=c%
}{
  \DeltaRow{Mengen}{A\dsep B\dsep a_0\dsep a\dsep b\dsep c}
  \DeltaRow{Funktionen}{F}
}
\begin{tabproof}
  \proofstep{1}{F\colon A\inj B}{\rA}
  \proofstep{2}{(b,a)\in \Gsurjfrominj{F}{a_0}}{\rA}
  \proofstep{3}{(b,c)\in \Gsurjfrominj{F}{a_0}}{\rA}
  \proofstep{1,2}{a\in A}{%
    \FormulaRefAuto{F\colon A\inj B \dsep (b,a)\in \Gsurjfrominj{F}{a_0}\vdash a\in A}{1,2}}
  \proofstep{1,3}{c\in A}{%
    \FormulaRefAuto{F\colon A\inj B \dsep (b,a)\in \Gsurjfrominj{F}{a_0}\vdash a\in A}{1,3}}
  \proofstep{}{b\in F[A]\lor b\notin F[A]}{\FormulaRefAuto{P \lor \neg P}}
  \proofstep{7}{b\in F[A]}{\rA}
  \proofstep{1,2,7}{F(a)=b}{\FormulaRefAuto{F\colon A\inj B \dsep b\in F[A] \dsep (b,a)\in \Gsurjfrominj{F}{a_0}\vdash F(a)=b}{1,7,2}}
  \proofstep{1,3,7}{F(c)=b}{\FormulaRefAuto{F\colon A\inj B \dsep b\in F[A] \dsep (b,a)\in \Gsurjfrominj{F}{a_0}\vdash F(a)=b}{1,7,3}}
  \proofstep{1,2,3,7}{F(a)=F(c)}{\FormulaRefAuto{a=b,\, c=b \vdash a=c}{8,9}}
  \proofstep{1,2,3,7}{a=c}{%
    \FormulaRefAuto{F\colon A\inj B\dsep x\in A\dsep y\in A\dsep F(x)=F(y)\vdash x=y}{1,4,5,10}}
  \proofstep{12}{b\notin F[A]}{\rA}
  \proofstep{1,2,12}{a=a_0}{\FormulaRefAuto{F\colon A\inj B \dsep b\notin F[A] \dsep (b,a)\in \Gsurjfrominj{F}{a_0}\vdash a=a_0}{1,12,2}}
  \proofstep{1,3,12}{c=a_0}{\FormulaRefAuto{F\colon A\inj B \dsep b\notin F[A] \dsep (b,a)\in \Gsurjfrominj{F}{a_0}\vdash a=a_0}{1,12,3}}
  \proofstep{1,2,3,12}{a=c}{\FormulaRefAuto{a=b,\, c=b \vdash a=c}{13,14}}
  \proofstep{1,2,3}{a=c}{\rOE{6,7,11,12,15}}
\end{tabproof}

\FormulaThmDeltaR{%
  \begin{gathered}[t]
    F\colon A\inj B \dsep a_0\in A \dsep b\in B \vdash{}\\
    \exists a\in A\,\bigl((b,a)\in \Gsurjfrominj{F}{a_0}\bigr)
  \end{gathered}
}{%
  F\colon A\inj B \dsep a_0\in A \dsep b\in B \vdash \exists a\in A\,\bigl((b,a)\in \Gsurjfrominj{F}{a_0}\bigr)
}{%
  \DeltaRow{Mengen}{A\dsep B\dsep a_0\dsep a\dsep b\dsep x}
  \DeltaRow{Funktionen}{F}
}
\begin{tabproofsplitwide}
  \proofpartwideindR[Fall \(b\in F[A]\)]{%
    F\colon A\inj B \dsep a_0\in A \dsep b\in B \dsep b\in F[A]\vdash
    \exists a\in A\,\bigl((b,a)\in \Gsurjfrominj{F}{a_0}\bigr)
  }{%
    F\colon A\inj B \dsep a_0\in A \dsep b\in B \dsep b\in F[A]\vdash
    \exists a\in A\,\bigl((b,a)\in \Gsurjfrominj{F}{a_0}\bigr)
  }
    \proofstepwidestar[1]{F\colon A\inj B}{\rA}
    \proofstepwidestar[2]{a_0\in A}{\rA}
    \proofstepwidestar[3]{b\in B}{\rA}
    \proofstepwidestar[4]{b\in F[A]}{\rA}
    \proofstepwidestar[1]{F\colon A\to B}{\FormulaRefAuto{Injektive Funktion}{1}}
    \proofstepwidestar[1,4]{\exists x\in A\,b=F(x)}{%
      \FormulaRefAuto{F\colon A\to B\dsep y\in F[A]\vdash \exists x\in A\,y=F(x)}{5,4}}

    \proofstepwidestar[7]{x\in A\land b=F(x)}{\rA}
    \proofstepwidestar[7]{x\in A}{\rAEa{7}}
    \proofstepwidestar[7]{b=F(x)}{\rAEb{7}}
    \proofstepwidestar[7]{F(x)=b}{\FormulaRefAuto{a=b\vdash b=a}{9}}
    \proofstepwidestar[4,7]{b\in F[A]\land F(x)=b}{\rAI{4,10}}
    \proofstepwidestar[4,7]{%
      \begin{gathered}[t]
        (b\in F[A]\land F(x)=b)\lor\\
        (b\notin F[A]\land x=a_0)
      \end{gathered}
    }{\rOIa{11}}
    \proofstepwidestar[1,3,7]{(b,x)\in \Gsurjfrominj{F}{a_0}}{%
      \FormulaRefAuto{F\colon A\inj B \dsep b\in B \dsep a\in A \dsep \bigl((b\in F[A]\land F(a)=b)\lor (b\notin F[A]\land a=a_0)\bigr)\vdash (b,a)\in \Gsurjfrominj{F}{a_0}}{1,3,8,12}}
    \proofstepwidestar[3,4,7]{\exists a\in A\,((b,a)\in \Gsurjfrominj{F}{a_0})}{\rEI{\rAI{8,13}}}

    \proofstepwidestar[1,3,4]{\exists a\in A\,((b,a)\in \Gsurjfrominj{F}{a_0})}{\rEE{6,7,14}}
  \closeproofpartwideindR

  \proofpartwideindR[Fall \(b\notin F[A]\)]{%
    F\colon A\inj B \dsep a_0\in A \dsep b\in B \dsep b\notin F[A]\vdash
    \exists a\in A\,\bigl((b,a)\in \Gsurjfrominj{F}{a_0}\bigr)
  }{%
    F\colon A\inj B \dsep a_0\in A \dsep b\in B \dsep b\notin F[A]\vdash
    \exists a\in A\,\bigl((b,a)\in \Gsurjfrominj{F}{a_0}\bigr)
  }
    \proofstepwidestar[1]{F\colon A\inj B}{\rA}
    \proofstepwidestar[2]{a_0\in A}{\rA}
    \proofstepwidestar[3]{b\in B}{\rA}
    \proofstepwidestar[4]{b\notin F[A]}{\rA}
    \proofstepwidestar[]{a_0=a_0}{\rII}
    \proofstepwidestar[2,4]{b\notin F[A]\land a_0=a_0}{\rAI{4,5}}
    \proofstepwidestar[2,4]{%
      \begin{gathered}[t]
        (b\in F[A]\land F(a_0)=b)\lor\\
        (b\notin F[A]\land a_0=a_0)
      \end{gathered}
    }{\rOIb{6}}
    \proofstepwidestar[1,2,3,4]{(b,a_0)\in \Gsurjfrominj{F}{a_0}}{%
      \FormulaRefAuto{F\colon A\inj B \dsep b\in B \dsep a\in A \dsep \bigl((b\in F[A]\land F(a)=b)\lor (b\notin F[A]\land a=a_0)\bigr)\vdash (b,a)\in \Gsurjfrominj{F}{a_0}}{1,3,2,7}}
    \proofstepwidestar[2,4]{\exists a\in A\,((b,a)\in \Gsurjfrominj{F}{a_0})}{\rEI{\rAI{2,8}}}
  \closeproofpartwideindR

  \proofpartwide{Schluss}
    \proofstepwidestar[1]{F\colon A\inj B}{\rA}
    \proofstepwidestar[2]{a_0\in A}{\rA}
    \proofstepwidestar[3]{b\in B}{\rA}
    \proofstepwidestar[]{b\in F[A]\lor b\notin F[A]}{\FormulaRefAuto{P \lor \neg P}}

    \proofstepwidestar[5]{b\in F[A]}{\rA}
    \proofstepwidestar[1,2,3,5]{\exists a\in A\,((b,a)\in \Gsurjfrominj{F}{a_0})}{%
      \FormulaRefAuto{F\colon A\inj B \dsep a_0\in A \dsep b\in B \dsep b\in F[A]\vdash \exists a\in A\,\bigl((b,a)\in \Gsurjfrominj{F}{a_0}\bigr)}{1,2,3,5}}

    \proofstepwidestar[7]{b\notin F[A]}{\rA}
    \proofstepwidestar[1,2,3,7]{\exists a\in A\,((b,a)\in \Gsurjfrominj{F}{a_0})}{%
      \FormulaRefAuto{F\colon A\inj B \dsep a_0\in A \dsep b\in B \dsep b\notin F[A]\vdash \exists a\in A\,\bigl((b,a)\in \Gsurjfrominj{F}{a_0}\bigr)}{1,2,3,7}}

    \proofstepwidestar[1,2,3]{\exists a\in A\,((b,a)\in \Gsurjfrominj{F}{a_0})}{\rOE{4,5,6,7,8}}
  \closeproofpartwide
\end{tabproofsplitwide}

\FormulaThmDelta{%
  F\colon A\inj B \dsep a_0\in A \vdash \Gsurjfrominj{F}{a_0}\colon B\to A%
}{
  \DeltaRow{Mengen}{A\dsep B\dsep a_0\dsep a\dsep b\dsep c}
  \DeltaRow{Funktionen}{F}
}
\begin{tabproofsplit}
\proofpart{Trägerrelation}
\proofstep{1}{F\colon A\inj B}{\rA}
  \proofstep{2}{a_0\in A}{\rA}
  \proofstep{1}{\Gsurjfrominj{F}{a_0}\subseteq B\times A}{%
    \FormulaRefAuto{F\colon A\inj B \vdash \Gsurjfrominj{F}{a_0}\subseteq B\times A}{1}}

\closeproofpart

\proofpart{Totalität}
  \setcounter{proofstepnr}{3}% Fortlaufende Schritte dieses Beweises.
\proofstep{4}{b\in B}{\rA}
  \proofstep{1,2,4}{\exists a\in A\,((b,a)\in \Gsurjfrominj{F}{a_0})}{%
    \FormulaRefAuto{F\colon A\inj B \dsep a_0\in A \dsep b\in B \vdash \exists a\in A\,\bigl((b,a)\in \Gsurjfrominj{F}{a_0}\bigr)}{1,2,4}}
  \proofstep{1,2}{b\in B\rightarrow \exists a\in A\,((b,a)\in \Gsurjfrominj{F}{a_0})}{\rRI{4,5}}

\closeproofpart

\proofpart{Funktionale Eindeutigkeit}
  \setcounter{proofstepnr}{6}% Fortlaufende Schritte dieses Beweises.
\proofstep{7}{(b,a)\in \Gsurjfrominj{F}{a_0}\land (b,c)\in \Gsurjfrominj{F}{a_0}}{\rA}
  \proofstep{7}{(b,a)\in \Gsurjfrominj{F}{a_0}}{\rAEa{7}}
  \proofstep{7}{(b,c)\in \Gsurjfrominj{F}{a_0}}{\rAEb{7}}
  \proofstep{1,7}{a=c}{%
    \FormulaRefAuto{F\colon A\inj B \dsep (b,a)\in \Gsurjfrominj{F}{a_0} \dsep (b,c)\in \Gsurjfrominj{F}{a_0}\vdash a=c}{1,8,9}}
  \proofstep{1}{((b,a)\in \Gsurjfrominj{F}{a_0}\land (b,c)\in \Gsurjfrominj{F}{a_0})\rightarrow a=c}{\rRI{7,10}}

\closeproofpart

\proofpart{Zusammenführung}
  \setcounter{proofstepnr}{11}% Fortlaufende Schritte dieses Beweises.
\proofstep{1,2}{\Gsurjfrominj{F}{a_0}\colon B\to A}{\FormulaRefAuto{Funktion}{3,6,11}}

\closeproofpart
\end{tabproofsplit}

\FormulaThmDelta{%
  F\colon A\inj B \dsep a_0\in A \vdash \Gsurjfrominj{F}{a_0}\colon B\sur A%
}{
  \DeltaRow{Mengen}{A\dsep B\dsep a_0\dsep b\dsep y\dsep x}
  \DeltaRow{Funktionen}{F}
}
\begin{tabproofsplit}
\proofpart{Funktionstyp}
\proofstep{1}{F\colon A\inj B}{\rA}
  \proofstep{2}{a_0\in A}{\rA}
  \proofstep{1}{F\colon A\to B}{\FormulaRefAuto{Injektive Funktion}{1}}
  \proofstep{1,2}{\Gsurjfrominj{F}{a_0}\colon B\to A}{%
    \FormulaRefAuto{F\colon A\inj B \dsep a_0\in A \vdash \Gsurjfrominj{F}{a_0}\colon B\to A}{1,2}}

\closeproofpart

\proofpart{Surjektivität}
  \setcounter{proofstepnr}{4}% Fortlaufende Schritte dieses Beweises.
\proofstep{5}{y\in A}{\rA}
  \proofstep{1,5}{(y,F(y))\in F}{%
    \FormulaRefAuto{F\colon A\to B\dsep x\in A\vdash (x,F(x))\in F}{3,5}}
  \proofstep{1,5}{F(y)\in B}{%
    \FormulaRefAuto{F\colon A\to B\dsep (x,y)\in F\vdash y\in B}{3,6}}
  \proofstep{}{F(y)=F(y)}{\rII}
  \proofstep{5}{y\in A\land F(y)=F(y)}{\rAI{5,8}}
  \proofstep{5}{\exists x\in A\,F(y)=F(x)}{\rEI{9}}
  \proofstep{1,5}{F(y)\in F[A]}{%
    \FormulaRefAuto{F\colon A\to B\dsep \exists x\in A\,y=F(x)\vdash y\in F[A]}{3,10}}
  \proofstep{1,5}{F(y)\in F[A]\land F(y)=F(y)}{\rAI{11,8}}
  \proofstep{1,5}{%
    \begin{aligned}[t]
      &(F(y)\in F[A]\land F(y)=F(y))\\
      &\lor (F(y)\notin F[A]\land y=a_0)
    \end{aligned}
  }{\rOIa{12}}
  \proofstep{1,5}{(F(y),y)\in \Gsurjfrominj{F}{a_0}}{%
    \FormulaRefAuto{F\colon A\inj B \dsep b\in B \dsep a\in A \dsep \bigl((b\in F[A]\land F(a)=b)\lor (b\notin F[A]\land a=a_0)\bigr)
    \vdash (b,a)\in \Gsurjfrominj{F}{a_0}}{1,7,5,13}}
  \proofstep{1,2,5}{\Gsurjfrominj{F}{a_0}(F(y))=y}{%
    \FormulaRefAuto{F\colon A\to B\dsep (x,y)\in F\vdash F(x)=y}{4,14}}
  \proofstep{1,2,5}{\exists b\in B\,\bigl(\Gsurjfrominj{F}{a_0}(b)=y\bigr)}{\rEI{\rAI{7,15}}}
  \proofstep{1,2}{y\in A\rightarrow \exists b\in B\,\bigl(\Gsurjfrominj{F}{a_0}(b)=y\bigr)}{\rRI{5,16}}

\closeproofpart

\proofpart{Zusammenführung}
  \setcounter{proofstepnr}{17}% Fortlaufende Schritte dieses Beweises.
\proofstep{1,2}{\Gsurjfrominj{F}{a_0}\colon B\sur A}{\FormulaRefAuto{Surjektive Funktion}{4,17}}

\closeproofpart
\end{tabproofsplit}

\FormulaThmDelta{%
  F\colon A\inj B \dsep A\neq \varnothing \vdash \exists G\colon B\sur A%
}{
  \DeltaRow{Mengen}{A\dsep B}
  \DeltaRow{Funktionen}{F\dsep G}
}
\begin{tabproof}
  \proofstep{1}{F\colon A\inj B}{\rA}
  \proofstep{2}{A\neq \varnothing}{\rA}
  \proofstep{2}{\exists a_0\,(a_0\in A)}{\FormulaRefAuto{A\neq\varnothing \eqvdash \exists x\,(x \in A)}{2}}
  \proofstep{4}{a_0\in A}{\rA}
  \proofstep{1,4}{\Gsurjfrominj{F}{a_0}\colon B\sur A}{\FormulaRefAuto{F\colon A\inj B \dsep a_0\in A \vdash \Gsurjfrominj{F}{a_0}\colon B\sur A}{1,4}}
  \proofstep{1,4}{\exists G\colon B\sur A}{\rEI{5}}
  \proofstep{1,2}{\exists G\colon B\sur A}{\rEE{3,4,6}}
\end{tabproof}

\FormulaThmDelta[Linksinverses Verhalten von \(\Gsurjfrominj{F}{a_0}\) an \(y\)]{%
  F\colon A\inj B \dsep a_0\in A \dsep y\in A \vdash \Gsurjfrominj{F}{a_0}(F(y))=y
}{%
  \DeltaRow{Mengen}{A\dsep B\dsep a_0\dsep y\dsep x}
  \DeltaRow{Funktionen}{F}
}
\begin{tabproof}
  \proofstep{1}{F\colon A\inj B}{\rA}
  \proofstep{2}{a_0\in A}{\rA}
  \proofstep{3}{y\in A}{\rA}
  \proofstep{1}{F\colon A\to B}{\FormulaRefAuto{Injektive Funktion}{1}}
  \proofstep{1,2}{\Gsurjfrominj{F}{a_0}\colon B\to A}{%
    \FormulaRefAuto{F\colon A\inj B \dsep a_0\in A \vdash \Gsurjfrominj{F}{a_0}\colon B\to A}{1,2}}
  \proofstep{1,3}{F(y)\in B}{%
    \FormulaRefAuto{F\colon A\to B\dsep x\in A\vdash F(x)\in B}{4,3}}
  \proofstep{}{F(y)=F(y)}{\rII}
  \proofstep{3}{y\in A\land F(y)=F(y)}{\rAI{3,7}}
  \proofstep{3}{\exists x\in A\,F(y)=F(x)}{\rEI{8}}
  \proofstep{1,3}{F(y)\in F[A]}{%
    \FormulaRefAuto{F\colon A\to B\dsep \exists x\in A\,z=F(x)\vdash z\in F[A]}{4,9}}
  \proofstep{1,3}{%
    \begin{gathered}[t]
      (F(y)\in F[A]\land F(y)=F(y))\lor\\
      (F(y)\notin F[A]\land y=a_0)
    \end{gathered}
  }{\rOIa{\rAI{10,7}}}
  \proofstep{1,3}{(F(y),y)\in \Gsurjfrominj{F}{a_0}}{%
    \FormulaRefAuto{F\colon A\inj B \dsep b\in B \dsep a\in A \dsep \bigl((b\in F[A]\land F(a)=b)\lor (b\notin F[A]\land a=a_0)\bigr)\vdash (b,a)\in \Gsurjfrominj{F}{a_0}}{1,6,3,11}}
  \proofstep{1,2,3}{\Gsurjfrominj{F}{a_0}(F(y))=y}{%
    \FormulaRefAuto{F\colon A\to B\dsep (x,y)\in F\vdash F(x)=y}{5,12}}
\end{tabproof}

\FormulaThmDelta[Dualer Hilfssatz bei Nichtleere]{%
  F\colon A\inj B \dsep A\neq \varnothing \vdash \exists H\colon B\sur A\,\forall y\in A\,\bigl(H(F(y))=y\bigr)
}{%
  \DeltaRow{Mengen}{A\dsep B\dsep a_0\dsep y}
  \DeltaRow{Funktionen}{F\dsep H}
}
\begin{tabproof}
  \proofstep{1}{F\colon A\inj B}{\rA}
  \proofstep{2}{A\neq \varnothing}{\rA}
  \proofstep{2}{\exists a_0\,(a_0\in A)}{\FormulaRefAuto{A\neq\varnothing \eqvdash \exists x\,(x \in A)}{2}}
  \proofstep{4}{a_0\in A}{\rA}
  \proofstep{1,4}{\Gsurjfrominj{F}{a_0}\colon B\sur A}{%
    \FormulaRefAuto{F\colon A\inj B \dsep a_0\in A \vdash \Gsurjfrominj{F}{a_0}\colon B\sur A}{1,4}}
  \proofstep{6}{y\in A}{\rA}
  \proofstep{1,4,6}{\Gsurjfrominj{F}{a_0}(F(y))=y}{%
    \FormulaRefAuto{F\colon A\inj B \dsep a_0\in A \dsep y\in A \vdash \Gsurjfrominj{F}{a_0}(F(y))=y}{1,4,6}}
  \proofstep{1,4}{\forall y\in A\,\bigl(\Gsurjfrominj{F}{a_0}(F(y))=y\bigr)}{\rUI{\rRI{6,7}}}
  \proofstep{1,4}{%
    \Gsurjfrominj{F}{a_0}\colon B\sur A \land
    \forall y\in A\,\bigl(\Gsurjfrominj{F}{a_0}(F(y))=y\bigr)
  }{\rAI{5,8}}
  \proofstep{1,4}{\exists H\colon B\sur A\,\forall y\in A\,\bigl(H(F(y))=y\bigr)}{\rEI{9}}
  \proofstep{1,2}{\exists H\colon B\sur A\,\forall y\in A\,\bigl(H(F(y))=y\bigr)}{\rEE{3,4,10}}
\end{tabproof}

\subsection{Urbildmengen unter surjektiven Funktionen}

\FormulaThmDelta[Urbilder liegen im Definitionsbereich surjektiver Funktionen]{%
  C\subseteq B\dsep F\colon A\sur B \vdash F^{-1}[C]\subseteq A
}{
  \DeltaRow{Mengen}{A\dsep B\dsep C}
  \DeltaRow{Funktionen}{F}
}
\begin{tabproof}
  \proofstep{1}{C\subseteq B}{\rA}
  \proofstep{2}{F\colon A\sur B}{\rA}
  \proofstep{2}{F\colon A\to B}{\FormulaRefAuto{Surjektive Funktion}{2}}
  \proofstep{1,2}{F^{-1}[C]\subseteq A}{%
    \FormulaRefAuto{F\colon A\to B\dsep C\subseteq B \vdash F^{-1}[C]\subseteq A}{3,1}}
\end{tabproof}


\chapter{Beispiele und Konstruktionen}

\section{Projektionen}

\subsection{Die Projektionen auf die Komponenten}

\subsubsection{Projektion auf die erste Komponente}

\FormulaDefDeltaK[Projektion auf die erste Komponente]{%
  \begin{aligned}[t]
    &\pi_1\colon A\times B\to A,\\[-2pt]
    &\forall (x,y)\in A\times B\,
      \bigl(\pi_1(x,y)\coloneqq x\bigr)
  \end{aligned}%
}{FirstProjectionDef}{%
  \DeltaRow{Mengen}{A\dsep B}
  \DeltaRow{Funktionen}{\pi_1}
}
\begin{tabproof}
  \proofstep{1}{(x,y)\in A\times B}{\rA}
  \proofstep{1}{x\in A\land y\in B}{%
    \FormulaRefAuto{(a,b)\in A\times B\eqvdash a\in A\land b\in B}[thm]{1}}
  \proofstep{1}{x\in A}{\rAEa{2}}
\end{tabproof}

\FormulaThmDelta[Projektion auf die erste Komponente]
{\pi_1\colon A\times B\to A}
{
  \DeltaRow{Mengen}{A\dsep B}
}
\begin{tabproof}
  \proofstep{}{\pi_1\colon A\times B\to A}{%
    \FormulaRefAuto{FirstProjectionDef}[def]}
\end{tabproof}

\FormulaThmDelta
{\TotRel{\pi_1,A\times B,A}}
{
  \DeltaRow{Mengen}{A\dsep B}
}
\begin{tabproof}
    \proofstep{}{\pi_1\colon A\times B\to A}{\FormulaRefAuto{\pi_1\colon A\times B\to A}}
    \proofstep{}{\TotRel{\pi_1,A\times B,A}}{\FormulaRefAuto{F\colon A\to B\vdash\TotRel{F,A,B}}{1}}
\end{tabproof}

\FormulaThmDelta[Grundgleichung der Projektion auf die erste Komponente]
{(x,y)\in A\times B \vdash \pi_1(x,y)=x}
{
  \DeltaRow{Mengen}{A\dsep B\dsep x\dsep y}
}
\begin{tabproof}
  \proofstep{1}{(x,y)\in A\times B}{\rA}
  \proofstep{}{\forall (u,v)\in A\times B\,
    \bigl(\pi_1(u,v)=u\bigr)}{%
    \FormulaRefAuto{FirstProjectionDef}[def]}
  \proofstep{1}{\pi_1(x,y)=x}{%
    \rRE{\rUE{2},1}}
\end{tabproof}

\FormulaThmDelta
{x\in A\dsep y\in B \vdash \pi_1(x,y)=x}
{
  \DeltaRow{Mengen}{A\dsep B\dsep x\dsep y}
}
\begin{tabproof}
  \proofstep{1}{x\in A}{\rA}
  \proofstep{2}{y\in B}{\rA}
  \proofstep{1,2}{(x,y)\in A\times B}{%
    \FormulaRefAuto{(a,b)\in A\times B\eqvdash a\in A\land b\in B}{1,2}}
  \proofstep{1,2}{\pi_1(x,y)=x}{%
    \FormulaRefAuto{(x,y)\in A\times B\vdash \pi_1(x,y)=x}{3}}
\end{tabproof}


\subsubsection{Projektion auf die zweite Komponente}

\FormulaDefDeltaK[Projektion auf die zweite Komponente]{%
  \begin{aligned}[t]
    &\pi_2\colon A\times B\to B,\\[-2pt]
    &\forall (x,y)\in A\times B\,
      \bigl(\pi_2(x,y)\coloneqq y\bigr)
  \end{aligned}%
}{SecondProjectionDef}{%
  \DeltaRow{Mengen}{A\dsep B}
  \DeltaRow{Funktionen}{\pi_2}
}
\begin{tabproof}
  \proofstep{1}{(x,y)\in A\times B}{\rA}
  \proofstep{1}{x\in A\land y\in B}{%
    \FormulaRefAuto{(a,b)\in A\times B\eqvdash a\in A\land b\in B}[thm]{1}}
  \proofstep{1}{y\in B}{\rAEb{2}}
\end{tabproof}

\FormulaThmDelta[Projektion auf die zweite Komponente]
{\pi_2\colon A\times B\to B}
{
  \DeltaRow{Mengen}{A\dsep B}
}
\begin{tabproof}
  \proofstep{}{\pi_2\colon A\times B\to B}{%
    \FormulaRefAuto{SecondProjectionDef}[def]}
\end{tabproof}

\FormulaThmDelta
{\TotRel{\pi_2,A\times B,B}}
{
  \DeltaRow{Mengen}{A\dsep B}
}
\begin{tabproof}
    \proofstep{}{\pi_2\colon A\times B\to B}{\FormulaRefAuto{\pi_2\colon A\times B\to B}}
    \proofstep{}{\TotRel{\pi_2,A\times B,B}}{\FormulaRefAuto{F\colon A\to B\vdash\TotRel{F,A,B}}{1}}
\end{tabproof}

\FormulaThmDelta[Grundgleichung der Projektion auf die zweite Komponente]
{(x,y)\in A\times B \vdash \pi_2(x,y)=y}
{
  \DeltaRow{Mengen}{A\dsep B\dsep x\dsep y}
}
\begin{tabproof}
  \proofstep{1}{(x,y)\in A\times B}{\rA}
  \proofstep{}{\forall (u,v)\in A\times B\,
    \bigl(\pi_2(u,v)=v\bigr)}{%
    \FormulaRefAuto{SecondProjectionDef}[def]}
  \proofstep{1}{\pi_2(x,y)=y}{%
    \rRE{\rUE{2},1}}
\end{tabproof}

\FormulaThmDelta
{x\in A\dsep y\in B \vdash \pi_2(x,y)=y}
{
  \DeltaRow{Mengen}{A\dsep B\dsep x\dsep y}
}
\begin{tabproof}
  \proofstep{1}{x\in A}{\rA}
  \proofstep{2}{y\in B}{\rA}
  \proofstep{1,2}{(x,y)\in A\times B}{%
    \FormulaRefAuto{(a,b)\in A\times B\eqvdash a\in A\land b\in B}{1,2}}
  \proofstep{1,2}{\pi_2(x,y)=y}{%
    \FormulaRefAuto{(x,y)\in A\times B \vdash \pi_2(x,y)=y}{3}}
\end{tabproof}


\subsubsection{Eigenschaften von Projektionsabbildungen}

\FormulaThmDelta[Rekonstruktion aus Projektionen]{%
  (x,y)\in A\times B \vdash (x,y) = \bigl(\pi_1(x,y),\pi_2(x,y)\bigr)
}{
  \DeltaRow{Mengen}{A\dsep B\dsep x\dsep y}
  \DeltaRow{Projektion auf erste Komponente}{\pi_1:A\times B\to A}
  \DeltaRow{Projektion auf zweite Komponente}{\pi_2:A\times B\to B}
}
\begin{tabproofwide}
  \proofstepwidestar[1]{(x,y)\in A\times B}{\rA}
  \proofstepwidestar[1]{\pi _ 1 ( x , y ) = x}{%
      \FormulaRefAuto{(x,y)\in A\times B \vdash \pi_1(x,y)=x}{1}}
  \proofstepwidestar[1]{x=\pi_1(x,y)}{%
    \FormulaRefAuto{a=b\vdash b=a}{2}}
  \proofstepwidestar[1]{\pi _ 2 ( x , y ) = y}{%
      \FormulaRefAuto{(x,y)\in A\times B \vdash \pi_2(x,y)=y}{1}}
  \proofstepwidestar[1]{y=\pi_2(x,y)}{%
    \FormulaRefAuto{a=b\vdash b=a}{4}}
  \proofstepwide[1]{(x,y)}{=}{(x,y)}{\rII}
  \proofstepwide[1]{}{=}{(\pi_1(x,y),y)}{\rIE{3,6}}
  \proofstepwide[1]{}{=}{(\pi_1(x,y),\pi_2(x,y))}{\rIE{5,7}}
  \proofstepwide[1]{(x,y)}{=}{(\pi_1(x,y),\pi_2(x,y))}{\rChain{6,8}}
\end{tabproofwide}

\FormulaThmDelta[Element des Produkts als Paar]{%
  z\in A\times B \vdash z = \bigl(\pi_1(z),\pi_2(z)\bigr)
}{
  \DeltaRow{Mengen}{A\dsep B\dsep z}
  \DeltaRow{Projektion auf erste Komponente}{\pi_1:A\times B\to A}
  \DeltaRow{Projektion auf zweite Komponente}{\pi_2:A\times B\to B}
}
\begin{tabproof}
  \proofstep{1}{z\in A\times B}{\rA}

  \proofstep{1}{\exists x\in A\,\exists y\in B\, z=(x,y)}{%
    \FormulaRefAuto{x\in A\times B\eqvdash \exists a \in A\, \exists b \in B\,x=(a,b)}{1}}

  \proofstep{3}{x\in A\land y\in B\land z=(x,y)}{\rA}

  \proofstep{3}{x\in A}{\rAEa{3}}

  \proofstep{3}{y\in B}{\FormulaRefAuto{P\land Q\land R\vdash Q}{3}}

  \proofstep{3}{z=(x,y)}{\rAEb{3}}

  \proofstep{3}{(x,y)\in A\times B}{%
    \FormulaRefAuto{(a,b)\in A\times B\eqvdash a\in A\land b\in B}{4,5}}

  \proofstep{3}{(x,y)=\bigl(\pi_1(x,y),\pi_2(x,y)\bigr)}{%
    \FormulaRefAuto{(x,y)\in A\times B \vdash (x,y)=\bigl(\pi_1(x,y),\pi_2(x,y)\bigr)}{7}}

  \proofstep{3}{z=\bigl(\pi_1(z),\pi_2(z)\bigr)}{%
    \rIE{6,8}}

  \proofstep{1}{z=\bigl(\pi_1(z),\pi_2(z)\bigr)}{%
    \rEE{2,3,9}}
\end{tabproof}

\FormulaThmDeltaK[Elementkriterium einer ausgesonderten Paarmenge]{%
  (x,y)\in\{p\in A\times B\mid P(\pi_1(p),\pi_2(p))\}
  \eqvdash x\in A\land y\in B\land P(x,y)%
}{ProductComprehensionMembership}{%
  \DeltaRow{Mengen}{A\dsep B\dsep p\dsep x\dsep y}
  \DeltaRow{Prädikate}{P}
  \DeltaRow{Projektionen}{\pi_1\dsep\pi_2}
}
\begin{tabproofwide}
  \proofstepwidestar[]{%
    \begin{aligned}[t]
      &(x,y)\in\{p\in A\times B\mid P(\pi_1(p),\pi_2(p))\}\\[-2pt]
      &\quad\leftrightarrow (x,y)\in A\times B\land
        P(\pi_1(x,y),\pi_2(x,y))
    \end{aligned}}{%
    \FormulaRefAuto{x\in\{u\in A\mid P(u)\}\eqvdash
      x\in A\land P(x)}[thm]}
  \proofstepwidestar[2]{(x,y)\in A\times B}{\rA}
  \proofstepwidestar[2]{\pi_1(x,y)=x}{%
    \FormulaRefAuto{(x,y)\in A\times B\vdash\pi_1(x,y)=x}[thm]{2}}
  \proofstepwidestar[2]{\pi_2(x,y)=y}{%
    \FormulaRefAuto{(x,y)\in A\times B\vdash\pi_2(x,y)=y}[thm]{2}}
  \proofstepwidestar[2]{%
    P(\pi_1(x,y),\pi_2(x,y))\leftrightarrow P(x,y)}{%
    \rIE{3,4,\FormulaRefAuto{P\leftrightarrow P}[thm]}}
  \proofstepwidestar[]{%
    (x,y)\in A\times B\rightarrow
    \bigl(P(\pi_1(x,y),\pi_2(x,y))\leftrightarrow P(x,y)\bigr)}{%
    \rRI{2,5}}
  \proofstepwidestar[]{(x,y)\in A\times B\leftrightarrow x\in A\land y\in B}{%
    \FormulaRefAuto{(a,b)\in A\times B\eqvdash a\in A\land b\in B}[thm]}
  \proofstepwidestar[8]{%
    (x,y)\in A\times B\land P(\pi_1(x,y),\pi_2(x,y))}{\rA}
  \proofstepwidestar[8]{%
    P(\pi_1(x,y),\pi_2(x,y))\leftrightarrow P(x,y)}{%
    \rRE{\rAEa{8},6}}
  \proofstepwidestar[8]{(x,y)\in A\times B\land P(x,y)}{%
    \rAI{\rAEa{8},\rRE{\rAEb{8},\rLREa{9}}}}
  \proofstepwidestar[11]{(x,y)\in A\times B\land P(x,y)}{\rA}
  \proofstepwidestar[11]{%
    P(\pi_1(x,y),\pi_2(x,y))\leftrightarrow P(x,y)}{%
    \rRE{\rAEa{11},6}}
  \proofstepwidestar[11]{%
    (x,y)\in A\times B\land P(\pi_1(x,y),\pi_2(x,y))}{%
    \rAI{\rAEa{11},\rRE{\rAEb{11},\rLREb{12}}}}
  \proofstepwidestar[]{\begin{aligned}[t]
    &((x,y)\in A\times B\land
      P(\pi_1(x,y),\pi_2(x,y)))\\[-2pt]
    &\quad\leftrightarrow((x,y)\in A\times B\land P(x,y))
  \end{aligned}}{\rLRI{\rRI{8,10},\rRI{11,13}}}
  \proofstepwidestar[]{%
    \begin{aligned}[t]
      &(x,y)\in\{p\in A\times B\mid P(\pi_1(p),\pi_2(p))\}\\[-2pt]
      &\quad\leftrightarrow x\in A\land y\in B\land P(x,y)
    \end{aligned}}{%
    \rLRS{\FormulaRefAuto{P\land(Q\land R)\eqvdash
      (P\land Q)\land R}[thm],\rLRS{7,\rLRS{14,1}}}}
\end{tabproofwide}

\FormulaThmDelta[Erste Komponente eines Produkteelements]{%
  z\in A\times B\vdash \pi_1(z)\in A%
}{%
  \DeltaRow{Mengen}{z\dsep A\dsep B}%
  \DeltaRow{Funktionensymbol}{\pi_1}%
}
\begin{tabproof}
  \proofstep{1}{z\in A\times B}{\rA}
  \proofstep{}{\pi_1\colon A\times B\to A}{%
    \FormulaRefAuto{\pi_1\colon A\times B\to A}}
  \proofstep{1}{\pi_1(z)\in A}{%
    \FormulaRefAuto{F\colon A\to B\dsep x\in A\vdash F(x)\in B}{2,1}}
\end{tabproof}

\FormulaThmDelta[Zweite Komponente eines Produkteelements]{%
  z\in A\times B\vdash \pi_2(z)\in B%
}{%
  \DeltaRow{Mengen}{z\dsep A\dsep B}%
  \DeltaRow{Funktionensymbol}{\pi_2}%
}
\begin{tabproof}
  \proofstep{1}{z\in A\times B}{\rA}
  \proofstep{}{\pi_2\colon A\times B\to B}{%
    \FormulaRefAuto{\pi_2\colon A\times B\to B}}
  \proofstep{1}{\pi_2(z)\in B}{%
    \FormulaRefAuto{F\colon A\to B\dsep x\in A\vdash F(x)\in B}{2,1}}
\end{tabproof}

Die Grundgleichungen der Projektionen lassen sich unmittelbar auf Terme
übertragen, die als Paare identifiziert sind.

\FormulaThmDelta[Erste Projektion einer Paaridentität]{%
  a=(x,y)\dsep (x,y)\in A\times B\vdash \pi_1(a)=x%
}{%
  \DeltaRow{Mengen}{a\dsep x\dsep y\dsep A\dsep B}%
  \DeltaRow{Funktionensymbol}{\pi_1}%
}
\begin{tabproofwide}
  \proofstepwidestar[1]{a=(x,y)}{\rA}
  \proofstepwidestar[2]{(x,y)\in A\times B}{\rA}
  \proofstepwidestar[1]{\pi_1(a)=\pi_1(a)}{\rII}
  \proofstepwidestar[1]{\pi_1(a)=\pi_1(x,y)}{\rIE{1,3}}
  \proofstepwidestar[2]{\pi_1(x,y)=x}{%
    \FormulaRefAuto{(x,y)\in A\times B \vdash \pi_1(x,y)=x}{2}}
  \proofstepwidestar[1,2]{\pi_1(a)=x}{%
    \FormulaRefAuto{a = b,\, b = c \vdash a = c}{4,5}}
\end{tabproofwide}

\FormulaThmDelta[Zweite Projektion einer Paaridentität]{%
  a=(x,y)\dsep (x,y)\in A\times B\vdash \pi_2(a)=y%
}{%
  \DeltaRow{Mengen}{a\dsep x\dsep y\dsep A\dsep B}%
  \DeltaRow{Funktionensymbol}{\pi_2}%
}
\begin{tabproofwide}
  \proofstepwidestar[1]{a=(x,y)}{\rA}
  \proofstepwidestar[2]{(x,y)\in A\times B}{\rA}
  \proofstepwidestar[1]{\pi_2(a)=\pi_2(a)}{\rII}
  \proofstepwidestar[1]{\pi_2(a)=\pi_2(x,y)}{\rIE{1,3}}
  \proofstepwidestar[2]{\pi_2(x,y)=y}{%
    \FormulaRefAuto{(x,y)\in A\times B \vdash \pi_2(x,y)=y}{2}}
  \proofstepwidestar[1,2]{\pi_2(a)=y}{%
    \FormulaRefAuto{a = b,\, b = c \vdash a = c}{4,5}}
\end{tabproofwide}

Ist das identifizierte Objekt selbst bereits als Element eines Produkts
typisiert, muss die Zugehörigkeit des dargestellten Paares nicht gesondert aus
den Komponenten hergeleitet werden.

\FormulaThmDelta[Erste Projektion eines identifizierten Produkteelements]{%
  a\in A\times B\dsep a=(x,y)\vdash \pi_1(a)=x%
}{%
  \DeltaRow{Mengen}{a\dsep x\dsep y\dsep A\dsep B}%
  \DeltaRow{Funktionensymbol}{\pi_1}%
}
\begin{tabproofwide}
  \proofstepwidestar[1]{a\in A\times B}{\rA}
  \proofstepwidestar[2]{a=(x,y)}{\rA}
  \proofstepwidestar[2]{(x,y)=a}{%
    \FormulaRefAuto{a = b \vdash b = a}{2}}
  \proofstepwidestar[1,2]{(x,y)\in A\times B}{%
    \FormulaRefAuto{a=b \dsep b\in A \vdash a\in A}{3,1}}
  \proofstepwidestar[1,2]{\pi_1(a)=x}{%
    \FormulaRefAuto{a=(x,y)\dsep (x,y)\in A\times B\vdash \pi_1(a)=x}{2,4}}
\end{tabproofwide}

\FormulaThmDelta[Zweite Projektion eines identifizierten Produkteelements]{%
  a\in A\times B\dsep a=(x,y)\vdash \pi_2(a)=y%
}{%
  \DeltaRow{Mengen}{a\dsep x\dsep y\dsep A\dsep B}%
  \DeltaRow{Funktionensymbol}{\pi_2}%
}
\begin{tabproofwide}
  \proofstepwidestar[1]{a\in A\times B}{\rA}
  \proofstepwidestar[2]{a=(x,y)}{\rA}
  \proofstepwidestar[2]{(x,y)=a}{%
    \FormulaRefAuto{a = b \vdash b = a}{2}}
  \proofstepwidestar[1,2]{(x,y)\in A\times B}{%
    \FormulaRefAuto{a=b \dsep b\in A \vdash a\in A}{3,1}}
  \proofstepwidestar[1,2]{\pi_2(a)=y}{%
    \FormulaRefAuto{a=(x,y)\dsep (x,y)\in A\times B\vdash \pi_2(a)=y}{2,4}}
\end{tabproofwide}

\subsection{Surjektivität}

Die Definitionen der Projektionen stehen bewusst in diesem Band, weil ihre
kennzeichnende Standard-Eigenschaft die Surjektivität ist. Für die erste
Projektion genügt ein bewohnter zweiter Faktor, für die zweite ein bewohnter
erster Faktor; bei leerer Zielmenge ist die Surjektivität ohnehin trivial.

\subsubsection{Erste Projektion}

\FormulaThmDelta[Surjektivität von \(\pi_1\)]{%
x\in A \dsep b_0\in B \vdash \exists z\in A\times B\; \pi_1(z)=x
}{
  \DeltaRow{Mengen}{A \dsep B \dsep x \dsep b_0 \dsep z}
}
\begin{tabproof}
  \proofstep{1}{x\in A}{\rA}
  \proofstep{2}{b_0\in B}{\rA}
  % Paarbildung
  \proofstep{1,2}{(x,b_0)\in A\times B}{%
    \FormulaRefAuto{(a,b)\in A\times B\eqvdash a\in A\land b\in B}{1,2}}
  % Grundgleichung der Projektion: (x,b0) in A×B ⇒ π1(x,b0)=x
  \proofstep{1,2}{\pi_1(x,b_0)=x}{%
    \FormulaRefAuto{(x,y)\in A\times B \vdash \pi_1(x,y)=x}{3}}
  % Existenzzeuge z = (x,b0) für Surjektivität
  \proofstep{1,2}{\exists z\in A\times B\; \pi_1(z)=x}{%
    \rEI{\rAI{3,4}}}
\end{tabproof}

\FormulaThmDeltaK[Erste Projektion als surjektive Funktion]{%
  b_0\in B \vdash \pi_1\colon A\times B\sur A%
}{FirstProjectionSurjective}{%
  \DeltaRow{Mengen}{A\dsep B\dsep b_0\dsep x\dsep z}%
}
\begin{tabproofsplit}
\proofpart{Funktionstyp}
\proofstep{1}{b_0\in B}{\rA}
  \proofstep{}{\pi_1\colon A\times B\to A}{%
    \FormulaRefAuto{\pi_1\colon A\times B\to A}}

\closeproofpart

\proofpart{Surjektivität}
  \setcounter{proofstepnr}{2}% Fortlaufende Schritte dieses Beweises.
\proofstep{3}{x\in A}{\rA}
  \proofstep{1,3}{\exists z\in A\times B\,\pi_1(z)=x}{%
    \FormulaRefAuto{%
      x\in A\dsep b_0\in B\vdash
      \exists z\in A\times B\,\pi_1(z)=x%
    }{3,1}}
  \proofstep{1}{\forall x\in A\,\exists z\in A\times B\,\pi_1(z)=x}{%
    \rUI{\rRI{3,4}}}

\closeproofpart

\proofpart{Zusammenführung}
  \setcounter{proofstepnr}{5}% Fortlaufende Schritte dieses Beweises.
\proofstep{1}{\pi_1\colon A\times B\sur A}{%
    \FormulaRefAuto{Surjektive Funktion}{2,5}}

\closeproofpart
\end{tabproofsplit}

\subsubsection{Zweite Projektion}

\FormulaThmDelta[Surjektivität von \(\pi_2\)]{%
y\in B \dsep a_0\in A \vdash \exists z\in A\times B\; \pi_2(z)=y
}{
  \DeltaRow{Mengen}{A \dsep B \dsep y \dsep a_0 \dsep z}
}
\begin{tabproof}
  \proofstep{1}{y\in B}{\rA}
  \proofstep{2}{a_0\in A}{\rA}
  % Paarbildung
  \proofstep{1,2}{(a_0,y)\in A\times B}{%
    \FormulaRefAuto{a\in A \dsep b\in B \vdash (a,b)\in A\times B}{2,1}}
  % Grundgleichung der Projektion: (a0,y) in A×B ⇒ π2(a0,y)=y
  \proofstep{1,2}{\pi_2(a_0,y)=y}{%
    \FormulaRefAuto{(x,y)\in A\times B \vdash \pi_2(x,y)=y}{3}}
  % Existenzzeuge z = (a0,y) für Surjektivität
  \proofstep{1,2}{\exists z\in A\times B\; \pi_2(z)=y}{%
    \rEI{\rAI{3,4}}}
\end{tabproof}

\FormulaThmDeltaK[Zweite Projektion als surjektive Funktion]{%
  a_0\in A \vdash \pi_2\colon A\times B\sur B%
}{SecondProjectionSurjective}{%
  \DeltaRow{Mengen}{A\dsep B\dsep a_0\dsep y\dsep z}%
}
\begin{tabproofsplit}
\proofpart{Funktionstyp}
\proofstep{1}{a_0\in A}{\rA}
  \proofstep{}{\pi_2\colon A\times B\to B}{%
    \FormulaRefAuto{\pi_2\colon A\times B\to B}}

\closeproofpart

\proofpart{Surjektivität}
  \setcounter{proofstepnr}{2}% Fortlaufende Schritte dieses Beweises.
\proofstep{3}{y\in B}{\rA}
  \proofstep{1,3}{\exists z\in A\times B\,\pi_2(z)=y}{%
    \FormulaRefAuto{%
      y\in B\dsep a_0\in A\vdash
      \exists z\in A\times B\,\pi_2(z)=y%
    }{3,1}}
  \proofstep{1}{\forall y\in B\,\exists z\in A\times B\,\pi_2(z)=y}{%
    \rUI{\rRI{3,4}}}

\closeproofpart

\proofpart{Zusammenführung}
  \setcounter{proofstepnr}{5}% Fortlaufende Schritte dieses Beweises.
\proofstep{1}{\pi_2\colon A\times B\sur B}{%
    \FormulaRefAuto{Surjektive Funktion}{2,5}}

\closeproofpart
\end{tabproofsplit}


% Nach den Projektionen und den allgemeinen Termbildregeln einbinden.
\subsection{Termbildung über zwei Variablen}

\noindent\begin{minipage}{\linewidth}
\FormulaDefDeltaK[Termbildung über zwei gebundenen Variablen]{%
 \{t(x,y)\mid x,y\in D\}\coloneqq
 \{t(\pi_1(p),\pi_2(p))\mid p\in D\times D\}
}{TreeCodeBinaryTermImageDef}{%
 \DeltaRow{Mengen}{D\dsep x\dsep y\dsep p}
 \DeltaRow{Mengenterm}{t(x,y)}
 \DeltaRow{Projektionen auf \(D\times D\)}{\pi_1\dsep\pi_2}}
\end{minipage}\par

\noindent\begin{minipage}{\linewidth}
\FormulaThmDeltaK[Elementkriterium einer zweistelligen Termbildung]{%
 z\in\{t(x,y)\mid x,y\in D\}
 \eqvdash\exists x,y\in D\;z=t(x,y)
}{TreeCodeBinaryTermImageMembership}{%
 \DeltaRow{Mengen}{D\dsep x\dsep y\dsep p\dsep z}
 \DeltaRow{Mengenterm}{t(x,y)}}
\end{minipage}\par
\begin{tabproofwide}
  \proofstepwidestar[1]{z\in\{t(x,y)\mid x,y\in D\}}{%
    \rA}
  \proofstepwidestar[1]{\exists p\in D\times D\;z=t(\pi_1(p),\pi_2(p))}{%
    \FormulaRefAuto{B27TermImageMembership}[thm]{\rIE{\FormulaRefAuto{TreeCodeBinaryTermImageDef}[def],1}}}
  \proofstepwidestar[3]{p\in D\times D\land z=t(\pi_1(p),\pi_2(p))}{%
    \rA}
  \proofstepwidestar[3]{\pi_1(p)\in D}{%
    \FormulaRefAuto{z\in A\times B\vdash\pi_1(z)\in A}[thm]{\rAEa{3}}}
  \proofstepwidestar[3]{\pi_2(p)\in D}{%
    \FormulaRefAuto{z\in A\times B\vdash\pi_2(z)\in B}[thm]{\rAEa{3}}}
  \proofstepwidestar[3]{\exists y\in D\;z=t(\pi_1(p),y)}{%
    \rEI{\rAI{5,\rAEb{3}}}}
  \proofstepwidestar[3]{\exists x,y\in D\;z=t(x,y)}{%
    \rEI{\rAI{4,6}}}
  \proofstepwidestar[1]{\exists x,y\in D\;z=t(x,y)}{%
    \rEE{2,3,7}}
  \proofstepwidestar[9]{\exists x,y\in D\;z=t(x,y)}{%
    \rA}
  \proofstepwidestar[10]{x\in D\land\exists y\in D\;z=t(x,y)}{%
    \rA}
  \proofstepwidestar[11]{y\in D\land z=t(x,y)}{%
    \rA}
  \proofstepwidestar[10,11]{(x,y)\in D\times D}{%
    \FormulaRefAuto{a\in A\dsep b\in B\vdash(a,b)\in A\times B}[thm]{\rAEa{10},\rAEa{11}}}
  \proofstepwidestar[10,11]{\pi_1(x,y)=x}{%
    \FormulaRefAuto{(x,y)\in A\times B\vdash\pi_1(x,y)=x}[thm]{12}}
  \proofstepwidestar[10,11]{\pi_2(x,y)=y}{%
    \FormulaRefAuto{(x,y)\in A\times B\vdash\pi_2(x,y)=y}[thm]{12}}
  \proofstepwidestar[10,11]{z=t(\pi_1(x,y),\pi_2(x,y))}{%
    \rIE{\FormulaRefAuto{a=b\vdash b=a}[thm]{13},\rIE{\FormulaRefAuto{a=b\vdash b=a}[thm]{14},\rAEb{11}}}}
  \proofstepwidestar[10,11]{\exists p\in D\times D\;z=t(\pi_1(p),\pi_2(p))}{%
    \rEI{\rAI{12,15}}}
  \proofstepwidestar[10,11]{z\in\{t(\pi_1(p),\pi_2(p))\mid p\in D\times D\}}{%
    \FormulaRefAuto{B27TermImageMembership}[thm]{16}}
  \proofstepwidestar[10,11]{z\in\{t(x,y)\mid x,y\in D\}}{%
    \rIE{\FormulaRefAuto{a=b\vdash b=a}[thm]{\FormulaRefAuto{TreeCodeBinaryTermImageDef}[def]},17}}
  \proofstepwidestar[10]{z\in\{t(x,y)\mid x,y\in D\}}{%
    \rEE{\rAEb{10},11,18}}
  \proofstepwidestar[9]{z\in\{t(x,y)\mid x,y\in D\}}{%
    \rEE{9,10,19}}
  \proofstepwidestar[]{z\in\{t(x,y)\mid x,y\in D\}\leftrightarrow\exists x,y\in D\;z=t(x,y)}{%
    \rLRI{\rRI{1,8},\rRI{9,20}}}
\end{tabproofwide}


\subsection{Binäre Operatoren als Funktionen}

Ein binärer Operator auf einer Menge bestimmt eine Funktion auf dem
kartesischen Quadrat dieser Menge. Die Projektionen geben die beiden
Argumente des Operators an.

\FormulaDefDeltaK[Funktionsgraph eines binären Operators]{%
  \mathsf{BinOp}(A,\star)\vdash
  \begin{aligned}[t]
    &\text{(i)}\;\BinaryOpFunction{A}{\star}\colon A\times A\to A\\[-2pt]
    &\text{(ii)}\;\forall x,y\in A\,
      \bigl(\BinaryOpFunction{A}{\star}(x,y)\coloneqq x\star y\bigr)
  \end{aligned}%
}{BinaryOperationFunctionDef}{%
  \DeltaRow{Mengen}{A\dsep p\dsep x\dsep y}
  \DeltaRow{Funktionen}{\star\dsep G\dsep\BinaryOpFunction{A}{\star}}
  \DeltaRow{\textbf{Neues Symbol}}{}
  \DeltaPrem{Operatorfunktion}{\BinaryOpFunction{A}{\star}}
}
\begin{tabproofsplitwide}
  \proofpartwide{Wohldefiniertheit und Funktionstypisierung}
  \proofstepwidestar[1]{\mathsf{BinOp}(A,\star)}{\rA}
  \proofstepwidestar[2]{p\in A\times A}{\rA}
  \proofstepwidestar[2]{\pi_1(p)\in A}{%
    \FormulaRefAuto{z\in A\times B\vdash\pi_1(z)\in A}[thm]{2}}
  \proofstepwidestar[2]{\pi_2(p)\in A}{%
    \FormulaRefAuto{z\in A\times B\vdash\pi_2(z)\in B}[thm]{2}}
  \proofstepwidestar[1,2]{\pi_1(p)\star\pi_2(p)\in A}{%
    \FormulaRefAuto{BinaryOperationClosureAxiom}[ax]{1,3,4}}
  \proofstepwidestar[1]{%
    \forall p\in A\times A\,
    \pi_1(p)\star\pi_2(p)\in A}{\rUI{\rRI{2,5}}}
  \proofstepwidestar[1]{%
    \begin{aligned}[t]
      &\exists!G\colon A\times A\to A\,\\[-2pt]
      &\quad\forall p\in A\times A\,
        G(p)=\pi_1(p)\star\pi_2(p)
    \end{aligned}}{%
    \FormulaRefAuto{TypedFunctionDefinitionPrinciple}[thm]{6}}
  \proofstepwidestar{%
    \begin{aligned}[t]
      &\BinaryOpFunction{A}{\star}\coloneqq
        \iota G\colon A\times A\to A\,\\[-2pt]
      &\quad\forall p\in A\times A\,
        G(p)=\pi_1(p)\star\pi_2(p)
    \end{aligned}}{\FormulaRefAuto{IotaDefinitionDef}[def]{7}}
  \proofstepwidestar[1]{%
    \begin{aligned}[t]
      &\BinaryOpFunction{A}{\star}\colon A\times A\to A\land{}\\[-2pt]
      &\forall p\in A\times A\,
        \BinaryOpFunction{A}{\star}(p)=\pi_1(p)\star\pi_2(p)
    \end{aligned}}{%
    \FormulaRefAuto{IotaDefinitionDef}[def]{7,8}}
  \proofstepwidestar[1]{%
    \BinaryOpFunction{A}{\star}\colon A\times A\to A}{\rAEa{9}}
  \proofstepwidestar[1]{%
    \forall p\in A\times A\,
    \BinaryOpFunction{A}{\star}(p)
      =\pi_1(p)\star\pi_2(p)}{\rAEb{9}}
  \closeproofpartwide

  \proofpartwide{Auswertung an geordneten Paaren}
  \setcounter{proofstepnr}{11}
  \proofstepwidestar[12]{x\in A}{\rA}
  \proofstepwidestar[13]{y\in A}{\rA}
  \proofstepwidestar[12,13]{x\in A\land y\in A}{\rAI{12,13}}
  \proofstepwidestar[12,13]{(x,y)\in A\times A}{%
    \FormulaRefAuto{(a,b)\in A\times B\eqvdash
      a\in A\land b\in B}[thm]{14}}
  \proofstepwidestar[12,13]{\pi_1(x,y)=x}{%
    \FormulaRefAuto{(x,y)\in A\times B\vdash
      \pi_1(x,y)=x}[thm]{15}}
  \proofstepwidestar[12,13]{\pi_2(x,y)=y}{%
    \FormulaRefAuto{(x,y)\in A\times B\vdash
      \pi_2(x,y)=y}[thm]{15}}
  \proofstepwidestar[1]{%
    (x,y)\in A\times A\rightarrow
    \BinaryOpFunction{A}{\star}(x,y)
      =\pi_1(x,y)\star\pi_2(x,y)}{\rUE{11}}
  \proofstepwidestar[1,12,13]{%
    \BinaryOpFunction{A}{\star}(x,y)
      =\pi_1(x,y)\star\pi_2(x,y)}{\rRE{18,15}}
  \proofstepwidestar[12,13]{%
    \pi_1(x,y)\star\pi_2(x,y)=x\star y}{\rIE{16,17}}
  \proofstepwidestar[1,12,13]{%
    \BinaryOpFunction{A}{\star}(x,y)=x\star y}{\rChain{19,20}}
  \proofstepwidestar[1]{%
    \forall x,y\in A\,
    \BinaryOpFunction{A}{\star}(x,y)=x\star y}{%
    \rUI{\rRI{12,\rUI{\rRI{13,21}}}}}
  \closeproofpartwide
\end{tabproofsplitwide}

\subsection{Filter einer einzelnen Faser}

Für beliebige Mengen \(U,V,K\) behält der folgende Filter genau die
geordneten Paare aus \(K\), die in \(U\times V\) liegen und deren erste
Komponente von \(t\) verschieden oder deren zweite Komponente gleich
\(c\) ist. Dabei müssen \(t\) und \(c\) nicht in den Trägermengen liegen.

\FormulaDefDeltaK[Filter einer einzelnen Faser]{%
  \begin{aligned}[t]
    &\Phi^{U,V}_{t,c}(K)\coloneqq
      \{p\in U\times V\mid
        (\pi_1(p),\pi_2(p))\in K\land{}\\[-2pt]
    &\qquad(\pi_1(p)\ne t\lor\pi_2(p)=c)\}
  \end{aligned}
}{WordRecursionFiberFilterDef}{%
  \DeltaRow{Mengen}{U\dsep V\dsep K\dsep t\dsep c\dsep p}
  \DeltaRow{Projektionen auf \(U\times V\)}{\pi_1\dsep\pi_2}
  \DeltaRow{Neues Symbol}{\Phi^{U,V}_{t,c}(K)}}

Die definierte Menge ist eine Aussonderung aus \(U\times V\).
Ihre Existenz und Eindeutigkeit liegen in
\FormulaRefAuto{\{x\in A\,\mid\,P(x)\}\coloneq
  \iota B\bigl(\forall x\,(x\in B\leftrightarrow
    (x\in A\land P(x)))\bigr)}[def].

\noindent\begin{minipage}{\linewidth}
\FormulaThmDeltaK[Ein gefilterter Graph ist eine Teilmenge]{%
  \Phi^{U,V}_{t,c}(K)\subseteq K
}{WordRecursionFilterSubset}{%
  \DeltaRow{Mengen}{U\dsep V\dsep K\dsep t\dsep c\dsep p}
  \DeltaRow{Projektionen auf \(U\times V\)}{\pi_1\dsep\pi_2}}
\end{minipage}\par
\begin{tabproofwide}
  \proofstepwidestar[1]{p\in\Phi^{U,V}_{t,c}(K)}{\rA}
  \proofstepwidestar[]{\begin{aligned}[t]
    &\Phi^{U,V}_{t,c}(K)=\{p\in U\times V\mid
      (\pi_1(p),\pi_2(p))\in K\land{}\\[-2pt]
    &\qquad(\pi_1(p)\ne t\lor\pi_2(p)=c)\}
  \end{aligned}}{\FormulaRefAuto{WordRecursionFiberFilterDef}[def]}
  \proofstepwidestar[1]{\begin{aligned}[t]
    &p\in U\times V\land
      ((\pi_1(p),\pi_2(p))\in K\land{}\\[-2pt]
    &\qquad(\pi_1(p)\ne t\lor\pi_2(p)=c))
  \end{aligned}}{%
    \FormulaRefAuto{x\in\{u\in A\mid P(u)\}\eqvdash
      x\in A\land P(x)}[thm]{\rIE{2,1}}}
  \proofstepwidestar[1]{p=(\pi_1(p),\pi_2(p))}{%
    \FormulaRefAuto{z\in A\times B\vdash
      z=\bigl(\pi_1(z),\pi_2(z)\bigr)}[thm]{\rAEa{3}}}
  \proofstepwidestar[1]{p\in K}{%
    \rIE{\FormulaRefAuto{a=b\vdash b=a}[thm]{4},\rAEa{\rAEb{3}}}}
  \proofstepwidestar[]{\Phi^{U,V}_{t,c}(K)\subseteq K}{%
    \FormulaRefAuto{SubsetIntroductionRule}[thm]{[1]\vdots 5}}
\end{tabproofwide}

\noindent\begin{minipage}{\linewidth}
\FormulaThmDeltaK[Elementkriterium des Faserfilters]{%
  \begin{aligned}[t]
    &K\subseteq U\times V\vdash{}\\[-2pt]
    &(u,x)\in\Phi^{U,V}_{t,c}(K)\leftrightarrow
      (u,x)\in K\land(u\ne t\lor x=c)
  \end{aligned}
}{WordRecursionFilterMembership}{%
  \DeltaRow{Mengen}{U\dsep V\dsep K\dsep t\dsep c\dsep u\dsep x}
  \DeltaRow{Projektionen auf \(U\times V\)}{\pi_1\dsep\pi_2}}
\end{minipage}\par
\begin{tabproofwide}
  \proofstepwidestar[1]{K\subseteq U\times V}{\rA}
  \proofstepwidestar[]{\begin{aligned}[t]
    &\Phi^{U,V}_{t,c}(K)=\{p\in U\times V\mid
      (\pi_1(p),\pi_2(p))\in K\land{}\\[-2pt]
    &\qquad(\pi_1(p)\ne t\lor\pi_2(p)=c)\}
  \end{aligned}}{\FormulaRefAuto{WordRecursionFiberFilterDef}[def]}
  \proofstepwidestar[]{\begin{aligned}[t]
    &(u,x)\in\Phi^{U,V}_{t,c}(K)\leftrightarrow{}\\[-2pt]
    &u\in U\land x\in V\land
      ((u,x)\in K\land(u\ne t\lor x=c))
  \end{aligned}}{%
    \rIE{\FormulaRefAuto{a=b\vdash b=a}[thm]{2},
      \FormulaRefAuto{ProductComprehensionMembership}[thm]}}
  \proofstepwidestar[4]{(u,x)\in K\land(u\ne t\lor x=c)}{\rA}
  \proofstepwidestar[1,4]{(u,x)\in U\times V}{%
    \FormulaRefAuto{A\subseteq B\dsep x\in A\vdash x\in B}[thm]{1,\rAEa{4}}}
  \proofstepwidestar[1,4]{u\in U\land x\in V}{%
    \FormulaRefAuto{(a,b)\in A\times B\eqvdash
      a\in A\land b\in B}[thm]{5}}
  \proofstepwidestar[1]{\begin{aligned}[t]
    &((u,x)\in K\land(u\ne t\lor x=c))\\[-2pt]
    &\qquad\rightarrow(u\in U\land x\in V)
  \end{aligned}}{%
    \rRI{4,6}}
  \proofstepwidestar[1]{\begin{aligned}[t]
    &(u,x)\in\Phi^{U,V}_{t,c}(K)\leftrightarrow{}\\[-2pt]
    &(u,x)\in K\land(u\ne t\lor x=c)
  \end{aligned}}{%
    \rLRS{\FormulaRefAuto{P\rightarrow Q\vdash
      (Q\land P)\leftrightarrow P}[thm]{7},
      \rLRS{\FormulaRefAuto{P\land(Q\land R)\eqvdash
        (P\land Q)\land R}[thm],3}}}
\end{tabproofwide}

\noindent\begin{minipage}{\linewidth}
\FormulaThmDeltaK[Der Filter erzwingt seinen vorgeschriebenen Faserwert]{%
  K\subseteq U\times V\dsep
    (t,z)\in\Phi^{U,V}_{t,c}(K)\vdash z=c
}{WordRecursionFilterOwnFiber}{%
  \DeltaRow{Mengen}{U\dsep V\dsep K\dsep t\dsep c\dsep z}}
\end{minipage}\par
\begin{tabproofwide}
  \proofstepwidestar[1]{K\subseteq U\times V}{\rA}
  \proofstepwidestar[2]{(t,z)\in\Phi^{U,V}_{t,c}(K)}{\rA}
  \proofstepwidestar[1,2]{(t,z)\in K\land(t\ne t\lor z=c)}{%
    \FormulaRefAuto{WordRecursionFilterMembership}[thm]{1,2}}
  \proofstepwidestar[1,2]{t=t\rightarrow z=c}{%
    \FormulaRefAuto{P\rightarrow Q\eqvdash\neg P\lor Q}[thm]{\rAEb{3}}}
  \proofstepwidestar[]{t=t}{\rII}
  \proofstepwidestar[1,2]{z=c}{\rRE{5,4}}
\end{tabproofwide}


\Needspace{16\baselineskip}
\section{Fasern surjektiver Funktionen}

\FormulaThmDeltaR{y\in B \vdash \Fib_F(y)\neq\varnothing}
{F\colon A\sur B \dsep y\in B \vdash \Fib_F(y)\neq\varnothing}
{
  \DeltaRow{Mengen}{A\dsep B\dsep x\dsep y}
  \DeltaRow{Surjektive Funktion}{F\colon A\sur B}
}
\begin{tabproof}
  \proofstep{1}{y\in B}{\rA}
  \proofstep{1}{\exists x\in A\,F(x)=y}{\FormulaRefAuto{Surjektivität}{1}}
  \proofstep{3}{x\in A \land F(x)=y}{\rA}
  \proofstep{3}{x\in A}{\rAEa{3}}
  \proofstep{3}{F(x)=y}{\rAEb{3}}
  \proofstep{1,3}{x\in \Fib_F(y)}{%
    \FormulaRefAuto{y\in B\dsep x\in A\dsep F(x)=y \vdash x\in\Fib_F(y)}{1,4,5}}
  \proofstep{1,3}{\Fib_F(y)\neq\varnothing}{\FormulaRefAuto{a\in A\vdash A\neq\varnothing}{6}}
  \proofstep{1}{\Fib_F(y)\neq\varnothing}{\rEE{2,3,7}}
\end{tabproof}

\FormulaThmDelta
{X\in \Fib_F[B] \vdash X\neq\varnothing}
{
  \DeltaRow{Mengen}{A\dsep B\dsep X\dsep y}
  \DeltaRow{Surjektive Funktion}{F\colon A\sur B}
}
\begin{tabproof}
  \proofstep{1}{X\in \Fib_F[B]}{\rA}
  \proofstep{}{\Fib_F\colon B\to \powerset(A)}{%
    \FormulaRefAuto{\Fib_F\colon B\to \powerset(A)}}

  \proofstep{1}{\exists y\in B\,X=\Fib_F(y)}{%
    \FormulaRefAuto{F\colon A\to B\dsep y\in F[A] \vdash \exists x\in A\,y=F(x)}{2,1}}

  \proofstep{4}{y\in B \land X=\Fib_F(y)}{\rA}
  \proofstep{4}{y\in B}{\rAEa{4}}
  \proofstep{4}{X=\Fib_F(y)}{\rAEb{4}}

  \proofstep{5}{\Fib_F(y)\neq\varnothing}{%
    \FormulaRefAuto{F\colon A\sur B \dsep y\in B \vdash \Fib_F(y)\neq\varnothing}{5}}

  \proofstep{4}{X\neq\varnothing}{\rIE{6,7}}
  \proofstep{1}{X\neq\varnothing}{\rEE{3,4,8}}
\end{tabproof}

\FormulaThmDelta
{\DisjFam(\Fib_F[B])}
{
  \DeltaRow{Mengen}{A\dsep B\dsep X\dsep Y}
  \DeltaRow{Surjektive Funktion}{F\colon A\sur B}
}
\begin{tabproofsplit}
\proofpart{Nichtleerheit der Familienglieder}
\proofstep{1}{X\in \Fib_F[B]}{\rA}
  \proofstep{1}{X\neq\varnothing}{%
    \FormulaRefAuto{X\in \Fib_F[B] \vdash X\neq\varnothing}{1}}

\closeproofpart

\proofpart{Paarweise Disjunktheit}
  \setcounter{proofstepnr}{2}% Fortlaufende Schritte dieses Beweises.
\proofstep{3}{Y\in \Fib_F[B]}{\rA}
  \proofstep{4}{X\neq Y}{\rA}
  \proofstep{1,3,4}{X\cap Y=\varnothing}{%
    \FormulaRefAuto{X\in \Fib_F[B] \dsep Y\in \Fib_F[B] \dsep X\neq Y \vdash X\cap Y=\varnothing}{1,3,4}}

\closeproofpart

\proofpart{Zusammenführung}
  \setcounter{proofstepnr}{5}% Fortlaufende Schritte dieses Beweises.
\proofstep{}{\DisjFam(\Fib_F[B])}{%
    \FormulaRefAuto{Disjunkte Familie}{2,5}}

\closeproofpart
\end{tabproofsplit}


\section{Einschränkungen und Korestriktionen}

\subsection{Einschränkung auf das Bild}

\FormulaThmDelta[Surjektivität der Einschränkung auf das Bild]{%
  C\subseteq A \dsep y\in F[C] \vdash \exists x\in C\, F\restriction_C(x)=y
}{
  \DeltaRow{Mengen}{A\dsep B\dsep C\dsep x\dsep y}
  \DeltaPrem{Funktionen}{F\colon A\to B}
}
\begin{tabproof}
  \proofstep{1}{C\subseteq A}{\rA}
  \proofstep{2}{y\in F[C]}{\rA}
  \proofstep{1,2}{\exists x\in C\, y=F(x)}{%
    \FormulaRefAuto{C\subseteq A\dsep F\colon A\to B\dsep y\in F[C]\vdash \exists x\in C\,y=F(x)}{1,2}}

  \proofstep{4}{x\in C\land y=F(x)}{\rA}
  \proofstep{4}{x\in C}{\rAEa{4}}
  \proofstep{4}{y=F(x)}{\rAEb{4}}
  \proofstep{1,4}{F\restriction_C(x)=F(x)}{%
    \FormulaRefAuto{C\subseteq A \dsep x\in C \vdash F\restriction_C(x)=F(x)}{1,5}}
  \proofstep{4}{F(x)=y}{\FormulaRefAuto{a=b\vdash b=a}{6}}
  \proofstep{1,4}{F\restriction_C(x)=y}{%
    \FormulaRefAuto{a=b,\, b=c \vdash a=c}{7,8}}
  \proofstep{1,4}{\exists x\in C\,F\restriction_C(x)=y}{\rEI{\rAI{5,9}}}

  \proofstep{1,2}{\exists x\in C\,F\restriction_C(x)=y}{\rEE{3,4,10}}
\end{tabproof}

\FormulaThmDelta[Einschränkung als surjektive Funktion auf das Bild]{%
C\subseteq A \vdash F\restriction_C\colon C\twoheadrightarrow F[C]
}{
  \DeltaRow{Mengen}{A\dsep B\dsep C}
  \DeltaPrem{Funktionen}{F\colon A\to B}
}
\begin{tabproof}
  \proofstep{1}{C\subseteq A}{\rA}
  \proofstep{1}{F\restriction_C\colon C\to B}{%
    \FormulaRefAuto{C\subseteq A \vdash F\restriction_C\colon C\to B}{1}}
  \proofstep{1}{\forall y\in F[C]\;\exists x\in C\,F\restriction_C(x)=y}{%
    \FormulaRefAuto{C\subseteq A \dsep y\in F[C] \vdash \exists x\in C\, F\restriction_C(x)=y}{1}}
  \proofstep{1}{F\restriction_C\colon C\twoheadrightarrow F[C]}{%
    \FormulaRefAuto{Surjektive Funktion}{2,3}}
\end{tabproof}

\subsection{Einschränkung auf Urbilder}

\FormulaThmDelta[Zeugen im Urbild einer Teilmenge]{%
  F\colon A\sur B \dsep T\subseteq B \dsep y\in T \vdash
  \exists x\in F^{-1}[T]\,F\restriction_{F^{-1}[T]}(x)=y
}{
  \DeltaRow{Mengen}{A\dsep B\dsep T\dsep x\dsep y}
  \DeltaRow{Funktionen}{F}
}
\begin{tabproof}
  \proofstep{1}{F\colon A\sur B}{\rA}
  \proofstep{2}{T\subseteq B}{\rA}
  \proofstep{3}{y\in T}{\rA}
  \proofstep{1}{F\colon A\to B}{\FormulaRefAuto{Surjektive Funktion}{1}}
  \proofstep{2,4}{F^{-1}[T]\subseteq A}{%
    \FormulaRefAuto{F\colon A\to B\dsep C\subseteq B \vdash F^{-1}[C]\subseteq A}{4,2}}
  \proofstep{2,3}{y\in B}{%
    \FormulaRefAuto{A\subseteq B, x\in A \vdash x\in B}{2,3}}
  \proofstep{1,2,3}{\exists x\in A\,F(x)=y}{%
    \FormulaRefAuto{y\in B\vdash \exists x\in A\; F(x)=y}{6}}
  \proofstep{8}{x\in A \land F(x)=y}{\rA}
  \proofstep{8}{x\in A}{\rAEa{8}}
  \proofstep{8}{F(x)=y}{\rAEb{8}}
  \proofstep{3,8}{F(x)\in T}{\rIE{10,3}}
  \proofstep{4,2}{x\in F^{-1}[T]\leftrightarrow (x\in A \land F(x)\in T)}{%
    \FormulaRefAuto{F\colon A\to B\dsep C\subseteq B
    \vdash \bigl(x\in F^{-1}[C] \leftrightarrow
    (x\in A \land F(x)\in C)\bigr)}{4,2}}
  \proofstep{8}{x\in A \land F(x)\in T}{\rAI{9,11}}
  \proofstep{2,8}{x\in F^{-1}[T]}{%
    \FormulaRefAuto{P \leftrightarrow Q, Q \vdash P}{12,13}}
  \proofstep{5,8}{F\restriction_{F^{-1}[T]}(x)=F(x)}{%
    \FormulaRefAuto{C\subseteq A \dsep x\in C \vdash F\restriction_C(x)=F(x)}{5,14}}
  \proofstep{5,8}{F\restriction_{F^{-1}[T]}(x)=y}{%
    \FormulaRefAuto{a = b,\, b = c \vdash a = c}{15,10}}
  \proofstep{5,8}{\exists x\in F^{-1}[T]\,F\restriction_{F^{-1}[T]}(x)=y}{%
    \rEI{\rAI{14,16}}}
  \proofstep{1,2,3}{\exists x\in F^{-1}[T]\,F\restriction_{F^{-1}[T]}(x)=y}{%
    \rEE{7,8,17}}
\end{tabproof}

\FormulaThmDelta[Einschränkung auf ein Urbild als surjektive Funktion]{%
  F\colon A\sur B \dsep T\subseteq B \vdash F\restriction_{F^{-1}[T]}\colon F^{-1}[T]\sur T
}{
  \DeltaRow{Mengen}{A\dsep B\dsep T\dsep y}
  \DeltaRow{Funktionen}{F}
}
\begin{tabproof}
  \proofstep{1}{F\colon A\sur B}{\rA}
  \proofstep{2}{T\subseteq B}{\rA}
  \proofstep{1}{F\colon A\to B}{\FormulaRefAuto{Surjektive Funktion}{1}}
  \proofstep{2,3}{F\restriction_{F^{-1}[T]}\colon F^{-1}[T]\to T}{%
    \FormulaRefAuto{F\colon A\to B\dsep T\subseteq B \vdash F\restriction_{F^{-1}[T]}\colon F^{-1}[T]\to T}{3,2}}
  \proofstep{1,2}{\forall y\in T\,\exists x\in F^{-1}[T]\,F\restriction_{F^{-1}[T]}(x)=y}{%
    \FormulaRefAuto{F\colon A\sur B \dsep T\subseteq B \dsep y\in T \vdash \exists x\in F^{-1}[T]\,F\restriction_{F^{-1}[T]}(x)=y}{1,2}}
  \proofstep{1,2}{F\restriction_{F^{-1}[T]}\colon F^{-1}[T]\sur T}{%
    \FormulaRefAuto{Surjektive Funktion}{4,5}}
\end{tabproof}


\subsection{Eingeschränkte Projektionen}

\paragraph{Projektionen}

\FormulaThmDelta[Surjektivität]
{x\in A\vdash \exists y\in R(\pi_1\restriction_R(y)=x)}
{
  \DeltaRow{Mengen}{A\dsep B\dsep x}
  \DeltaPrem{Totale Relationen}{\TotRel{R,A,B}}
}
\begin{tabproofwide}
  \proofstepwidestar[1]{x\in A}{\rA}
  \proofstepwidestar[]{R\subseteq A\times B}{\FormulaRefAuto{F \subseteq A \times B}}
  \proofstepwidestar[1]{\exists x' \, (x,x')\in R}{\FormulaRefAuto{x\in A \vdash \exists y\,(x,y)\in F}{1}}
  \proofstepwidestar[4]{(x,x')\in R}{\rA}
  \proofstepwidestar[4]{(x,x')\in A\times B}{\FormulaRefAuto{A\subseteq B, x\in A\vdash x\in B}{2,4}}
  \proofstepwide[4]{\pi_1\restriction_R(x,x')}{=}{\pi_1(x,x')}{\FormulaRefAuto{C\subseteq A \dsep x\in C \vdash F\restriction_C(x)=F(x)}{2,4}}
  \proofstepwide[4]{}{=}{x}{\FormulaRefAuto{(x,y)\in A\times B \vdash \pi_1(x,y)=x}{5}}
  \proofstepwide[4]{\pi_1\restriction_R(x,x')}{=}{x}{\rChain{6,7}}
  \proofstepwidestar[4]{\exists y\in R \, \pi_1\restriction_R(y)=x}{\rEI{\rAI{4,8}}}
  \proofstepwidestar[1]{\exists y\in R \, \pi_1\restriction_R(y)=x}{\rEE{3,4,9}}
\end{tabproofwide}

\FormulaThmDelta[Surjektive Funktion]{%
\pi_1\restriction_R\colon R\sur A
}{
  \DeltaRow{Mengen}{A\dsep B}
  \DeltaPrem{Totale Relationen}{\TotRel{R,A,B}}
}
\begin{tabproof}
    \proofstep{}{R\subseteq A\times B}{%
      \FormulaRefAuto{R\subseteq A\times B}}
    \proofstep{}{\pi_1\restriction_R\colon R\to A}{\FormulaRefAuto{C\subseteq A \vdash F\restriction_C\colon C\to B}{1}}
    \proofstep{}{\forall x\in A \exists y\in R\,(\pi_1\restriction_R(y)=x)}{\FormulaRefAuto{x\in A\vdash \exists y\in R(\pi_1\restriction_R(y)=x)}}
    \proofstep{}{\pi_1\restriction_R\colon R\sur A}{\FormulaRefAuto{Surjektive Funktion}{2,3}}
\end{tabproof}


\subsection{Korestriktion auf das Bild}

Die Korestriktion einer Funktion auf ihr Bild wurde bereits im grundlegenden
Aufbau dieses Bandes als surjektive Standardkonstruktion nachgewiesen. Sie wird
im Folgenden ohne erneute Herleitung verwendet.

\section{Kompositionen}

\subsection{Erhalt der Surjektivität}

\FormulaThmDelta[Erhalt der Surjektivität]{%
z\in C \vdash \exists x\in A\, (G\circ F)(x)=z
}{
  \DeltaRow{Mengen}{x \dsep y \dsep z \dsep A \dsep B \dsep C}
  \DeltaPrem{Surjektive Funktionen}{F\colon A\sur B \dsep G\colon B\sur C}
}
\begin{tabproof}
  \proofstep{1}{z\in C}{\rA}
  \proofstep{1}{\exists y\in B\, G(y)=z}{\FormulaRefAuto{Surjektivität}{1}} % (Schema: Surj. von G)
  \proofstep{3}{y\in B \land G(y)=z}{\rA}
  \proofstep{3}{y\in B}{\rAEa{3}}
  \proofstep{3}{G(y)=z}{\rAEb{3}}
  \proofstep{3}{\exists x\in A\, F(x)=y}{\FormulaRefAuto{y\in B \vdash \exists x\in A\, F(x)=y}{4}}
  \proofstep{7}{x\in A \land F(x)=y}{\rA}
  \proofstep{7}{x\in A}{\rAEa{7}}
  \proofstep{7}{F(x)=y}{\rAEb{7}}
  \proofstep{3,7}{G(F(x))=z}{\rIE{9,5}}
  \proofstep{3,7}{(G\circ F)(x)=G(F(x))}{\FormulaRefAuto{x\in A \vdash (G\circ F)(x)=G(F(x))}{8}}
  \proofstep{3,7}{(G\circ F)(x)=z}{\rIE{11,10}}
  \proofstep{3,7}{\exists x\in A\,(G\circ F)(x)=z}{\rEI{\rAI{8,12}}}
  \proofstep{3}{\exists x\in A\,(G\circ F)(x)=z}{\rEE{6,7,13}}
  \proofstep{1}{\exists x\in A\,(G\circ F)(x)=z}{\rEE{2,3,14}}
\end{tabproof}

\FormulaThmDelta[Die Komposition als surjektive Funktion]{%
G\circ F\colon A\sur C
}{
  \DeltaRow{Mengen}{A\dsep B\dsep C}
  \DeltaPrem{Surjektive Funktionen}{F\colon A\sur B \dsep G\colon B\sur C}
}
\begin{tabproof}
    \proofstep{}{G\circ F\colon A\to C}{\FormulaRefAuto{G\circ F\colon A\to C}}
    \proofstep{}{\forall z\in C\exists x\in A\, (G\circ F)(x)=z}{\FormulaRefAuto{z\in C \vdash \exists x\in A\, (G\circ F)(x)=z}}
    \proofstep{}{G\circ F\colon A\sur C}{\FormulaRefAuto{Surjektive Funktion}{1,2}}
\end{tabproof}


\section{Abstieg entlang von Surjektionen}

% Logical access lemmas for representative descent; no new inference rules.
\subsection{Quantoren über Repräsentanten}

Die folgenden Schemata verbinden beschränkte Repräsentantenquantoren mit
ausgeschriebenen Trägerbedingungen. Die Umformungen werden bewiesen;
insbesondere wird dabei kein Quantor stillschweigend vertauscht.
Die Trägermengen enthalten die gebundenen Variablen nicht frei.
Zusätzliche Parameter der Prädikate bleiben fest; Zeugenvariablen werden
frisch gewählt. In den jeweiligen Beweisen stehen \(\mathcal E\) und
\(\mathcal G\) für die linke und rechte Seite der angegebenen Äquivalenz.

\FormulaThmDeltaK[Beschränkte Existenzquantoren über ein Paar]{%
  \begin{aligned}[t]&\exists a\in A\,\exists b\in B\,(H(a,b))\\[-2pt]&\quad\leftrightarrow \exists a\exists b\,\bigl((a\in A\land b\in B)\land H(a,b)\bigr)\end{aligned}
}{BoundedPairExistsGuard}{%
  \DeltaRow{Mengen}{A\dsep B}
  \DeltaRow{Prädikate}{H\dsep L\dsep T}
}
\begin{tabproofwide}
  \proofstepwidestar[1]{\mathcal E}{%
    \rA}
  \proofstepwidestar[2]{a\in A}{%
    \rA}
  \proofstepwidestar[3]{b\in B}{%
    \rA}
  \proofstepwidestar[4]{H(a,b)}{%
    \rA}
  \proofstepwidestar[2,3]{a\in A\land b\in B}{%
    \rAI{2,3}}
  \proofstepwidestar[2,3,4]{(a\in A\land b\in B)\land H(a,b)}{%
    \rAI{5,4}}
  \proofstepwidestar[2,3,4]{\mathcal G}{%
    \rEIStar{6}}
  \proofstepwidestar[1]{\mathcal G}{%
    \rEEStar{1,2,3,4,7}}
  \proofstepwidestar[]{\mathcal E\rightarrow\mathcal G}{%
    \rRI{1,8}}
  \proofstepwidestar[10]{\mathcal G}{%
    \rA}
  \proofstepwidestar[11]{(a\in A\land b\in B)\land H(a,b)}{%
    \rA}
  \proofstepwidestar[11]{a\in A\land b\in B}{%
    \rAEa{11}}
  \proofstepwidestar[11]{a\in A}{%
    \rAEa{12}}
  \proofstepwidestar[11]{b\in B}{%
    \rAEb{12}}
  \proofstepwidestar[11]{H(a,b)}{%
    \rAEb{11}}
  \proofstepwidestar[11]{\mathcal E}{%
    \rEIStar{13,14,15}}
  \proofstepwidestar[10]{\mathcal E}{%
    \rEEStar{10,11,16}}
  \proofstepwidestar[]{\mathcal G\rightarrow\mathcal E}{%
    \rRI{10,17}}
  \proofstepwidestar[]{\mathcal E\leftrightarrow\mathcal G}{%
    \rLRI{9,18}}
\end{tabproofwide}

\FormulaThmDeltaK[Vier Existenzquantoren mit paarweisen Trägerbedingungen]{%
  \begin{aligned}[t]&\exists a\in A\,\exists b\in B\,\exists c\in A\,\exists d\in B\,(H(a,b,c,d))\\[-2pt]&\quad\leftrightarrow \exists a\exists b\exists c\exists d\,\bigl((a\in A\land b\in B)\land (c\in A\land d\in B)\land H(a,b,c,d)\bigr)\end{aligned}
}{BoundedFourExistsGuard}{%
  \DeltaRow{Mengen}{A\dsep B}
  \DeltaRow{Prädikate}{H\dsep L\dsep T}
}
\begin{tabproofwide}
  \proofstepwidestar[1]{\mathcal E}{%
    \rA}
  \proofstepwidestar[2]{a\in A}{%
    \rA}
  \proofstepwidestar[3]{b\in B}{%
    \rA}
  \proofstepwidestar[4]{c\in A}{%
    \rA}
  \proofstepwidestar[5]{d\in B}{%
    \rA}
  \proofstepwidestar[6]{H(a,b,c,d)}{%
    \rA}
  \proofstepwidestar[2,3]{a\in A\land b\in B}{%
    \rAI{2,3}}
  \proofstepwidestar[4,5]{c\in A\land d\in B}{%
    \rAI{4,5}}
  \proofstepwidestar[2,3,4,5,6]{(a\in A\land b\in B)\land (c\in A\land d\in B)\land H(a,b,c,d)}{%
    \rAIStar{7,8,6}}
  \proofstepwidestar[2,3,4,5,6]{\mathcal G}{%
    \rEIStar{9}}
  \proofstepwidestar[1]{\mathcal G}{%
    \rEEStar{1,2,3,4,5,6,10}}
  \proofstepwidestar[]{\mathcal E\rightarrow\mathcal G}{%
    \rRI{1,11}}
  \proofstepwidestar[13]{\mathcal G}{%
    \rA}
  \proofstepwidestar[14]{(a\in A\land b\in B)\land (c\in A\land d\in B)\land H(a,b,c,d)}{%
    \rA}
  \proofstepwidestar[14]{a\in A\land b\in B}{%
    \rAEa{14}}
  \proofstepwidestar[14]{a\in A}{%
    \rAEa{15}}
  \proofstepwidestar[14]{b\in B}{%
    \rAEb{15}}
  \proofstepwidestar[14]{c\in A\land d\in B}{%
    \rAEa{\rAEb{14}}}
  \proofstepwidestar[14]{c\in A}{%
    \rAEa{18}}
  \proofstepwidestar[14]{d\in B}{%
    \rAEb{18}}
  \proofstepwidestar[14]{H(a,b,c,d)}{%
    \rAEb{\rAEb{14}}}
  \proofstepwidestar[14]{\mathcal E}{%
    \rEIStar{16,17,19,20,21}}
  \proofstepwidestar[13]{\mathcal E}{%
    \rEEStar{13,14,22}}
  \proofstepwidestar[]{\mathcal G\rightarrow\mathcal E}{%
    \rRI{13,23}}
  \proofstepwidestar[]{\mathcal E\leftrightarrow\mathcal G}{%
    \rLRI{12,24}}
\end{tabproofwide}

\FormulaThmDeltaK[Vier Existenzquantoren nach Trägern gruppieren]{%
  \begin{aligned}[t]&\exists a\in A\,\exists c\in A\,\exists b\in B\,\exists d\in B\,(H(a,b,c,d))\\[-2pt]&\quad\leftrightarrow \exists a\exists b\exists c\exists d\,\bigl((a\in A\land b\in B)\land (c\in A\land d\in B)\land H(a,b,c,d)\bigr)\end{aligned}
}{BoundedFourExistsGroupedGuard}{%
  \DeltaRow{Mengen}{A\dsep B}
  \DeltaRow{Prädikate}{H\dsep L\dsep T}
}
\begin{tabproofwide}
  \proofstepwidestar[1]{\mathcal E}{%
    \rA}
  \proofstepwidestar[2]{a\in A}{%
    \rA}
  \proofstepwidestar[3]{c\in A}{%
    \rA}
  \proofstepwidestar[4]{b\in B}{%
    \rA}
  \proofstepwidestar[5]{d\in B}{%
    \rA}
  \proofstepwidestar[6]{H(a,b,c,d)}{%
    \rA}
  \proofstepwidestar[2,4]{a\in A\land b\in B}{%
    \rAI{2,4}}
  \proofstepwidestar[3,5]{c\in A\land d\in B}{%
    \rAI{3,5}}
  \proofstepwidestar[2,3,4,5,6]{(a\in A\land b\in B)\land (c\in A\land d\in B)\land H(a,b,c,d)}{%
    \rAIStar{7,8,6}}
  \proofstepwidestar[2,3,4,5,6]{\mathcal G}{%
    \rEIStar{9}}
  \proofstepwidestar[1]{\mathcal G}{%
    \rEEStar{1,2,3,4,5,6,10}}
  \proofstepwidestar[]{\mathcal E\rightarrow\mathcal G}{%
    \rRI{1,11}}
  \proofstepwidestar[13]{\mathcal G}{%
    \rA}
  \proofstepwidestar[14]{(a\in A\land b\in B)\land (c\in A\land d\in B)\land H(a,b,c,d)}{%
    \rA}
  \proofstepwidestar[14]{a\in A\land b\in B}{%
    \rAEa{14}}
  \proofstepwidestar[14]{a\in A}{%
    \rAEa{15}}
  \proofstepwidestar[14]{b\in B}{%
    \rAEb{15}}
  \proofstepwidestar[14]{c\in A\land d\in B}{%
    \rAEa{\rAEb{14}}}
  \proofstepwidestar[14]{c\in A}{%
    \rAEa{18}}
  \proofstepwidestar[14]{d\in B}{%
    \rAEb{18}}
  \proofstepwidestar[14]{H(a,b,c,d)}{%
    \rAEb{\rAEb{14}}}
  \proofstepwidestar[14]{\mathcal E}{%
    \rEIStar{16,19,17,20,21}}
  \proofstepwidestar[13]{\mathcal E}{%
    \rEEStar{13,14,22}}
  \proofstepwidestar[]{\mathcal G\rightarrow\mathcal E}{%
    \rRI{13,23}}
  \proofstepwidestar[]{\mathcal E\leftrightarrow\mathcal G}{%
    \rLRI{12,24}}
\end{tabproofwide}

\FormulaThmDeltaK[Zusatzbedingung einer Paarrepräsentation]{%
  \begin{aligned}[t]&\exists a\in A\,\exists b\in B\,(H(a,b)\land L(a,b)\land T(a,b))\\[-2pt]&\quad\leftrightarrow \exists a\exists b\,\bigl((a\in A\land b\in B\land L(a,b))\land H(a,b)\land T(a,b)\bigr)\end{aligned}
}{BoundedPairExistsGuardWithCondition}{%
  \DeltaRow{Mengen}{A\dsep B}
  \DeltaRow{Prädikate}{H\dsep L\dsep T}
}
\begin{tabproofwide}
  \proofstepwidestar[1]{\mathcal E}{%
    \rA}
  \proofstepwidestar[2]{a\in A}{%
    \rA}
  \proofstepwidestar[3]{b\in B}{%
    \rA}
  \proofstepwidestar[4]{H(a,b)\land L(a,b)\land T(a,b)}{%
    \rA}
  \proofstepwidestar[4]{H(a,b)}{%
    \rAEa{4}}
  \proofstepwidestar[4]{L(a,b)}{%
    \rAEa{\rAEb{4}}}
  \proofstepwidestar[4]{T(a,b)}{%
    \rAEb{\rAEb{4}}}
  \proofstepwidestar[2,3,4]{a\in A\land b\in B\land L(a,b)}{%
    \rAIStar{2,3,6}}
  \proofstepwidestar[2,3,4]{(a\in A\land b\in B\land L(a,b))\land H(a,b)\land T(a,b)}{%
    \rAIStar{8,5,7}}
  \proofstepwidestar[2,3,4]{\mathcal G}{%
    \rEIStar{9}}
  \proofstepwidestar[1]{\mathcal G}{%
    \rEEStar{1,2,3,4,10}}
  \proofstepwidestar[]{\mathcal E\rightarrow\mathcal G}{%
    \rRI{1,11}}
  \proofstepwidestar[13]{\mathcal G}{%
    \rA}
  \proofstepwidestar[14]{(a\in A\land b\in B\land L(a,b))\land H(a,b)\land T(a,b)}{%
    \rA}
  \proofstepwidestar[14]{a\in A\land b\in B\land L(a,b)}{%
    \rAEa{14}}
  \proofstepwidestar[14]{a\in A}{%
    \rAEa{15}}
  \proofstepwidestar[14]{b\in B}{%
    \rAEa{\rAEb{15}}}
  \proofstepwidestar[14]{L(a,b)}{%
    \rAEb{\rAEb{15}}}
  \proofstepwidestar[14]{H(a,b)}{%
    \rAEa{\rAEb{14}}}
  \proofstepwidestar[14]{T(a,b)}{%
    \rAEb{\rAEb{14}}}
  \proofstepwidestar[14]{H(a,b)\land L(a,b)\land T(a,b)}{%
    \rAIStar{19,18,20}}
  \proofstepwidestar[14]{\mathcal E}{%
    \rEIStar{16,17,21}}
  \proofstepwidestar[13]{\mathcal E}{%
    \rEEStar{13,14,22}}
  \proofstepwidestar[]{\mathcal G\rightarrow\mathcal E}{%
    \rRI{13,23}}
  \proofstepwidestar[]{\mathcal E\leftrightarrow\mathcal G}{%
    \rLRI{12,24}}
\end{tabproofwide}

\FormulaThmDeltaK[Beschränkte Allquantoren über ein Paar]{%
  \begin{aligned}[t]&\forall a\in A\,\forall b\in B\,H(a,b)\\[-2pt]&\quad\leftrightarrow \forall a\forall b\,((a\in A\land b\in B)\rightarrow H(a,b))\end{aligned}
}{BoundedPairForallGuard}{%
  \DeltaRow{Mengen}{A\dsep B}
  \DeltaRow{Prädikate}{H\dsep L\dsep T}
}
\begin{tabproofwide}
  \proofstepwidestar[1]{\mathcal E}{%
    \rA}
  \proofstepwidestar[2]{a\in A\land b\in B}{%
    \rA}
  \proofstepwidestar[2]{a\in A}{%
    \rAEa{2}}
  \proofstepwidestar[2]{b\in B}{%
    \rAEb{2}}
  \proofstepwidestar[1,2]{\forall b\in B\,H(a,b)}{%
    \rRE{\rUE{1},3}}
  \proofstepwidestar[1,2]{H(a,b)}{%
    \rRE{\rUE{5},4}}
  \proofstepwidestar[1]{(a\in A\land b\in B)\rightarrow H(a,b)}{%
    \rRI{2,6}}
  \proofstepwidestar[1]{\forall b\,((a\in A\land b\in B)\rightarrow H(a,b))}{%
    \rUI{7}}
  \proofstepwidestar[1]{\mathcal G}{%
    \rUI{8}}
  \proofstepwidestar[]{\mathcal E\rightarrow\mathcal G}{%
    \rRI{1,9}}
  \proofstepwidestar[11]{\mathcal G}{%
    \rA}
  \proofstepwidestar[12]{a\in A}{%
    \rA}
  \proofstepwidestar[13]{b\in B}{%
    \rA}
  \proofstepwidestar[12,13]{a\in A\land b\in B}{%
    \rAI{12,13}}
  \proofstepwidestar[11]{(a\in A\land b\in B)\rightarrow H(a,b)}{%
    \rUE{\rUE{11}}}
  \proofstepwidestar[11,12,13]{H(a,b)}{%
    \rRE{15,14}}
  \proofstepwidestar[11,12]{b\in B\rightarrow H(a,b)}{%
    \rRI{13,16}}
  \proofstepwidestar[11,12]{\forall b\in B\,H(a,b)}{%
    \rUI{17}}
  \proofstepwidestar[11]{a\in A\rightarrow\forall b\in B\,H(a,b)}{%
    \rRI{12,18}}
  \proofstepwidestar[11]{\mathcal E}{%
    \rUI{19}}
  \proofstepwidestar[]{\mathcal G\rightarrow\mathcal E}{%
    \rRI{11,20}}
  \proofstepwidestar[]{\mathcal E\leftrightarrow\mathcal G}{%
    \rLRI{10,21}}
\end{tabproofwide}

\FormulaThmDeltaK[Eindeutige Existenz bei punktweise äquivalenten Prädikaten]{%
  \forall z\,(P(z)\leftrightarrow Q(z))\dsep\exists!z\,P(z)
  \vdash\exists!z\,Q(z)
}{UniqueExistenceUnderPointwiseEquivalence}{%
  \DeltaRow{Prädikate}{P\dsep Q}
  \DeltaRow{Frische Variablen}{a\dsep b}
}
\begin{tabproofwide}
  \proofstepwidestar[1]{\forall z\,(P(z)\leftrightarrow Q(z))}{%
    \rA}
  \proofstepwidestar[2]{\exists!z\,P(z)}{%
    \rA}
  \proofstepwidestar[2]{\exists z\,P(z)}{%
    \FormulaRefAuto{\exists! x\,P(x)\vdash \exists x\,P(x)}{2}}
  \proofstepwidestar[4]{P(a)}{%
    \rA}
  \proofstepwidestar[1]{P(a)\leftrightarrow Q(a)}{%
    \rUE{1}}
  \proofstepwidestar[1,4]{Q(a)}{%
    \FormulaRefAuto{P\leftrightarrow Q,P\vdash Q}{5,4}}
  \proofstepwidestar[1,4]{\exists z\,Q(z)}{%
    \rEI{6}}
  \proofstepwidestar[1,2]{\exists z\,Q(z)}{%
    \rEE{3,4,7}}
  \proofstepwidestar[9]{Q(a)}{%
    \rA}
  \proofstepwidestar[10]{Q(b)}{%
    \rA}
  \proofstepwidestar[1]{P(a)\leftrightarrow Q(a)}{%
    \rUE{1}}
  \proofstepwidestar[1]{P(b)\leftrightarrow Q(b)}{%
    \rUE{1}}
  \proofstepwidestar[1,9]{P(a)}{%
    \FormulaRefAuto{P\leftrightarrow Q,Q\vdash P}{11,9}}
  \proofstepwidestar[1,10]{P(b)}{%
    \FormulaRefAuto{P\leftrightarrow Q,Q\vdash P}{12,10}}
  \proofstepwidestar[1,2,9,10]{a=b}{%
    \FormulaRefAuto{\exists! x\,P(x),\; P(a),\; P(b) \vdash a = b}{2,13,14}}
  \proofstepwidestar[1,2]{\exists!z\,Q(z)}{%
    \UEI{8,9,10,15}}
\end{tabproofwide}

Für beschränkte eindeutige Existenz wird das letzte Schema auf die
Prädikate \(z\in B\land P(z)\) und \(z\in B\land Q(z)\)
angewandt. Auch dann ist zuerst die punktweise Äquivalenz für alle
\(z\) nachzuweisen. Eine unter einer offenen Einzelannahme gewonnene
Äquivalenz genügt nicht.

\FormulaThmDeltaK[Prädikatwechsel bei eindeutiger Existenz mit fester Zusatzbedingung]{%
  \begin{aligned}[t]
    &\forall z\,(P(z)\leftrightarrow Q(z))\dsep
      \exists!z\,(C(z)\land P(z))\\[-2pt]
    &\qquad\vdash\exists!z\,(C(z)\land Q(z))
  \end{aligned}
}{UniqueExistenceUnderEquivalentConjuncts}{%
  \DeltaRow{Prädikate}{C\dsep P\dsep Q}
}
\begin{tabproofwide}
  \proofstepwidestar[1]{\forall z\,(P(z)\leftrightarrow Q(z))}{\rA}
  \proofstepwidestar[2]{\exists!z\,(C(z)\land P(z))}{\rA}
  \proofstepwidestar[3]{C(z)\land P(z)}{\rA}
  \proofstepwidestar[3]{C(z)}{\rAEa{3}}
  \proofstepwidestar[3]{P(z)}{\rAEb{3}}
  \proofstepwidestar[1]{P(z)\leftrightarrow Q(z)}{\rUE{1}}
  \proofstepwidestar[1,3]{Q(z)}{\FormulaRefAuto{P\leftrightarrow Q,P\vdash Q}{6,5}}
  \proofstepwidestar[1,3]{C(z)\land Q(z)}{\rAI{4,7}}
  \proofstepwidestar[1]{(C(z)\land P(z))\rightarrow(C(z)\land Q(z))}{\rRI{3,8}}
  \proofstepwidestar[10]{C(z)\land Q(z)}{\rA}
  \proofstepwidestar[10]{C(z)}{\rAEa{10}}
  \proofstepwidestar[10]{Q(z)}{\rAEb{10}}
  \proofstepwidestar[1]{P(z)\leftrightarrow Q(z)}{\rUE{1}}
  \proofstepwidestar[1,10]{P(z)}{\FormulaRefAuto{P\leftrightarrow Q,Q\vdash P}{13,12}}
  \proofstepwidestar[1,10]{C(z)\land P(z)}{\rAI{11,14}}
  \proofstepwidestar[1]{(C(z)\land Q(z))\rightarrow(C(z)\land P(z))}{\rRI{10,15}}
  \proofstepwidestar[1]{(C(z)\land P(z))\leftrightarrow(C(z)\land Q(z))}{\rLRI{9,16}}
  \proofstepwidestar[1]{\forall z\,((C(z)\land P(z))\leftrightarrow(C(z)\land Q(z)))}{\rUI{17}}
  \proofstepwidestar[1,2]{\exists!z\,(C(z)\land Q(z))}{%
    \FormulaRefAuto{UniqueExistenceUnderPointwiseEquivalence}[thm]{18,2}}
\end{tabproofwide}

\subsection{Eindeutige Werte und Faktorabbildungen}


Eine Vorschrift auf Repräsentanten bestimmt genau dann einen Wert auf
jedem Bildpunkt, wenn sie auf dessen Faser konstant ist. Die folgenden
Schemata halten diese Konstruktion ohne Auswahl von Repräsentanten fest.
Die Ausdrücke \(\kappa(a,b)\) und \(\tau(a,b)\), beziehungsweise
\(\tau(a,b,c,d)\), sind bereits gebildete Terme. Die angegebenen
Allquantoren binden deren ausgezeichnete Variablen; alle anderen Parameter
bleiben fest. Gebundene Variablen werden vor Einsetzungen so umbenannt,
dass keine freien Parameter gebunden werden. Insbesondere sind die
Zielvariablen \(u,v,z,w\) und die neu eingeführten Funktionsvariablen
frisch für diese Parameter. \(P\) darf die Zugehörigkeit der beiden
Komponenten zu verschiedenen Mengen ausdrücken. Die Variablen der
Existenzelimination werden stets frisch gewählt.


\FormulaThmDeltaK[Eindeutiger Wert aus Paarrepräsentanten]{%
  \begin{aligned}[t]&\forall a\forall b\,\bigl(P(a,b)\rightarrow\kappa(a,b)\in Q\bigr)\dsep{}\\[-2pt]&\forall u\in Q\,\exists a\exists b\,(P(a,b)\land u=\kappa(a,b))\dsep{}\\[-2pt]&\forall a\forall b\,\bigl(P(a,b)\rightarrow \tau(a,b)\in B\bigr)\dsep{}\\[-2pt]&\forall a\forall b\forall c\forall d\,\bigl((P(a,b)\land P(c,d)\land \kappa(a,b)=\kappa(c,d))\rightarrow \tau(a,b)=\tau(c,d)\bigr)\dsep{}\\[-2pt]&u\in Q\vdash{}\\[-2pt]&\exists!z\in B\,\exists a\exists b\,\bigl(P(a,b)\land u=\kappa(a,b)\land z=\tau(a,b)\bigr)\end{aligned}
}{RepresentativePairValueUnique}{%
  \DeltaRow{Mengen}{Q\dsep B\dsep u\dsep v\dsep z\dsep w}
  \DeltaRow{Prädikat}{P}
  \DeltaRow{Termschemata}{\kappa(a,b)\dsep\tau(a,b)}
}
\begin{tabproofsplitwide}
\proofpartwide{Existenz eines Werts}
  \proofstepwidestar[1]{\forall a\forall b\,\bigl(P(a,b)\rightarrow\kappa(a,b)\in Q\bigr)}{%
    \rA}
  \proofstepwidestar[2]{\forall u\in Q\,\exists a\exists b\,(P(a,b)\land u=\kappa(a,b))}{%
    \rA}
  \proofstepwidestar[3]{\forall a\forall b\,\bigl(P(a,b)\rightarrow \tau(a,b)\in B\bigr)}{%
    \rA}
  \proofstepwidestar[4]{\forall a\forall b\forall c\forall d\,\bigl((P(a,b)\land P(c,d)\land \kappa(a,b)=\kappa(c,d))\rightarrow \tau(a,b)=\tau(c,d)\bigr)}{%
    \rA}
  \proofstepwidestar[5]{u\in Q}{%
    \rA}
  \proofstepwidestar[2]{u\in Q\rightarrow\exists a\exists b\,(P(a,b)\land u=\kappa(a,b))}{%
    \rUE{2}}
  \proofstepwidestar[2,5]{\exists a\exists b\,(P(a,b)\land u=\kappa(a,b))}{%
    \rRE{6,5}}
  \proofstepwidestar[8]{P(a,b)\land u=\kappa(a,b)}{%
    \rA}
  \proofstepwidestar[8]{P(a,b)}{%
    \rAEa{8}}
  \proofstepwidestar[8]{u=\kappa(a,b)}{%
    \rAEb{8}}
  \proofstepwidestar[3]{P(a,b)\rightarrow \tau(a,b)\in B}{%
    \rUE{\rUE{3}}}
  \proofstepwidestar[3,8]{\tau(a,b)\in B}{%
    \rRE{11,9}}
  \proofstepwidestar[]{\tau(a,b)=\tau(a,b)}{%
    \rII}
  \proofstepwidestar[8]{P(a,b)\land u=\kappa(a,b)\land \tau(a,b)=\tau(a,b)}{%
    \rAIStar{9,10,13}}
  \proofstepwidestar[8]{\exists a_0\exists b_0\,\bigl(P(a_0,b_0)\land u=\kappa(a_0,b_0)\land \tau(a,b)=\tau(a_0,b_0)\bigr)}{%
    \rEIStar{14}}
  \proofstepwidestar[3,8]{\tau(a,b)\in B\land \exists a_0\exists b_0\,\bigl(P(a_0,b_0)\land u=\kappa(a_0,b_0)\land \tau(a,b)=\tau(a_0,b_0)\bigr)}{%
    \rAIStar{12,15}}
  \proofstepwidestar[3,8]{\begin{aligned}[t]&\exists z\in B\,\exists a\exists b\,\bigl(P(a,b)\\[-2pt]&\quad\land u=\kappa(a,b)\land z=\tau(a,b)\bigr)\end{aligned}}{%
    \rEIStar{16}}
  \proofstepwidestar[2,3,5]{\begin{aligned}[t]&\exists z\in B\,\exists a\exists b\,\bigl(P(a,b)\\[-2pt]&\quad\land u=\kappa(a,b)\land z=\tau(a,b)\bigr)\end{aligned}}{%
    \rEEStar{7,8,17}}
\closeproofpartwide
\proofpartwide{Eindeutigkeit des Werts}
\setcounter{proofstepnr}{18}
  \proofstepwidestar[19]{\begin{aligned}[t]&z\in B\land \exists a\exists b\,\bigl(P(a,b)\\[-2pt]&\quad\land u=\kappa(a,b)\land z=\tau(a,b)\bigr)\end{aligned}}{%
    \rA}
  \proofstepwidestar[20]{\begin{aligned}[t]&w\in B\land \exists c\exists d\,\bigl(P(c,d)\\[-2pt]&\quad\land u=\kappa(c,d)\land w=\tau(c,d)\bigr)\end{aligned}}{%
    \rA}
  \proofstepwidestar[19]{\exists a\exists b\,\bigl(P(a,b)\land u=\kappa(a,b)\land z=\tau(a,b)\bigr)}{%
    \rAEb{19}}
  \proofstepwidestar[20]{\exists c\exists d\,\bigl(P(c,d)\land u=\kappa(c,d)\land w=\tau(c,d)\bigr)}{%
    \rAEb{20}}
  \proofstepwidestar[23]{P(a,b)\land u=\kappa(a,b)\land z=\tau(a,b)}{%
    \rA}
  \proofstepwidestar[24]{P(c,d)\land u=\kappa(c,d)\land w=\tau(c,d)}{%
    \rA}
  \proofstepwidestar[23]{P(a,b)}{%
    \rAEa{23}}
  \proofstepwidestar[24]{P(c,d)}{%
    \rAEa{24}}
  \proofstepwidestar[23]{u=\kappa(a,b)}{%
    \rAEa{\rAEb{23}}}
  \proofstepwidestar[24]{u=\kappa(c,d)}{%
    \rAEa{\rAEb{24}}}
  \proofstepwidestar[23,24]{\kappa(a,b)=\kappa(c,d)}{%
    \FormulaRefAuto{a=b\dsep a=c\vdash b=c}[thm]{27,28}}
  \proofstepwidestar[23]{z=\tau(a,b)}{%
    \rAEb{\rAEb{23}}}
  \proofstepwidestar[24]{w=\tau(c,d)}{%
    \rAEb{\rAEb{24}}}
  \proofstepwidestar[4]{\begin{aligned}[t]&(P(a,b)\land P(c,d)\land \kappa(a,b)=\kappa(c,d))\\[-2pt]&\quad\rightarrow \tau(a,b)=\tau(c,d)\end{aligned}}{%
    \rUE{\rUE{\rUE{\rUE{4}}}}}
  \proofstepwidestar[23,24]{P(a,b)\land P(c,d)\land \kappa(a,b)=\kappa(c,d)}{%
    \rAIStar{25,26,29}}
  \proofstepwidestar[4,23,24]{\tau(a,b)=\tau(c,d)}{%
    \rRE{32,33}}
  \proofstepwidestar[4,23,24]{z=\tau(c,d)}{%
    \rIE{34,30}}
  \proofstepwidestar[4,23,24]{z=w}{%
    \FormulaRefAuto{a=b\dsep c=b\vdash a=c}[thm]{35,31}}
  \proofstepwidestar[4,20,23]{z=w}{%
    \rEEStar{22,24,36}}
  \proofstepwidestar[4,19,20]{z=w}{%
    \rEEStar{21,23,37}}
\closeproofpartwide
\proofpartwide{Eindeutige Existenz}
\setcounter{proofstepnr}{38}
  \proofstepwidestar[2,3,4,5]{\begin{aligned}[t]&\exists!z\in B\,\exists a\exists b\,\bigl(P(a,b)\\[-2pt]&\quad\land u=\kappa(a,b)\land z=\tau(a,b)\bigr)\end{aligned}}{%
    \UEI{18,19,20,38}}
\closeproofpartwide
\end{tabproofsplitwide}


\FormulaThmDeltaK[Eindeutiger binärer Wert aus Paarrepräsentanten]{%
  \begin{aligned}[t]&\forall a\forall b\,\bigl(P(a,b)\rightarrow\kappa(a,b)\in Q\bigr)\dsep{}\\[-2pt]&\forall u\in Q\,\exists a\exists b\,(P(a,b)\land u=\kappa(a,b))\dsep{}\\[-2pt]&\forall a\forall b\forall c\forall d\,\bigl(P(a,b)\land P(c,d)\rightarrow \tau(a,b,c,d)\in B\bigr)\dsep{}\\[-2pt]&\begin{aligned}[t]&\forall a\forall b\forall c\forall d\forall e\forall f\forall g\forall h\,\bigl((P(a,b)\land P(c,d)\\[-2pt]&\quad\land P(e,f)\land P(g,h)\\[-2pt]&\quad\land \kappa(a,b)=\kappa(e,f)\land \kappa(c,d)=\kappa(g,h))\\[-2pt]&\quad \rightarrow \tau(a,b,c,d)=\tau(e,f,g,h)\bigr)\end{aligned}\dsep{}\\[-2pt]&u\in Q\dsep{}\\[-2pt]&v\in Q\vdash{}\\[-2pt]&\exists!z\in B\,\exists a\exists b\exists c\exists d\,\bigl(P(a,b)\land P(c,d)\\[-2pt]&\quad\land u=\kappa(a,b)\land v=\kappa(c,d)\land z=\tau(a,b,c,d)\bigr)\end{aligned}
}{RepresentativePairBinaryValueUnique}{%
  \DeltaRow{Mengen}{Q\dsep B\dsep u\dsep v\dsep z\dsep w}
  \DeltaRow{Prädikat}{P}
  \DeltaRow{Termschemata}{\kappa(a,b)\dsep\tau(a,b,c,d)}
}
\begin{tabproofsplitwide}
\proofpartwide{Existenz eines Werts}
  \proofstepwidestar[1]{\forall a\forall b\,\bigl(P(a,b)\rightarrow\kappa(a,b)\in Q\bigr)}{%
    \rA}
  \proofstepwidestar[2]{\forall u\in Q\,\exists a\exists b\,(P(a,b)\land u=\kappa(a,b))}{%
    \rA}
  \proofstepwidestar[3]{\forall a\forall b\forall c\forall d\,\bigl(P(a,b)\land P(c,d)\rightarrow \tau(a,b,c,d)\in B\bigr)}{%
    \rA}
  \proofstepwidestar[4]{\begin{aligned}[t]&\forall a\forall b\forall c\forall d\forall e\forall f\forall g\forall h\,\bigl((P(a,b)\land P(c,d)\\[-2pt]&\quad\land P(e,f)\land P(g,h)\\[-2pt]&\quad\land \kappa(a,b)=\kappa(e,f)\land \kappa(c,d)=\kappa(g,h))\\[-2pt]&\quad \rightarrow \tau(a,b,c,d)=\tau(e,f,g,h)\bigr)\end{aligned}}{%
    \rA}
  \proofstepwidestar[5]{u\in Q}{%
    \rA}
  \proofstepwidestar[6]{v\in Q}{%
    \rA}
  \proofstepwidestar[2]{u\in Q\rightarrow\exists a\exists b\,(P(a,b)\land u=\kappa(a,b))}{%
    \rUE{2}}
  \proofstepwidestar[2,5]{\exists a\exists b\,(P(a,b)\land u=\kappa(a,b))}{%
    \rRE{7,5}}
  \proofstepwidestar[9]{P(a,b)\land u=\kappa(a,b)}{%
    \rA}
  \proofstepwidestar[9]{P(a,b)}{%
    \rAEa{9}}
  \proofstepwidestar[9]{u=\kappa(a,b)}{%
    \rAEb{9}}
  \proofstepwidestar[2]{v\in Q\rightarrow\exists c\exists d\,(P(c,d)\land v=\kappa(c,d))}{%
    \rUE{2}}
  \proofstepwidestar[2,6]{\exists c\exists d\,(P(c,d)\land v=\kappa(c,d))}{%
    \rRE{12,6}}
  \proofstepwidestar[14]{P(c,d)\land v=\kappa(c,d)}{%
    \rA}
  \proofstepwidestar[14]{P(c,d)}{%
    \rAEa{14}}
  \proofstepwidestar[14]{v=\kappa(c,d)}{%
    \rAEb{14}}
  \proofstepwidestar[9,14]{P(a,b)\land P(c,d)}{%
    \rAIStar{10,15}}
  \proofstepwidestar[3]{P(a,b)\land P(c,d)\rightarrow \tau(a,b,c,d)\in B}{%
    \rUE{\rUE{\rUE{\rUE{3}}}}}
  \proofstepwidestar[3,9,14]{\tau(a,b,c,d)\in B}{%
    \rRE{18,17}}
  \proofstepwidestar[]{\tau(a,b,c,d)=\tau(a,b,c,d)}{%
    \rII}
  \proofstepwidestar[9,14]{P(a,b)\land P(c,d)\land u=\kappa(a,b)\land v=\kappa(c,d)\land \tau(a,b,c,d)=\tau(a,b,c,d)}{%
    \rAIStar{10,15,11,16,20}}
  \proofstepwidestar[9,14]{\begin{aligned}[t]&\exists a_0\exists b_0\exists c_0\exists d_0\,\bigl(P(a_0,b_0)\land P(c_0,d_0)\\[-2pt]&\quad\land u=\kappa(a_0,b_0)\land v=\kappa(c_0,d_0)\\[-2pt]&\quad\land \tau(a,b,c,d)=\tau(a_0,b_0,c_0,d_0)\bigr)\end{aligned}}{%
    \rEIStar{21}}
  \proofstepwidestar[3,9,14]{\begin{aligned}[t]&\tau(a,b,c,d)\in B\land \exists a_0\exists b_0\exists c_0\exists d_0\\[-2pt]&\quad\bigl(P(a_0,b_0)\land P(c_0,d_0)\\[-2pt]&\quad \land u=\kappa(a_0,b_0)\land v=\kappa(c_0,d_0)\\[-2pt]&\quad\land \tau(a,b,c,d)=\tau(a_0,b_0,c_0,d_0)\bigr)\end{aligned}}{%
    \rAIStar{19,22}}
  \proofstepwidestar[3,9,14]{\exists z\in B\,\exists a\exists b\exists c\exists d\,\bigl(P(a,b)\land P(c,d)\land u=\kappa(a,b)\land v=\kappa(c,d)\land z=\tau(a,b,c,d)\bigr)}{%
    \rEIStar{23}}
  \proofstepwidestar[2,3,6,9]{\exists z\in B\,\exists a\exists b\exists c\exists d\,\bigl(P(a,b)\land P(c,d)\land u=\kappa(a,b)\land v=\kappa(c,d)\land z=\tau(a,b,c,d)\bigr)}{%
    \rEEStar{13,14,24}}
  \proofstepwidestar[2,3,5,6]{\exists z\in B\,\exists a\exists b\exists c\exists d\,\bigl(P(a,b)\land P(c,d)\land u=\kappa(a,b)\land v=\kappa(c,d)\land z=\tau(a,b,c,d)\bigr)}{%
    \rEEStar{8,9,25}}
\closeproofpartwide
\proofpartwide{Eindeutigkeit des Werts}
\setcounter{proofstepnr}{26}
  \proofstepwidestar[27]{z\in B\land \exists a\exists b\exists c\exists d\,\bigl(P(a,b)\land P(c,d)\land u=\kappa(a,b)\land v=\kappa(c,d)\land z=\tau(a,b,c,d)\bigr)}{%
    \rA}
  \proofstepwidestar[28]{w\in B\land \exists e\exists f\exists g\exists h\,\bigl(P(e,f)\land P(g,h)\land u=\kappa(e,f)\land v=\kappa(g,h)\land w=\tau(e,f,g,h)\bigr)}{%
    \rA}
  \proofstepwidestar[27]{\exists a\exists b\exists c\exists d\,\bigl(P(a,b)\land P(c,d)\land u=\kappa(a,b)\land v=\kappa(c,d)\land z=\tau(a,b,c,d)\bigr)}{%
    \rAEb{27}}
  \proofstepwidestar[28]{\exists e\exists f\exists g\exists h\,\bigl(P(e,f)\land P(g,h)\land u=\kappa(e,f)\land v=\kappa(g,h)\land w=\tau(e,f,g,h)\bigr)}{%
    \rAEb{28}}
  \proofstepwidestar[31]{P(a,b)\land P(c,d)\land u=\kappa(a,b)\land v=\kappa(c,d)\land z=\tau(a,b,c,d)}{%
    \rA}
  \proofstepwidestar[32]{P(e,f)\land P(g,h)\land u=\kappa(e,f)\land v=\kappa(g,h)\land w=\tau(e,f,g,h)}{%
    \rA}
  \proofstepwidestar[31]{P(a,b)}{%
    \rAEa{31}}
  \proofstepwidestar[31]{P(c,d)}{%
    \rAEa{\rAEb{31}}}
  \proofstepwidestar[32]{P(e,f)}{%
    \rAEa{32}}
  \proofstepwidestar[32]{P(g,h)}{%
    \rAEa{\rAEb{32}}}
  \proofstepwidestar[31]{u=\kappa(a,b)}{%
    \rAEa{\rAEb{\rAEb{31}}}}
  \proofstepwidestar[32]{u=\kappa(e,f)}{%
    \rAEa{\rAEb{\rAEb{32}}}}
  \proofstepwidestar[31,32]{\kappa(a,b)=\kappa(e,f)}{%
    \FormulaRefAuto{a=b\dsep a=c\vdash b=c}[thm]{37,38}}
  \proofstepwidestar[31]{v=\kappa(c,d)}{%
    \rAEa{\rAEb{\rAEb{\rAEb{31}}}}}
  \proofstepwidestar[32]{v=\kappa(g,h)}{%
    \rAEa{\rAEb{\rAEb{\rAEb{32}}}}}
  \proofstepwidestar[31,32]{\kappa(c,d)=\kappa(g,h)}{%
    \FormulaRefAuto{a=b\dsep a=c\vdash b=c}[thm]{40,41}}
  \proofstepwidestar[31]{z=\tau(a,b,c,d)}{%
    \rAEb{\rAEb{\rAEb{\rAEb{31}}}}}
  \proofstepwidestar[32]{w=\tau(e,f,g,h)}{%
    \rAEb{\rAEb{\rAEb{\rAEb{32}}}}}
  \proofstepwidestar[4]{(P(a,b)\land P(c,d)\land P(e,f)\land P(g,h)\land \kappa(a,b)=\kappa(e,f)\land \kappa(c,d)=\kappa(g,h))\rightarrow \tau(a,b,c,d)=\tau(e,f,g,h)}{%
    \rUE{\rUE{\rUE{\rUE{\rUE{\rUE{\rUE{\rUE{4}}}}}}}}}
  \proofstepwidestar[31,32]{P(a,b)\land P(c,d)\land P(e,f)\land P(g,h)\land \kappa(a,b)=\kappa(e,f)\land \kappa(c,d)=\kappa(g,h)}{%
    \rAIStar{33,34,35,36,39,42}}
  \proofstepwidestar[4,31,32]{\tau(a,b,c,d)=\tau(e,f,g,h)}{%
    \rRE{45,46}}
  \proofstepwidestar[4,31,32]{z=\tau(e,f,g,h)}{%
    \rIE{47,43}}
  \proofstepwidestar[4,31,32]{z=w}{%
    \FormulaRefAuto{a=b\dsep c=b\vdash a=c}[thm]{48,44}}
  \proofstepwidestar[4,28,31]{z=w}{%
    \rEEStar{30,32,49}}
  \proofstepwidestar[4,27,28]{z=w}{%
    \rEEStar{29,31,50}}
\closeproofpartwide
\Needspace{7\baselineskip}
\proofpartwide{Eindeutige Existenz}
\setcounter{proofstepnr}{51}
  \proofstepwidestar[2,3,4,5,6]{\exists!z\in B\,\exists a\exists b\exists c\exists d\,\bigl(P(a,b)\land P(c,d)\land u=\kappa(a,b)\land v=\kappa(c,d)\land z=\tau(a,b,c,d)\bigr)}{%
    \UEI{26,27,28,51}}
\closeproofpartwide
\end{tabproofsplitwide}


\FormulaThmDeltaK[Eindeutiger Wert auf Surjektionsfasern]{%
  \begin{aligned}[t]&q\colon S\sur Q\dsep{}\\[-2pt]&h\colon S\to B\dsep{}\\[-2pt]&\forall s\in S\,\forall s^{\prime}\in S\,\bigl((q(s)=q(s^{\prime}))\rightarrow h(s)=h(s^{\prime})\bigr)\dsep{}\\[-2pt]&u\in Q\vdash{}\\[-2pt]&\exists!z\in B\,\exists s\in S\,\bigl(u=q(s)\land z=h(s)\bigr)\end{aligned}
}{SurjectionFiberValueUnique}{%
  \DeltaRow{Mengen}{S\dsep Q\dsep B\dsep u\dsep v\dsep z\dsep w}
  \DeltaRow{Funktionen}{q\dsep h}
}
\begin{tabproofsplitwide}
\proofpartwide{Existenz eines Werts}
  \proofstepwidestar[1]{q\colon S\sur Q}{%
    \rA}
  \proofstepwidestar[2]{h\colon S\to B}{%
    \rA}
  \proofstepwidestar[3]{\forall s\in S\,\forall s^{\prime}\in S\,\bigl((q(s)=q(s^{\prime}))\rightarrow h(s)=h(s^{\prime})\bigr)}{%
    \rA}
  \proofstepwidestar[4]{u\in Q}{%
    \rA}
  \proofstepwidestar[1]{q\colon S\to Q}{%
    \FormulaRefAuto{Surjektive Funktion}[def]{1}}
  \proofstepwidestar[1,4]{\exists s\in S\,q(s)=u}{%
    \FormulaRefAuto{Surjektivität}[ax]{1,4}}
  \proofstepwidestar[7]{s\in S\land q(s)=u}{%
    \rA}
  \proofstepwidestar[7]{s\in S}{%
    \rAEa{7}}
  \proofstepwidestar[7]{q(s)=u}{%
    \rAEb{7}}
  \proofstepwidestar[7]{u=q(s)}{%
    \FormulaRefAuto{a=b\vdash b=a}[thm]{9}}
  \proofstepwidestar[2,7]{h(s)\in B}{%
    \FormulaRefAuto{F\colon A\to B\dsep x\in A\vdash F(x)\in B}[thm]{2,8}}
  \proofstepwidestar[]{h(s)=h(s)}{%
    \rII}
  \proofstepwidestar[7]{u=q(s)\land h(s)=h(s)}{%
    \rAIStar{10,12}}
  \proofstepwidestar[7]{\exists s_0\in S\,\bigl(u=q(s_0)\land h(s)=h(s_0)\bigr)}{%
    \rEIStar{8,13}}
  \proofstepwidestar[2,7]{h(s)\in B\land \exists s_0\in S\,\bigl(u=q(s_0)\land h(s)=h(s_0)\bigr)}{%
    \rAIStar{11,14}}
  \proofstepwidestar[2,7]{\exists z\in B\,\exists s\in S\,\bigl(u=q(s)\land z=h(s)\bigr)}{%
    \rEIStar{15}}
  \proofstepwidestar[1,2,4]{\exists z\in B\,\exists s\in S\,\bigl(u=q(s)\land z=h(s)\bigr)}{%
    \rEEStar{6,7,16}}
\closeproofpartwide
\proofpartwide{Eindeutigkeit des Werts}
\setcounter{proofstepnr}{17}
  \proofstepwidestar[18]{z\in B\land \exists s\in S\,\bigl(u=q(s)\land z=h(s)\bigr)}{%
    \rA}
  \proofstepwidestar[19]{w\in B\land \exists s^{\prime}\in S\,\bigl(u=q(s^{\prime})\land w=h(s^{\prime})\bigr)}{%
    \rA}
  \proofstepwidestar[18]{\exists s\in S\,\bigl(u=q(s)\land z=h(s)\bigr)}{%
    \rAEb{18}}
  \proofstepwidestar[19]{\exists s^{\prime}\in S\,\bigl(u=q(s^{\prime})\land w=h(s^{\prime})\bigr)}{%
    \rAEb{19}}
  \proofstepwidestar[22]{s\in S\land u=q(s)\land z=h(s)}{%
    \rA}
  \proofstepwidestar[23]{s^{\prime}\in S\land u=q(s^{\prime})\land w=h(s^{\prime})}{%
    \rA}
  \proofstepwidestar[22]{s\in S}{%
    \rAEa{22}}
  \proofstepwidestar[23]{s^{\prime}\in S}{%
    \rAEa{23}}
  \proofstepwidestar[22]{u=q(s)}{%
    \rAEa{\rAEb{22}}}
  \proofstepwidestar[23]{u=q(s^{\prime})}{%
    \rAEa{\rAEb{23}}}
  \proofstepwidestar[22,23]{q(s)=q(s^{\prime})}{%
    \FormulaRefAuto{a=b\dsep a=c\vdash b=c}[thm]{26,27}}
  \proofstepwidestar[22]{z=h(s)}{%
    \rAEb{\rAEb{22}}}
  \proofstepwidestar[23]{w=h(s^{\prime})}{%
    \rAEb{\rAEb{23}}}
  \proofstepwidestar[3]{s\in S\rightarrow \forall s^{\prime}\in S\,\bigl((q(s)=q(s^{\prime}))\rightarrow h(s)=h(s^{\prime})\bigr)}{%
    \rUE{3}}
  \proofstepwidestar[3,22]{\forall s^{\prime}\in S\,\bigl((q(s)=q(s^{\prime}))\rightarrow h(s)=h(s^{\prime})\bigr)}{%
    \rRE{31,24}}
  \proofstepwidestar[3,22]{s^{\prime}\in S\rightarrow \bigl((q(s)=q(s^{\prime}))\rightarrow h(s)=h(s^{\prime})\bigr)}{%
    \rUE{32}}
  \proofstepwidestar[3,22,23]{\bigl((q(s)=q(s^{\prime}))\rightarrow h(s)=h(s^{\prime})\bigr)}{%
    \rRE{33,25}}
  \proofstepwidestar[3,22,23]{h(s)=h(s^{\prime})}{%
    \rRE{34,28}}
  \proofstepwidestar[3,22,23]{z=h(s^{\prime})}{%
    \rIE{35,29}}
  \proofstepwidestar[3,22,23]{z=w}{%
    \FormulaRefAuto{a=b\dsep c=b\vdash a=c}[thm]{36,30}}
  \proofstepwidestar[3,19,22]{z=w}{%
    \rEEStar{21,23,37}}
  \proofstepwidestar[3,18,19]{z=w}{%
    \rEEStar{20,22,38}}
\closeproofpartwide
\proofpartwide{Eindeutige Existenz}
\setcounter{proofstepnr}{39}
  \proofstepwidestar[1,2,3,4]{\exists!z\in B\,\exists s\in S\,\bigl(u=q(s)\land z=h(s)\bigr)}{%
    \UEI{17,18,19,39}}
\closeproofpartwide
\end{tabproofsplitwide}


\FormulaThmDeltaK[Eindeutiger binärer Wert auf Surjektionsfasern]{%
  \begin{aligned}[t]&q\colon S\sur Q\dsep{}\\[-2pt]&h\colon S\times S\to B\dsep{}\\[-2pt]&\begin{aligned}[t]&\forall s\in S\,\forall t\in S\,\forall s^{\prime}\in S\,\forall t^{\prime}\in S\,\bigl((q(s)=q(s^{\prime})\\[-2pt]&\quad \land q(t)=q(t^{\prime}))\rightarrow h(s,t)=h(s^{\prime},t^{\prime})\bigr)\end{aligned}\dsep{}\\[-2pt]&u\in Q\dsep{}\\[-2pt]&v\in Q\vdash{}\\[-2pt]&\exists!z\in B\,\exists s\in S\,\exists t\in S\,\bigl(u=q(s)\land v=q(t)\land z=h(s,t)\bigr)\end{aligned}
}{SurjectionBinaryFiberValueUnique}{%
  \DeltaRow{Mengen}{S\dsep Q\dsep B\dsep u\dsep v\dsep z\dsep w}
  \DeltaRow{Funktionen}{q\dsep h}
}
\Needspace{8\baselineskip}
\begin{tabproofsplitwide}
\proofpartwide{Existenz eines Werts}
  \proofstepwidestar[1]{q\colon S\sur Q}{%
    \rA}
  \proofstepwidestar[2]{h\colon S\times S\to B}{%
    \rA}
  \proofstepwidestar[3]{\begin{aligned}[t]&\forall s\in S\,\forall t\in S\,\forall s^{\prime}\in S\,\forall t^{\prime}\in S\,\bigl((q(s)=q(s^{\prime})\\[-2pt]&\quad \land q(t)=q(t^{\prime}))\rightarrow h(s,t)=h(s^{\prime},t^{\prime})\bigr)\end{aligned}}{%
    \rA}
  \proofstepwidestar[4]{u\in Q}{%
    \rA}
  \proofstepwidestar[5]{v\in Q}{%
    \rA}
  \proofstepwidestar[1]{q\colon S\to Q}{%
    \FormulaRefAuto{Surjektive Funktion}[def]{1}}
  \proofstepwidestar[1,4]{\exists s\in S\,q(s)=u}{%
    \FormulaRefAuto{Surjektivität}[ax]{1,4}}
  \proofstepwidestar[8]{s\in S\land q(s)=u}{%
    \rA}
  \proofstepwidestar[8]{s\in S}{%
    \rAEa{8}}
  \proofstepwidestar[8]{q(s)=u}{%
    \rAEb{8}}
  \proofstepwidestar[8]{u=q(s)}{%
    \FormulaRefAuto{a=b\vdash b=a}[thm]{10}}
  \proofstepwidestar[1,5]{\exists t\in S\,q(t)=v}{%
    \FormulaRefAuto{Surjektivität}[ax]{1,5}}
  \proofstepwidestar[13]{t\in S\land q(t)=v}{%
    \rA}
  \proofstepwidestar[13]{t\in S}{%
    \rAEa{13}}
  \proofstepwidestar[13]{q(t)=v}{%
    \rAEb{13}}
  \proofstepwidestar[13]{v=q(t)}{%
    \FormulaRefAuto{a=b\vdash b=a}[thm]{15}}
  \proofstepwidestar[8,13]{(s,t)\in S\times S}{%
    \FormulaRefAuto{a\in A\dsep b\in B\vdash(a,b)\in A\times B}[thm]{9,14}}
  \proofstepwidestar[2,8,13]{h(s,t)\in B}{%
    \FormulaRefAuto{F\colon A\to B\dsep x\in A\vdash F(x)\in B}[thm]{2,17}}
  \proofstepwidestar[]{h(s,t)=h(s,t)}{%
    \rII}
  \proofstepwidestar[8,13]{u=q(s)\land v=q(t)\land h(s,t)=h(s,t)}{%
    \rAIStar{11,16,19}}
  \proofstepwidestar[8,13]{\begin{aligned}[t]&\exists s_0\in S\,\exists t_0\in S\,\bigl(u=q(s_0)\\[-2pt]&\quad\land v=q(t_0)\land h(s,t)=h(s_0,t_0)\bigr)\end{aligned}}{%
    \rEIStar{9,14,20}}
  \proofstepwidestar[2,8,13]{h(s,t)\in B\land \exists s_0\in S\,\exists t_0\in S\,\bigl(u=q(s_0)\land v=q(t_0)\land h(s,t)=h(s_0,t_0)\bigr)}{%
    \rAIStar{18,21}}
  \proofstepwidestar[2,8,13]{\exists z\in B\,\exists s\in S\,\exists t\in S\,\bigl(u=q(s)\land v=q(t)\land z=h(s,t)\bigr)}{%
    \rEIStar{22}}
  \proofstepwidestar[1,2,5,8]{\exists z\in B\,\exists s\in S\,\exists t\in S\,\bigl(u=q(s)\land v=q(t)\land z=h(s,t)\bigr)}{%
    \rEEStar{12,13,23}}
  \proofstepwidestar[1,2,4,5]{\exists z\in B\,\exists s\in S\,\exists t\in S\,\bigl(u=q(s)\land v=q(t)\land z=h(s,t)\bigr)}{%
    \rEEStar{7,8,24}}
\closeproofpartwide
\proofpartwide{Eindeutigkeit des Werts}
\setcounter{proofstepnr}{25}
  \proofstepwidestar[26]{z\in B\land \exists s\in S\,\exists t\in S\,\bigl(u=q(s)\land v=q(t)\land z=h(s,t)\bigr)}{%
    \rA}
  \proofstepwidestar[27]{w\in B\land \exists s^{\prime}\in S\,\exists t^{\prime}\in S\,\bigl(u=q(s^{\prime})\land v=q(t^{\prime})\land w=h(s^{\prime},t^{\prime})\bigr)}{%
    \rA}
  \proofstepwidestar[26]{\exists s\in S\,\exists t\in S\,\bigl(u=q(s)\land v=q(t)\land z=h(s,t)\bigr)}{%
    \rAEb{26}}
  \proofstepwidestar[27]{\exists s^{\prime}\in S\,\exists t^{\prime}\in S\,\bigl(u=q(s^{\prime})\land v=q(t^{\prime})\land w=h(s^{\prime},t^{\prime})\bigr)}{%
    \rAEb{27}}
  \proofstepwidestar[30]{s\in S\land t\in S\land u=q(s)\land v=q(t)\land z=h(s,t)}{%
    \rA}
  \proofstepwidestar[31]{\begin{aligned}[t]&s^{\prime}\in S\land t^{\prime}\in S\land u=q(s^{\prime})\\[-2pt]&\quad\land v=q(t^{\prime})\land w=h(s^{\prime},t^{\prime})\end{aligned}}{%
    \rA}
  \proofstepwidestar[30]{s\in S}{%
    \rAEa{30}}
  \proofstepwidestar[30]{t\in S}{%
    \rAEa{\rAEb{30}}}
  \proofstepwidestar[31]{s^{\prime}\in S}{%
    \rAEa{31}}
  \proofstepwidestar[31]{t^{\prime}\in S}{%
    \rAEa{\rAEb{31}}}
  \proofstepwidestar[30]{u=q(s)}{%
    \rAEa{\rAEb{\rAEb{30}}}}
  \proofstepwidestar[31]{u=q(s^{\prime})}{%
    \rAEa{\rAEb{\rAEb{31}}}}
  \proofstepwidestar[30,31]{q(s)=q(s^{\prime})}{%
    \FormulaRefAuto{a=b\dsep a=c\vdash b=c}[thm]{36,37}}
  \proofstepwidestar[30]{v=q(t)}{%
    \rAEa{\rAEb{\rAEb{\rAEb{30}}}}}
  \proofstepwidestar[31]{v=q(t^{\prime})}{%
    \rAEa{\rAEb{\rAEb{\rAEb{31}}}}}
  \proofstepwidestar[30,31]{q(t)=q(t^{\prime})}{%
    \FormulaRefAuto{a=b\dsep a=c\vdash b=c}[thm]{39,40}}
  \proofstepwidestar[30]{z=h(s,t)}{%
    \rAEb{\rAEb{\rAEb{\rAEb{30}}}}}
  \proofstepwidestar[31]{w=h(s^{\prime},t^{\prime})}{%
    \rAEb{\rAEb{\rAEb{\rAEb{31}}}}}
  \proofstepwidestar[3]{\begin{aligned}[t]&s\in S\rightarrow \forall t\in S\,\forall s^{\prime}\in S\,\forall t^{\prime}\in S\,\bigl((q(s)=q(s^{\prime})\\[-2pt]&\quad \land q(t)=q(t^{\prime}))\rightarrow h(s,t)=h(s^{\prime},t^{\prime})\bigr)\end{aligned}}{%
    \rUE{3}}
  \proofstepwidestar[3,30]{\begin{aligned}[t]&\forall t\in S\,\forall s^{\prime}\in S\,\forall t^{\prime}\in S\,\bigl((q(s)=q(s^{\prime})\\[-2pt]&\quad \land q(t)=q(t^{\prime}))\rightarrow h(s,t)=h(s^{\prime},t^{\prime})\bigr)\end{aligned}}{%
    \rRE{44,32}}
  \proofstepwidestar[3,30]{\begin{aligned}[t]&t\in S\rightarrow \forall s^{\prime}\in S\,\forall t^{\prime}\in S\,\bigl((q(s)=q(s^{\prime})\\[-2pt]&\quad \land q(t)=q(t^{\prime}))\rightarrow h(s,t)=h(s^{\prime},t^{\prime})\bigr)\end{aligned}}{%
    \rUE{45}}
  \proofstepwidestar[3,30]{\begin{aligned}[t]&\forall s^{\prime}\in S\,\forall t^{\prime}\in S\,\bigl((q(s)=q(s^{\prime})\land q(t)=q(t^{\prime}))\\[-2pt]&\quad \rightarrow h(s,t)=h(s^{\prime},t^{\prime})\bigr)\end{aligned}}{%
    \rRE{46,33}}
  \proofstepwidestar[3,30]{\begin{aligned}[t]&s^{\prime}\in S\rightarrow \forall t^{\prime}\in S\,\bigl((q(s)=q(s^{\prime})\land q(t)=q(t^{\prime}))\\[-2pt]&\quad \rightarrow h(s,t)=h(s^{\prime},t^{\prime})\bigr)\end{aligned}}{%
    \rUE{47}}
  \proofstepwidestar[3,30,31]{\forall t^{\prime}\in S\,\bigl((q(s)=q(s^{\prime})\land q(t)=q(t^{\prime}))\rightarrow h(s,t)=h(s^{\prime},t^{\prime})\bigr)}{%
    \rRE{48,34}}
  \proofstepwidestar[3,30,31]{t^{\prime}\in S\rightarrow \bigl((q(s)=q(s^{\prime})\land q(t)=q(t^{\prime}))\rightarrow h(s,t)=h(s^{\prime},t^{\prime})\bigr)}{%
    \rUE{49}}
  \proofstepwidestar[3,30,31]{\begin{aligned}[t]&\bigl((q(s)=q(s^{\prime})\land q(t)=q(t^{\prime}))\\[-2pt]&\quad\rightarrow h(s,t)=h(s^{\prime},t^{\prime})\bigr)\end{aligned}}{%
    \rRE{50,35}}
  \proofstepwidestar[30,31]{q(s)=q(s^{\prime})\land q(t)=q(t^{\prime})}{%
    \rAIStar{38,41}}
  \proofstepwidestar[3,30,31]{h(s,t)=h(s^{\prime},t^{\prime})}{%
    \rRE{51,52}}
  \proofstepwidestar[3,30,31]{z=h(s^{\prime},t^{\prime})}{%
    \rIE{53,42}}
  \proofstepwidestar[3,30,31]{z=w}{%
    \FormulaRefAuto{a=b\dsep c=b\vdash a=c}[thm]{54,43}}
  \proofstepwidestar[3,27,30]{z=w}{%
    \rEEStar{29,31,55}}
  \proofstepwidestar[3,26,27]{z=w}{%
    \rEEStar{28,30,56}}
\closeproofpartwide
\Needspace{7\baselineskip}
\proofpartwide{Eindeutige Existenz}
\setcounter{proofstepnr}{57}
  \proofstepwidestar[1,2,3,4,5]{\begin{aligned}[t]&\exists!z\in B\,\exists s\in S\,\exists t\in S\,\bigl(u=q(s)\\[-2pt]&\quad\land v=q(t)\land z=h(s,t)\bigr)\end{aligned}}{%
    \UEI{25,26,27,57}}
\closeproofpartwide
\end{tabproofsplitwide}


\subsection{Funktionen aus eindeutigen Relationswerten}

Im folgenden Beweis steht \(D_R\) für den Mengenterm
\[
  \{p\in Q\times B\mid R(\pi_1(p),\pi_2(p))\}.
\]
Die Graphkonstruktion verwendet nur Aussonderung im kartesischen Produkt.


\FormulaThmDeltaK[Funktion aus einer Relation mit eindeutigem Wert]{%
  \begin{aligned}[t]&\forall u\in Q\,\exists!z\in B\,R(u,z)\vdash{}\\[-2pt]&\exists F\,\bigl(F\colon Q\to B\land\forall u\in Q\,\forall z\in B\,(R(u,z)\rightarrow F(u)=z)\bigr)\end{aligned}
}{UniqueRelationFunctionExists}{%
  \DeltaRow{Mengen}{Q\dsep B\dsep u\dsep z\dsep w}
  \DeltaRow{Prädikat}{R}
  \DeltaRow{Funktionen}{F}
}
\begin{tabproofsplitwide}
\proofpartwide{Der Relationgraph ist eine Funktion}
  \proofstepwidestar[1]{\forall u\in Q\,\exists!z\in B\,R(u,z)}{%
    \rA}
  \proofstepwidestar[]{D_R\subseteq Q\times B}{%
    \FormulaRefAuto{\{x\in A\mid P(x)\}\subseteq A}[thm]{}}
  \proofstepwidestar[3]{u\in Q}{%
    \rA}
  \proofstepwidestar[1]{u\in Q\rightarrow\exists!z\in B\,R(u,z)}{%
    \rUE{1}}
  \proofstepwidestar[1,3]{\exists!z\in B\,R(u,z)}{%
    \rRE{4,3}}
  \proofstepwidestar[1,3]{\exists z\in B\,R(u,z)}{%
    \FormulaRefAuto{\exists! x\,P(x)\vdash \exists x\,P(x)}[thm]{5}}
  \proofstepwidestar[7]{z\in B\land R(u,z)}{%
    \rA}
  \proofstepwidestar[7]{z\in B}{%
    \rAEa{7}}
  \proofstepwidestar[7]{R(u,z)}{%
    \rAEb{7}}
  \proofstepwidestar[]{(u,z)\in D_R\leftrightarrow u\in Q\land z\in B\land R(u,z)}{%
    \FormulaRefAuto{ProductComprehensionMembership}[thm]{}}
  \proofstepwidestar[3,7]{u\in Q\land z\in B\land R(u,z)}{%
    \rAIStar{3,8,9}}
  \proofstepwidestar[3,7]{(u,z)\in D_R}{%
    \rRE{\rLREb{10},11}}
  \proofstepwidestar[3,7]{z\in B\land(u,z)\in D_R}{%
    \rAIStar{8,12}}
  \proofstepwidestar[3,7]{\exists z\in B\,(u,z)\in D_R}{%
    \rEIStar{13}}
  \proofstepwidestar[1,3]{\exists z\in B\,(u,z)\in D_R}{%
    \rEEStar{6,7,14}}
  \proofstepwidestar[16]{z\in B\land(u,z)\in D_R}{%
    \rA}
  \proofstepwidestar[17]{w\in B\land(u,w)\in D_R}{%
    \rA}
  \proofstepwidestar[16]{(u,z)\in D_R}{%
    \rAEb{16}}
  \proofstepwidestar[17]{(u,w)\in D_R}{%
    \rAEb{17}}
  \proofstepwidestar[16]{u\in Q\land z\in B\land R(u,z)}{%
    \rRE{\rLREa{10},18}}
  \proofstepwidestar[]{(u,w)\in D_R\leftrightarrow u\in Q\land w\in B\land R(u,w)}{%
    \FormulaRefAuto{ProductComprehensionMembership}[thm]{}}
  \proofstepwidestar[17]{u\in Q\land w\in B\land R(u,w)}{%
    \rRE{\rLREa{21},19}}
  \proofstepwidestar[16]{z\in B}{%
    \rAEa{\rAEb{20}}}
  \proofstepwidestar[16]{R(u,z)}{%
    \rAEb{\rAEb{20}}}
  \proofstepwidestar[17]{w\in B}{%
    \rAEa{\rAEb{22}}}
  \proofstepwidestar[17]{R(u,w)}{%
    \rAEb{\rAEb{22}}}
  \proofstepwidestar[16]{z\in B\land R(u,z)}{%
    \rAIStar{23,24}}
  \proofstepwidestar[17]{w\in B\land R(u,w)}{%
    \rAIStar{25,26}}
  \proofstepwidestar[1,3,16,17]{z=w}{%
    \FormulaRefAuto{\exists! x\,P(x),\; P(a),\; P(b) \vdash a = b}[thm]{5,27,28}}
  \proofstepwidestar[1,3]{\exists!z\in B\,(u,z)\in D_R}{%
    \UEI{15,16,17,29}}
  \proofstepwidestar[1]{\forall u\in Q\,\exists!z\in B\,(u,z)\in D_R}{%
    \rUI{\rRI{3,30}}}
  \proofstepwidestar[1]{D_R\colon Q\to B}{%
    \FormulaRefAuto{UniqueValuedGraphFunction}[thm]{2,31}}
\closeproofpartwide
\proofpartwide{Auswertung des Relationgraphen}
\setcounter{proofstepnr}{32}
  \proofstepwidestar[33]{u\in Q}{%
    \rA}
  \proofstepwidestar[34]{z\in B}{%
    \rA}
  \proofstepwidestar[35]{R(u,z)}{%
    \rA}
  \proofstepwidestar[]{(u,z)\in D_R\leftrightarrow u\in Q\land z\in B\land R(u,z)}{%
    \FormulaRefAuto{ProductComprehensionMembership}[thm]{}}
  \proofstepwidestar[33,34,35]{u\in Q\land z\in B\land R(u,z)}{%
    \rAIStar{33,34,35}}
  \proofstepwidestar[33,34,35]{(u,z)\in D_R}{%
    \rRE{\rLREb{36},37}}
  \proofstepwidestar[1,33,34,35]{D_R(u)=z}{%
    \FormulaRefAuto{F\colon A\to B\dsep(x,y)\in F\vdash F(x)=y}[thm]{32,38}}
  \proofstepwidestar[1,33,34]{R(u,z)\rightarrow D_R(u)=z}{%
    \rRI{35,39}}
  \proofstepwidestar[1,33]{\forall z\in B\,(R(u,z)\rightarrow D_R(u)=z)}{%
    \rUI{\rRI{34,40}}}
  \proofstepwidestar[1]{\forall u\in Q\,\forall z\in B\,(R(u,z)\rightarrow D_R(u)=z)}{%
    \rUI{\rRI{33,41}}}
  \proofstepwidestar[1]{D_R\colon Q\to B\land \forall u\in Q\,\forall z\in B\,(R(u,z)\rightarrow D_R(u)=z)}{%
    \rAIStar{32,42}}
  \proofstepwidestar[1]{\exists F\,\bigl(F\colon Q\to B\land\forall u\in Q\,\forall z\in B\,(R(u,z)\rightarrow F(u)=z)\bigr)}{%
    \rEIStar{43}}
\closeproofpartwide
\end{tabproofsplitwide}


\Needspace{18\baselineskip}
\FormulaThmDeltaK[Faktorabbildung einer paarweise dargestellten Vorschrift]{%
  \begin{aligned}[t]&\forall a\forall b\,\bigl(P(a,b)\rightarrow\kappa(a,b)\in Q\bigr)\dsep{}\\[-2pt]&\forall u\in Q\,\exists a\exists b\,(P(a,b)\land u=\kappa(a,b))\dsep{}\\[-2pt]&\forall a\forall b\,\bigl(P(a,b)\rightarrow\tau(a,b)\in B\bigr)\dsep{}\\[-2pt]&\forall a\forall b\forall c\forall d\,\bigl((P(a,b)\land P(c,d)\land\kappa(a,b)=\kappa(c,d))\rightarrow\tau(a,b)=\tau(c,d)\bigr)\vdash{}\\[-2pt]&\exists!F\,\bigl(F\colon Q\to B\land \forall a\forall b\,\bigl(P(a,b)\rightarrow F(\kappa(a,b))=\tau(a,b)\bigr)\bigr)\end{aligned}
}{RepresentativePairFactorization}{%
  \DeltaRow{Mengen}{Q\dsep B\dsep u\dsep v\dsep z\dsep w}
  \DeltaRow{Prädikat}{P}
  \DeltaRow{Termschemata}{\kappa(a,b)\dsep\tau(a,b)}
  \DeltaRow{Funktionen}{F\dsep G\dsep H}
}
\begin{tabproofsplitwide}
\proofpartwide{Existenz der Faktorabbildung}
  \proofstepwidestar[1]{\forall a\forall b\,\bigl(P(a,b)\rightarrow\kappa(a,b)\in Q\bigr)}{%
    \rA}
  \proofstepwidestar[2]{\forall u\in Q\,\exists a\exists b\,(P(a,b)\land u=\kappa(a,b))}{%
    \rA}
  \proofstepwidestar[3]{\forall a\forall b\,\bigl(P(a,b)\rightarrow\tau(a,b)\in B\bigr)}{%
    \rA}
  \proofstepwidestar[4]{\forall a\forall b\forall c\forall d\,\bigl((P(a,b)\land P(c,d)\land\kappa(a,b)=\kappa(c,d))\rightarrow\tau(a,b)=\tau(c,d)\bigr)}{%
    \rA}
  \proofstepwidestar[5]{u\in Q}{%
    \rA}
  \proofstepwidestar[1,2,3,4,5]{\begin{aligned}[t]&\exists!z\in B\,\exists c\exists d\,(P(c,d)\\[-2pt]&\quad\land u=\kappa(c,d)\land z=\tau(c,d))\end{aligned}}{%
    \FormulaRefAuto{RepresentativePairValueUnique}[thm]{1,2,3,4,5}}
  \proofstepwidestar[1,2,3,4]{\begin{aligned}[t]&\forall u\in Q\,\exists!z\in B\,\exists c\exists d\,(P(c,d)\\[-2pt]&\quad\land u=\kappa(c,d)\land z=\tau(c,d))\end{aligned}}{%
    \rUI{\rRI{5,6}}}
  \proofstepwidestar[1,2,3,4]{\begin{aligned}[t]&\exists F\,(F\colon Q\to B\land \forall u\in Q\,\forall z\in B\,\bigl((\exists c\exists d\,(P(c,d)\\[-2pt]&\quad \land u=\kappa(c,d)\land z=\tau(c,d)))\rightarrow F(u)=z\bigr))\end{aligned}}{%
    \FormulaRefAuto{UniqueRelationFunctionExists}[thm]{7}}
  \proofstepwidestar[9]{\begin{aligned}[t]&F\colon Q\to B\land \forall u\in Q\,\forall z\in B\,\bigl((\exists c\exists d\,(P(c,d)\\[-2pt]&\quad\land u=\kappa(c,d)\land z=\tau(c,d)))\rightarrow F(u)=z\bigr)\end{aligned}}{%
    \rA}
  \proofstepwidestar[9]{F\colon Q\to B}{%
    \rAEa{9}}
  \proofstepwidestar[9]{\begin{aligned}[t]&\forall u\in Q\,\forall z\in B\,\bigl((\exists c\exists d\,(P(c,d)\\[-2pt]&\quad\land u=\kappa(c,d)\land z=\tau(c,d)))\\[-2pt]&\quad\rightarrow F(u)=z\bigr)\end{aligned}}{%
    \rAEb{9}}
  \proofstepwidestar[12]{P(a,b)}{%
    \rA}
  \proofstepwidestar[1]{P(a,b)\rightarrow\kappa(a,b)\in Q}{%
    \rUE{\rUE{1}}}
  \proofstepwidestar[1,12]{\kappa(a,b)\in Q}{%
    \rRE{13,12}}
  \proofstepwidestar[3]{P(a,b)\rightarrow\tau(a,b)\in B}{%
    \rUE{\rUE{3}}}
  \proofstepwidestar[3,12]{\tau(a,b)\in B}{%
    \rRE{15,12}}
  \proofstepwidestar[]{\kappa(a,b)=\kappa(a,b)}{%
    \rII}
  \proofstepwidestar[]{\tau(a,b)=\tau(a,b)}{%
    \rII}
  \proofstepwidestar[12]{P(a,b)\land\kappa(a,b)=\kappa(a,b)\land\tau(a,b)=\tau(a,b)}{%
    \rAIStar{12,17,18}}
  \proofstepwidestar[12]{\begin{aligned}[t]&\exists c\exists d\,(P(c,d)\land \kappa(a,b)=\kappa(c,d)\\[-2pt]&\quad\land \tau(a,b)=\tau(c,d))\end{aligned}}{%
    \rEIStar{19}}
  \proofstepwidestar[9]{\begin{aligned}[t]&\kappa(a,b)\in Q\rightarrow \forall z\in B\\[-2pt]&\quad\bigl((\exists c\exists d\,(P(c,d)\\[-2pt]&\quad\land \kappa(a,b)=\kappa(c,d)\land z=\tau(c,d)))\\[-2pt]&\quad\rightarrow F(\kappa(a,b))=z\bigr)\end{aligned}}{%
    \rUE{11}}
  \proofstepwidestar[1,9,12]{\begin{aligned}[t]&\forall z\in B\,\bigl((\exists c\exists d\,(P(c,d)\\[-2pt]&\quad\land \kappa(a,b)=\kappa(c,d)\land z=\tau(c,d)))\\[-2pt]&\quad\rightarrow F(\kappa(a,b))=z\bigr)\end{aligned}}{%
    \rRE{21,14}}
  \proofstepwidestar[1,9,12]{\begin{aligned}[t]&\tau(a,b)\in B\rightarrow((\exists c\exists d\,(P(c,d)\\[-2pt]&\quad\land \kappa(a,b)=\kappa(c,d)\land \tau(a,b)=\tau(c,d)))\\[-2pt]&\quad\rightarrow F(\kappa(a,b))=\tau(a,b))\end{aligned}}{%
    \rUE{22}}
  \proofstepwidestar[1,3,9,12]{\begin{aligned}[t]&(\exists c\exists d\,(P(c,d)\land \kappa(a,b)=\kappa(c,d)\\[-2pt]&\quad\land \tau(a,b)=\tau(c,d)))\\[-2pt]&\quad\rightarrow F(\kappa(a,b))=\tau(a,b)\end{aligned}}{%
    \rRE{23,16}}
  \proofstepwidestar[1,3,9,12]{F(\kappa(a,b))=\tau(a,b)}{%
    \rRE{24,20}}
  \proofstepwidestar[1,3,9]{\forall a\forall b\,\bigl(P(a,b)\rightarrow F(\kappa(a,b))=\tau(a,b)\bigr)}{%
    \rUI{\rUI{\rRI{12,25}}}}
  \proofstepwidestar[1,3,9]{F\colon Q\to B\land \forall a\forall b\,\bigl(P(a,b)\rightarrow F(\kappa(a,b))=\tau(a,b)\bigr)}{%
    \rAIStar{10,26}}
  \proofstepwidestar[1,3,9]{\exists F\,\bigl(F\colon Q\to B\land \forall a\forall b\,\bigl(P(a,b)\rightarrow F(\kappa(a,b))=\tau(a,b)\bigr)\bigr)}{%
    \rEIStar{27}}
  \proofstepwidestar[1,2,3,4]{\exists F\,\bigl(F\colon Q\to B\land \forall a\forall b\,\bigl(P(a,b)\rightarrow F(\kappa(a,b))=\tau(a,b)\bigr)\bigr)}{%
    \rEEStar{8,9,28}}
\closeproofpartwide
\proofpartwide{Eindeutigkeit der Faktorabbildung}
\setcounter{proofstepnr}{29}
  \proofstepwidestar[30]{G\colon Q\to B\land \forall a\forall b\,\bigl(P(a,b)\rightarrow G(\kappa(a,b))=\tau(a,b)\bigr)}{%
    \rA}
  \proofstepwidestar[31]{H\colon Q\to B\land \forall a\forall b\,\bigl(P(a,b)\rightarrow H(\kappa(a,b))=\tau(a,b)\bigr)}{%
    \rA}
  \proofstepwidestar[30]{G\colon Q\to B}{%
    \rAEa{30}}
  \proofstepwidestar[31]{H\colon Q\to B}{%
    \rAEa{31}}
  \proofstepwidestar[30]{\forall a\forall b\,\bigl(P(a,b)\rightarrow G(\kappa(a,b))=\tau(a,b)\bigr)}{%
    \rAEb{30}}
  \proofstepwidestar[31]{\forall a\forall b\,\bigl(P(a,b)\rightarrow H(\kappa(a,b))=\tau(a,b)\bigr)}{%
    \rAEb{31}}
  \proofstepwidestar[36]{u\in Q}{%
    \rA}
  \proofstepwidestar[2]{u\in Q\rightarrow\exists a\exists b\,(P(a,b)\land u=\kappa(a,b))}{%
    \rUE{2}}
  \proofstepwidestar[2,36]{\exists a\exists b\,(P(a,b)\land u=\kappa(a,b))}{%
    \rRE{37,36}}
  \proofstepwidestar[39]{P(a,b)\land u=\kappa(a,b)}{%
    \rA}
  \proofstepwidestar[39]{P(a,b)}{%
    \rAEa{39}}
  \proofstepwidestar[39]{u=\kappa(a,b)}{%
    \rAEb{39}}
  \proofstepwidestar[30]{P(a,b)\rightarrow G(\kappa(a,b))=\tau(a,b)}{%
    \rUE{\rUE{34}}}
  \proofstepwidestar[30,39]{G(\kappa(a,b))=\tau(a,b)}{%
    \rRE{42,40}}
  \proofstepwidestar[31]{P(a,b)\rightarrow H(\kappa(a,b))=\tau(a,b)}{%
    \rUE{\rUE{35}}}
  \proofstepwidestar[31,39]{H(\kappa(a,b))=\tau(a,b)}{%
    \rRE{44,40}}
  \proofstepwidestar[30,31,39]{G(\kappa(a,b))=H(\kappa(a,b))}{%
    \FormulaRefAuto{a=b\dsep c=b\vdash a=c}[thm]{43,45}}
  \proofstepwidestar[39]{\kappa(a,b)=u}{%
    \FormulaRefAuto{a=b\vdash b=a}[thm]{41}}
  \proofstepwidestar[30,31,39]{G(u)=H(u)}{%
    \rIE{47,46}}
  \proofstepwidestar[2,30,31,36]{G(u)=H(u)}{%
    \rEEStar{38,39,48}}
  \proofstepwidestar[2,30,31]{\forall u\in Q\,(G(u)=H(u))}{%
    \rUI{\rRI{36,49}}}
  \proofstepwidestar[2,30,31]{G=H}{%
    \FormulaRefAuto{FunctionExtensionality}[thm]{32,33,50}}
\closeproofpartwide
\proofpartwide{Eindeutige Existenz}
\setcounter{proofstepnr}{51}
  \proofstepwidestar[1,2,3,4]{\exists!F\,\bigl(F\colon Q\to B\land \forall a\forall b\,\bigl(P(a,b)\rightarrow F(\kappa(a,b))=\tau(a,b)\bigr)\bigr)}{%
    \UEI{29,30,31,51}}
\closeproofpartwide
\end{tabproofsplitwide}


\FormulaThmDeltaK[Eindeutige Faktorisierung durch eine Surjektion]{%
  \begin{aligned}[t]&q\colon S\sur Q\dsep{}\\[-2pt]&h\colon S\to B\dsep{}\\[-2pt]&\forall s\in S\,\forall t\in S\,(q(s)=q(t)\rightarrow h(s)=h(t))\vdash{}\\[-2pt]&\exists!F\,(F\colon Q\to B\land F\circ q=h)\end{aligned}
}{SurjectionFactorization}{%
  \DeltaRow{Mengen}{S\dsep Q\dsep B\dsep s\dsep t\dsep u\dsep z}
  \DeltaRow{Funktionen}{q\dsep h\dsep F\dsep G\dsep H}
}
\begin{tabproofsplitwide}
\proofpartwide{Existenz der Faktorabbildung}
  \proofstepwidestar[1]{q\colon S\sur Q}{%
    \rA}
  \proofstepwidestar[2]{h\colon S\to B}{%
    \rA}
  \proofstepwidestar[3]{\forall s\in S\,\forall t\in S\,(q(s)=q(t)\rightarrow h(s)=h(t))}{%
    \rA}
  \proofstepwidestar[4]{u\in Q}{%
    \rA}
  \proofstepwidestar[1,2,3,4]{\exists!z\in B\,\exists t\in S\,(u=q(t)\land z=h(t))}{%
    \FormulaRefAuto{SurjectionFiberValueUnique}[thm]{1,2,3,4}}
  \proofstepwidestar[1,2,3]{\forall u\in Q\,\exists!z\in B\,\exists t\in S\,(u=q(t)\land z=h(t))}{%
    \rUI{\rRI{4,5}}}
  \proofstepwidestar[1,2,3]{\exists F\,(F\colon Q\to B\land \forall u\in Q\,\forall z\in B\,\bigl((\exists t\in S\,(u=q(t)\land z=h(t)))\rightarrow F(u)=z\bigr))}{%
    \FormulaRefAuto{UniqueRelationFunctionExists}[thm]{6}}
  \proofstepwidestar[8]{F\colon Q\to B\land \forall u\in Q\,\forall z\in B\,\bigl((\exists t\in S\,(u=q(t)\land z=h(t)))\rightarrow F(u)=z\bigr)}{%
    \rA}
  \proofstepwidestar[8]{F\colon Q\to B}{%
    \rAEa{8}}
  \proofstepwidestar[8]{\begin{aligned}[t]&\forall u\in Q\,\forall z\in B\,\bigl((\exists t\in S\,(u=q(t)\\[-2pt]&\quad\land z=h(t)))\rightarrow F(u)=z\bigr)\end{aligned}}{%
    \rAEb{8}}
  \proofstepwidestar[1]{q\colon S\to Q}{%
    \FormulaRefAuto{Surjektive Funktion}[def]{1}}
  \proofstepwidestar[1,8]{F\circ q\colon S\to B}{%
    \FormulaRefAuto{G\circ F\colon A\to C}[thm]{11,9}}
  \proofstepwidestar[13]{s\in S}{%
    \rA}
  \proofstepwidestar[1,13]{q(s)\in Q}{%
    \FormulaRefAuto{F\colon A\to B\dsep x\in A\vdash F(x)\in B}[thm]{11,13}}
  \proofstepwidestar[2,13]{h(s)\in B}{%
    \FormulaRefAuto{F\colon A\to B\dsep x\in A\vdash F(x)\in B}[thm]{2,13}}
  \proofstepwidestar[]{q(s)=q(s)}{%
    \rII}
  \proofstepwidestar[]{h(s)=h(s)}{%
    \rII}
  \proofstepwidestar[]{q(s)=q(s)\land h(s)=h(s)}{%
    \rAIStar{16,17}}
  \proofstepwidestar[13]{\exists t\in S\,(q(s)=q(t)\land h(s)=h(t))}{%
    \rEIStar{13,18}}
  \proofstepwidestar[8]{q(s)\in Q\rightarrow \forall z\in B\,\bigl((\exists t\in S\,(q(s)=q(t)\land z=h(t)))\rightarrow F(q(s))=z\bigr)}{%
    \rUE{10}}
  \proofstepwidestar[1,8,13]{\forall z\in B\,\bigl((\exists t\in S\,(q(s)=q(t)\land z=h(t)))\rightarrow F(q(s))=z\bigr)}{%
    \rRE{20,14}}
  \proofstepwidestar[1,8,13]{h(s)\in B\rightarrow((\exists t\in S\,(q(s)=q(t)\land h(s)=h(t)))\rightarrow F(q(s))=h(s))}{%
    \rUE{21}}
  \proofstepwidestar[1,2,8,13]{\begin{aligned}[t]&(\exists t\in S\,(q(s)=q(t)\land h(s)=h(t)))\\[-2pt]&\quad\rightarrow F(q(s))=h(s)\end{aligned}}{%
    \rRE{22,15}}
  \proofstepwidestar[1,2,8,13]{F(q(s))=h(s)}{%
    \rRE{23,19}}
  \proofstepwidestar[1,8,13]{(F\circ q)(s)=F(q(s))}{%
    \FormulaRefAuto{x\in A\vdash(G\circ F)(x)=G(F(x))}[thm]{11,9,13}}
  \proofstepwidestar[1,2,8,13]{(F\circ q)(s)=h(s)}{%
    \rIE{24,25}}
  \proofstepwidestar[1,2,8]{\forall s\in S\,((F\circ q)(s)=h(s))}{%
    \rUI{\rRI{13,26}}}
  \proofstepwidestar[1,2,8]{F\circ q=h}{%
    \FormulaRefAuto{FunctionExtensionality}[thm]{12,2,27}}
  \proofstepwidestar[1,2,8]{F\colon Q\to B\land F\circ q=h}{%
    \rAIStar{9,28}}
  \proofstepwidestar[1,2,8]{\exists F\,(F\colon Q\to B\land F\circ q=h)}{%
    \rEIStar{29}}
  \proofstepwidestar[1,2,3]{\exists F\,(F\colon Q\to B\land F\circ q=h)}{%
    \rEEStar{7,8,30}}
\closeproofpartwide
\proofpartwide{Eindeutigkeit aus Surjektivität}
\setcounter{proofstepnr}{31}
  \proofstepwidestar[32]{G\colon Q\to B\land G\circ q=h}{%
    \rA}
  \proofstepwidestar[33]{H\colon Q\to B\land H\circ q=h}{%
    \rA}
  \proofstepwidestar[32]{G\colon Q\to B}{%
    \rAEa{32}}
  \proofstepwidestar[33]{H\colon Q\to B}{%
    \rAEa{33}}
  \proofstepwidestar[32]{G\circ q=h}{%
    \rAEb{32}}
  \proofstepwidestar[33]{H\circ q=h}{%
    \rAEb{33}}
  \proofstepwidestar[1,32]{G\circ q\colon S\to B}{%
    \FormulaRefAuto{G\circ F\colon A\to C}[thm]{11,34}}
  \proofstepwidestar[1,33]{H\circ q\colon S\to B}{%
    \FormulaRefAuto{G\circ F\colon A\to C}[thm]{11,35}}
  \proofstepwidestar[40]{u\in Q}{%
    \rA}
  \proofstepwidestar[1,40]{\exists s\in S\,q(s)=u}{%
    \FormulaRefAuto{Surjektivität}[ax]{1,40}}
  \proofstepwidestar[42]{s\in S\land q(s)=u}{%
    \rA}
  \proofstepwidestar[42]{s\in S}{%
    \rAEa{42}}
  \proofstepwidestar[42]{q(s)=u}{%
    \rAEb{42}}
  \proofstepwidestar[1,2,32,42]{(G\circ q)(s)=h(s)}{%
    \FormulaRefAuto{FunctionEqualityEvaluation}[thm]{38,2,43,36}}
  \proofstepwidestar[1,2,33,42]{(H\circ q)(s)=h(s)}{%
    \FormulaRefAuto{FunctionEqualityEvaluation}[thm]{39,2,43,37}}
  \proofstepwidestar[1,32,42]{(G\circ q)(s)=G(q(s))}{%
    \FormulaRefAuto{x\in A\vdash(G\circ F)(x)=G(F(x))}[thm]{11,34,43}}
  \proofstepwidestar[1,33,42]{(H\circ q)(s)=H(q(s))}{%
    \FormulaRefAuto{x\in A\vdash(G\circ F)(x)=G(F(x))}[thm]{11,35,43}}
  \proofstepwidestar[1,2,32,42]{G(q(s))=h(s)}{%
    \rIE{47,45}}
  \proofstepwidestar[1,2,33,42]{H(q(s))=h(s)}{%
    \rIE{48,46}}
  \proofstepwidestar[1,2,32,33,42]{G(q(s))=H(q(s))}{%
    \FormulaRefAuto{a=b\dsep c=b\vdash a=c}[thm]{49,50}}
  \proofstepwidestar[1,2,32,33,42]{G(u)=H(u)}{%
    \rIE{44,51}}
  \proofstepwidestar[1,2,32,33,40]{G(u)=H(u)}{%
    \rEEStar{41,42,52}}
  \proofstepwidestar[1,2,32,33]{\forall u\in Q\,(G(u)=H(u))}{%
    \rUI{\rRI{40,53}}}
  \proofstepwidestar[1,2,32,33]{G=H}{%
    \FormulaRefAuto{FunctionExtensionality}[thm]{34,35,54}}
\closeproofpartwide
\proofpartwide{Eindeutige Existenz}
\setcounter{proofstepnr}{55}
  \proofstepwidestar[1,2,3]{\exists!F\,(F\colon Q\to B\land F\circ q=h)}{%
    \UEI{31,32,33,55}}
\closeproofpartwide
\end{tabproofsplitwide}


\section{Leere Bereiche}

\FormulaThmDelta[Leere Relation ist surjektiv auf $\varnothing$]{%
  \varnothing\colon \varnothing\sur \varnothing%
}{%
  \DeltaDecl{Mengen}{y\dsep x}%
}
\begin{tabproofsplit}
\proofpart{Funktionstyp}
\proofstep{}{ \varnothing\colon \varnothing\to \varnothing }{%
    \FormulaRefAuto{\varnothing\colon \varnothing\to M}}


\closeproofpart

\proofpart{Surjektivität}
  \setcounter{proofstepnr}{1}% Fortlaufende Schritte dieses Beweises.
\proofstep{2}{y\in \varnothing}{\rA}
  \proofstep{}{y\notin \varnothing}{\FormulaRefAuto{x \not\in \varnothing}}
  \proofstep{2}{\bot}{\rBI{2,3}}
  \proofstep{2}{\exists x\in \varnothing\,(\varnothing(x)=y)}{%
    \FormulaRefAuto{P\dsep \neg P \vdash Q}{2,3}}
  \proofstep{}{y\in \varnothing\rightarrow \exists x\in \varnothing\,(\varnothing(x)=y)}{\rRI{2,5}}


\closeproofpart

\proofpart{Zusammenführung}
  \setcounter{proofstepnr}{6}% Fortlaufende Schritte dieses Beweises.
\proofstep{}{ \varnothing\colon \varnothing\sur \varnothing }{%
    \FormulaRefAuto{Surjektive Funktion}{1,6}}

\closeproofpart
\end{tabproofsplit}

\end{document}
